% !TEX program = lualatex
% ---------------------------------------------------------------------------------------------
%  Readback note 20: Theta functions of a CM elliptic curve: periods, tight frames, conjugates and divisibility of jets
%
%  Build:  python3 tools/rb.py note <slug>
%
%  One unit of a Lean artifact, rewritten as ordinary mathematics. The prose is hand-written. Lean names,
%  line ranges and counts come from data.tex; quoted Lean statements are copied from the artifact by the build.
%
%  Rules of the text: standard terminology; every symbol is explained where it first appears and listed again
%  in the notation at the end; no letter has two meanings; a proof is complete or says what it takes for granted;
%  Lean names appear only in the LEAN lines and in the concordance.
%
%  The signature rule (the build enforces it): a function, relation or predicate that the artifact defines is
%  written only through its notation macro, which the build makes from the spec (`notation`) with one slot for
%  each explicit argument of the Lean definition. No argument is fixed and dropped, also not "for brevity" and
%  not for one section. Every argument and every hypothesis of a stated declaration stands inside its
%  statement, or in a Setting that the statement cites by number; only implicit and instance arguments (the
%  ambient objects) may stand in the conventions.
% ---------------------------------------------------------------------------------------------
\documentclass[11pt,a4paper]{article}
\ifdefined\pdfvariable\pdfvariable suppressoptionalinfo \numexpr 1+2+4+8\relax\fi
\usepackage[a4paper,hmargin=35mm,top=22mm,bottom=27mm,footskip=12mm,nohead]{geometry}

\usepackage{fontspec}
\directlua{
  luaotfload.add_fallback("gfallback", {"STIXTwoMath-Regular.otf:mode=node;", "DejaVuSans.ttf:mode=node;"})
  luaotfload.add_fallback("gmonofallback", {"DejaVuSansMono.ttf:mode=node;", "STIXTwoMath-Regular.otf:mode=node;"})
}
\setmainfont{NewCM10}[Extension=.otf,UprightFont=*-Book,ItalicFont=*-BookItalic,
  BoldFont=*-Bold,BoldItalicFont=*-BoldItalic,SmallCapsFont=*-Book,SmallCapsFeatures={Letters=SmallCaps},
  RawFeature={fallback=gfallback}]
\setsansfont{NewCMSans10}[Extension=.otf,UprightFont=*-Book,ItalicFont=*-BookOblique,
  BoldFont=*-Bold,BoldItalicFont=*-BoldOblique,RawFeature={fallback=gfallback}]
\setmonofont{JetBrainsMono}[Extension=.otf,UprightFont=*-Regular,ItalicFont=*-Italic,
  BoldFont=*-SemiBold,BoldItalicFont=*-SemiBoldItalic,Scale=MatchLowercase,
  RawFeature={-calt;-liga;fallback=gmonofallback}]
\usepackage{amsmath}
\usepackage{amsthm}
\usepackage[math-style=TeX]{unicode-math}
\setmathfont{NewCMMath-Book.otf}
\emergencystretch=2.5em\relax

\usepackage{xcolor}
\usepackage{tikz}
\usetikzlibrary{calc,arrows.meta,positioning,backgrounds}
\usepackage[protrusion=true,expansion=true]{microtype}
\usepackage{enumitem}
\usepackage{needspace}
\usepackage{booktabs}
\usepackage{array}
\usepackage{longtable}
\usepackage{fancyhdr}
\usepackage{lastpage}
\usepackage{multicol}
\usepackage[hidelinks,pdfusetitle]{hyperref}

\definecolor{ink}{HTML}{0F172A}
\definecolor{inksoft}{HTML}{475569}
\definecolor{inkfaint}{HTML}{94A3B8}
\definecolor{rulegrey}{HTML}{CBD5E1}
\definecolor{arith}{HTML}{C2410C}
\definecolor{est}{HTML}{7C3AED}
\definecolor{ana}{HTML}{0891B2}
\definecolor{ok}{HTML}{047857}
\definecolor{leancol}{HTML}{3730A3}

\title{Theta functions of a CM elliptic curve: periods, tight frames, conjugates and divisibility of jets}
\author{grh-compact}
\date{}

% ------------------------------------------------------------------------------------- macros
\newcommand{\leanfont}{\ttfamily\fontsize{8.6}{10.5}\selectfont\color{leancol}}
\newcommand{\lean}[1]{\csname LN@#1\endcsname}
\newcommand{\fillin}[1]{{\color{arith}\itshape[#1]}}  % a place of the template that is not written yet; the build reports every one that is left
\newcommand{\leanquote}[1]{\csname LQ@#1\endcsname}
% generated by rb.py note; do not edit
\newcommand{\ConcDecls}{89}
\newcommand{\Pin}{72fead10}
\newcommand{\PinText}{that of commit {\leanfont 72fead10}}
\newcommand{\ArtifactName}{grh-compact}
\expandafter\newcommand\csname LN@bigOmegapos\endcsname{\href{https://github.com/lsernau/grh-lean/blob/72fead102f663d5c7a407964a13dc0cf9d977207/grh-compact/src/New/Curve/ConjugateGenerator.lean\#L15-L15}{{\leanfont Parameters.\allowbreak{}bigOmega\_pos}}}
\expandafter\newcommand\csname LN@jetScalenezero\endcsname{\href{https://github.com/lsernau/grh-lean/blob/72fead102f663d5c7a407964a13dc0cf9d977207/grh-compact/src/New/Curve/ConjugateGenerator.lean\#L18-L20}{{\leanfont Parameters.\allowbreak{}jetScale\_ne\_zero}}}
\expandafter\newcommand\csname LN@existsisJetConjugategenerator\endcsname{\href{https://github.com/lsernau/grh-lean/blob/72fead102f663d5c7a407964a13dc0cf9d977207/grh-compact/src/New/Curve/ConjugateGenerator.lean\#L26-L85}{{\leanfont Parameters.\allowbreak{}exists\_isJetConjugate\_generator}}}
\expandafter\newcommand\csname LN@ofintegral\endcsname{\href{https://github.com/lsernau/grh-lean/blob/72fead102f663d5c7a407964a13dc0cf9d977207/grh-compact/src/New/Curve/ConjugatesOfJets.lean\#L25-L26}{{\leanfont Parameters.\allowbreak{}integralAfter.\allowbreak{}of\_integral}}}
\expandafter\newcommand\csname LN@finHeisOpmemstandardData\endcsname{\href{https://github.com/lsernau/grh-lean/blob/72fead102f663d5c7a407964a13dc0cf9d977207/grh-compact/src/New/Curve/ConjugatesOfJets.lean\#L31-L36}{{\leanfont Parameters.\allowbreak{}finHeisOp\_mem\_standardData}}}
\expandafter\newcommand\csname LN@omegamemstandardData\endcsname{\href{https://github.com/lsernau/grh-lean/blob/72fead102f663d5c7a407964a13dc0cf9d977207/grh-compact/src/New/Curve/ConjugatesOfJets.lean\#L40-L50}{{\leanfont Parameters.\allowbreak{}omega\_mem\_standardData}}}
\expandafter\newcommand\csname LN@normjetSeriesle\endcsname{\href{https://github.com/lsernau/grh-lean/blob/72fead102f663d5c7a407964a13dc0cf9d977207/grh-compact/src/New/Curve/ConjugatesOfJets.lean\#L55-L84}{{\leanfont Parameters.\allowbreak{}norm\_jetSeries\_le}}}
\expandafter\newcommand\csname LN@normjetjetSeriesleone\endcsname{\href{https://github.com/lsernau/grh-lean/blob/72fead102f663d5c7a407964a13dc0cf9d977207/grh-compact/src/New/Curve/ConjugatesOfJets.lean\#L89-L96}{{\leanfont Parameters.\allowbreak{}norm\_jet\_jetSeries\_le\_one}}}
\expandafter\newcommand\csname LN@raises\endcsname{\href{https://github.com/lsernau/grh-lean/blob/72fead102f663d5c7a407964a13dc0cf9d977207/grh-compact/src/New/Curve/ConjugatesOfJets.lean\#L100-L102}{{\leanfont Kernel.\allowbreak{}raises}}}
\expandafter\newcommand\csname LN@smulmemraises\endcsname{\href{https://github.com/lsernau/grh-lean/blob/72fead102f663d5c7a407964a13dc0cf9d977207/grh-compact/src/New/Curve/ConjugatesOfJets.lean\#L106-L108}{{\leanfont Kernel.\allowbreak{}smul\_mem\_raises}}}
\expandafter\newcommand\csname LN@pairWeilmemstandardData\endcsname{\href{https://github.com/lsernau/grh-lean/blob/72fead102f663d5c7a407964a13dc0cf9d977207/grh-compact/src/New/Curve/ConjugatesOfJets.lean\#L112-L116}{{\leanfont Parameters.\allowbreak{}pairWeil\_mem\_standardData}}}
\expandafter\newcommand\csname LN@memraisesofjetbound\endcsname{\href{https://github.com/lsernau/grh-lean/blob/72fead102f663d5c7a407964a13dc0cf9d977207/grh-compact/src/New/Curve/ConjugatesOfJets.lean\#L121-L133}{{\leanfont Parameters.\allowbreak{}mem\_raises\_of\_jet\_bound}}}
\expandafter\newcommand\csname LN@memraisesofstandardData\endcsname{\href{https://github.com/lsernau/grh-lean/blob/72fead102f663d5c7a407964a13dc0cf9d977207/grh-compact/src/New/Curve/ConjugatesOfJets.lean\#L137-L143}{{\leanfont Parameters.\allowbreak{}mem\_raises\_of\_standardData}}}
\expandafter\newcommand\csname LN@memraisesofdvd\endcsname{\href{https://github.com/lsernau/grh-lean/blob/72fead102f663d5c7a407964a13dc0cf9d977207/grh-compact/src/New/Curve/ConjugatesOfJets.lean\#L148-L153}{{\leanfont Parameters.\allowbreak{}mem\_raises\_of\_dvd}}}
\expandafter\newcommand\csname LN@existssemilinearisJetConjugate\endcsname{\href{https://github.com/lsernau/grh-lean/blob/72fead102f663d5c7a407964a13dc0cf9d977207/grh-compact/src/New/Curve/ConjugatesOfJets.lean\#L164-L225}{{\leanfont Parameters.\allowbreak{}exists\_semilinear\_isJetConjugate}}}
\expandafter\newcommand\csname LN@IsFrobPacking\endcsname{\href{https://github.com/lsernau/grh-lean/blob/72fead102f663d5c7a407964a13dc0cf9d977207/grh-compact/src/New/Curve/DepthPairs.lean\#L14-L16}{{\leanfont IsFrobPacking}}}
\expandafter\newcommand\csname LN@nedofplaceClasseq\endcsname{\href{https://github.com/lsernau/grh-lean/blob/72fead102f663d5c7a407964a13dc0cf9d977207/grh-compact/src/New/Curve/DepthPairs.lean\#L18-L21}{{\leanfont Data.\allowbreak{}ne\_d\_of\_placeClass\_eq}}}
\expandafter\newcommand\csname LN@summablelocalAsk\endcsname{\href{https://github.com/lsernau/grh-lean/blob/72fead102f663d5c7a407964a13dc0cf9d977207/grh-compact/src/New/Curve/DepthPairs.lean\#L23-L26}{{\leanfont Data.\allowbreak{}summable\_localAsk}}}
\expandafter\newcommand\csname LN@depthLossInt\endcsname{\href{https://github.com/lsernau/grh-lean/blob/72fead102f663d5c7a407964a13dc0cf9d977207/grh-compact/src/New/Curve/DepthPairs.lean\#L29-L29}{{\leanfont Data.\allowbreak{}depthLossInt}}}
\expandafter\newcommand\csname LN@existsisFrobPacking\endcsname{\href{https://github.com/lsernau/grh-lean/blob/72fead102f663d5c7a407964a13dc0cf9d977207/grh-compact/src/New/Curve/DepthPairs.lean\#L34-L92}{{\leanfont Data.\allowbreak{}exists\_isFrobPacking}}}
\expandafter\newcommand\csname LN@existshasLargeSplitofpairs\endcsname{\href{https://github.com/lsernau/grh-lean/blob/72fead102f663d5c7a407964a13dc0cf9d977207/grh-compact/src/New/Curve/DepthPairs.lean\#L101-L146}{{\leanfont exists\_hasLargeSplit\_of\_pairs}}}
\expandafter\newcommand\csname LN@jetthetaprodintegral\endcsname{\href{https://github.com/lsernau/grh-lean/blob/72fead102f663d5c7a407964a13dc0cf9d977207/grh-compact/src/New/Curve/FormalParameter.lean\#L27-L55}{{\leanfont CMLattice.\allowbreak{}jet\_theta\_prod\_integral}}}
\expandafter\newcommand\csname LN@existsnormeqinvpowbounds\endcsname{\href{https://github.com/lsernau/grh-lean/blob/72fead102f663d5c7a407964a13dc0cf9d977207/grh-compact/src/New/Curve/FormalParameter.lean\#L59-L66}{{\leanfont exists\_norm\_eq\_inv\_pow\_bounds}}}
\expandafter\newcommand\csname LN@normcoeffleofderiv\endcsname{\href{https://github.com/lsernau/grh-lean/blob/72fead102f663d5c7a407964a13dc0cf9d977207/grh-compact/src/New/Curve/FormalParameter.lean\#L72-L95}{{\leanfont norm\_coeff\_le\_of\_deriv}}}
\expandafter\newcommand\csname LN@IsIntegralOnNullCone\endcsname{\href{https://github.com/lsernau/grh-lean/blob/72fead102f663d5c7a407964a13dc0cf9d977207/grh-compact/src/New/Curve/FormalParameter.lean\#L101-L105}{{\leanfont IsIntegralOnNullCone}}}
\expandafter\newcommand\csname LN@coeffsubsttangent\endcsname{\href{https://github.com/lsernau/grh-lean/blob/72fead102f663d5c7a407964a13dc0cf9d977207/grh-compact/src/New/Curve/FormalParameter.lean\#L119-L123}{{\leanfont coeff\_subst\_tangent}}}
\expandafter\newcommand\csname LN@existsformalparameter\endcsname{\href{https://github.com/lsernau/grh-lean/blob/72fead102f663d5c7a407964a13dc0cf9d977207/grh-compact/src/New/Curve/FormalParameter.lean\#L132-L248}{{\leanfont CMLattice.\allowbreak{}exists\_formal\_parameter}}}
\expandafter\newcommand\csname LN@taylorlineeqsubstjetSeries\endcsname{\href{https://github.com/lsernau/grh-lean/blob/72fead102f663d5c7a407964a13dc0cf9d977207/grh-compact/src/New/Curve/FrameSections.lean\#L22-L38}{{\leanfont Space.\allowbreak{}taylor\_line\_eq\_subst\_jetSeries}}}
\expandafter\newcommand\csname LN@eqzeroofjetSerieseqzero\endcsname{\href{https://github.com/lsernau/grh-lean/blob/72fead102f663d5c7a407964a13dc0cf9d977207/grh-compact/src/New/Curve/FrameSections.lean\#L42-L46}{{\leanfont Space.\allowbreak{}eq\_zero\_of\_jetSeries\_eq\_zero}}}
\expandafter\newcommand\csname LN@jetSerieseqprod\endcsname{\href{https://github.com/lsernau/grh-lean/blob/72fead102f663d5c7a407964a13dc0cf9d977207/grh-compact/src/New/Curve/FrameSections.lean\#L51-L67}{{\leanfont Space.\allowbreak{}jetSeries\_eq\_prod}}}
\expandafter\newcommand\csname LN@latticeModulecmCurve\endcsname{\href{https://github.com/lsernau/grh-lean/blob/72fead102f663d5c7a407964a13dc0cf9d977207/grh-compact/src/New/Curve/FrameSections.lean\#L72-L73}{{\leanfont Parameters.\allowbreak{}latticeModule\_cmCurve}}}
\expandafter\newcommand\csname LN@omegamemcmCurvelattice\endcsname{\href{https://github.com/lsernau/grh-lean/blob/72fead102f663d5c7a407964a13dc0cf9d977207/grh-compact/src/New/Curve/FrameSections.lean\#L76-L77}{{\leanfont Parameters.\allowbreak{}omega₀\_mem\_cmCurve\_lattice}}}
\expandafter\newcommand\csname LN@normmappsi\endcsname{\href{https://github.com/lsernau/grh-lean/blob/72fead102f663d5c7a407964a13dc0cf9d977207/grh-compact/src/New/Curve/FrameSections.lean\#L79-L81}{{\leanfont Parameters.\allowbreak{}norm\_map\_psi₀}}}
\expandafter\newcommand\csname LN@existsfinitePartofisTightFrame\endcsname{\href{https://github.com/lsernau/grh-lean/blob/72fead102f663d5c7a407964a13dc0cf9d977207/grh-compact/src/New/Curve/FrameSections.lean\#L89-L126}{{\leanfont Parameters.\allowbreak{}exists\_finitePart\_of\_isTightFrame}}}
\expandafter\newcommand\csname LN@isAdmissibleCombname\endcsname{\href{https://github.com/lsernau/grh-lean/blob/72fead102f663d5c7a407964a13dc0cf9d977207/grh-compact/src/New/Curve/IntegralVacuum.lean\#L19-L39}{{\leanfont isAdmissibleComb\_name}}}
\expandafter\newcommand\csname LN@vacuumScale\endcsname{\href{https://github.com/lsernau/grh-lean/blob/72fead102f663d5c7a407964a13dc0cf9d977207/grh-compact/src/New/Curve/IntegralVacuum.lean\#L43-L44}{{\leanfont Parameters.\allowbreak{}vacuumScale}}}
\expandafter\newcommand\csname LN@intVacuum\endcsname{\href{https://github.com/lsernau/grh-lean/blob/72fead102f663d5c7a407964a13dc0cf9d977207/grh-compact/src/New/Curve/IntegralVacuum.lean\#L49-L49}{{\leanfont Parameters.\allowbreak{}intVacuum}}}
\expandafter\newcommand\csname LN@integralAfter\endcsname{\href{https://github.com/lsernau/grh-lean/blob/72fead102f663d5c7a407964a13dc0cf9d977207/grh-compact/src/New/Curve/IntegralVacuum.lean\#L53-L55}{{\leanfont Parameters.\allowbreak{}integralAfter}}}
\expandafter\newcommand\csname LN@standardData\endcsname{\href{https://github.com/lsernau/grh-lean/blob/72fead102f663d5c7a407964a13dc0cf9d977207/grh-compact/src/New/Curve/IntegralVacuum.lean\#L59-L60}{{\leanfont Parameters.\allowbreak{}standardData}}}
\expandafter\newcommand\csname LN@kappanezero\endcsname{\href{https://github.com/lsernau/grh-lean/blob/72fead102f663d5c7a407964a13dc0cf9d977207/grh-compact/src/New/Curve/IntegralVacuum.lean\#L63-L63}{{\leanfont Parameters.\allowbreak{}kappa₀\_ne\_zero}}}
\expandafter\newcommand\csname LN@intVacuumfinVeceqprod\endcsname{\href{https://github.com/lsernau/grh-lean/blob/72fead102f663d5c7a407964a13dc0cf9d977207/grh-compact/src/New/Curve/IntegralVacuum.lean\#L69-L72}{{\leanfont Parameters.\allowbreak{}intVacuum\_finVec\_eq\_prod}}}
\expandafter\newcommand\csname LN@normintVacuumle\endcsname{\href{https://github.com/lsernau/grh-lean/blob/72fead102f663d5c7a407964a13dc0cf9d977207/grh-compact/src/New/Curve/IntegralVacuum.lean\#L77-L80}{{\leanfont Parameters.\allowbreak{}norm\_intVacuum\_le}}}
\expandafter\newcommand\csname LN@normintVacuumleone\endcsname{\href{https://github.com/lsernau/grh-lean/blob/72fead102f663d5c7a407964a13dc0cf9d977207/grh-compact/src/New/Curve/IntegralVacuum.lean\#L84-L86}{{\leanfont Parameters.\allowbreak{}norm\_intVacuum\_le\_one}}}
\expandafter\newcommand\csname LN@normdataInnerle\endcsname{\href{https://github.com/lsernau/grh-lean/blob/72fead102f663d5c7a407964a13dc0cf9d977207/grh-compact/src/New/Curve/IntegralVacuum.lean\#L90-L95}{{\leanfont Parameters.\allowbreak{}norm\_dataInner\_le}}}
\expandafter\newcommand\csname LN@existsintVacuumnezero\endcsname{\href{https://github.com/lsernau/grh-lean/blob/72fead102f663d5c7a407964a13dc0cf9d977207/grh-compact/src/New/Curve/IntegralVacuum.lean\#L97-L98}{{\leanfont Parameters.\allowbreak{}exists\_intVacuum\_ne\_zero}}}
\expandafter\newcommand\csname LN@existsnullcouplingframes\endcsname{\href{https://github.com/lsernau/grh-lean/blob/72fead102f663d5c7a407964a13dc0cf9d977207/grh-compact/src/New/Curve/JetDivisibility.lean\#L24-L53}{{\leanfont exists\_null\_coupling\_frames}}}
\expandafter\newcommand\csname LN@normschurJetle\endcsname{\href{https://github.com/lsernau/grh-lean/blob/72fead102f663d5c7a407964a13dc0cf9d977207/grh-compact/src/New/Curve/JetDivisibility.lean\#L65-L111}{{\leanfont norm\_schurJet\_le}}}
\expandafter\newcommand\csname LN@katzdivisibilitylimit\endcsname{\href{https://github.com/lsernau/grh-lean/blob/72fead102f663d5c7a407964a13dc0cf9d977207/grh-compact/src/New/Curve/KatzDivisibility.lean\#L15-L42}{{\leanfont CMLattice.\allowbreak{}katz\_divisibility\_limit}}}
\expandafter\newcommand\csname LN@localAsk\endcsname{\href{https://github.com/lsernau/grh-lean/blob/72fead102f663d5c7a407964a13dc0cf9d977207/grh-compact/src/New/Curve/LiveSum.lean\#L19-L20}{{\leanfont Data.\allowbreak{}localAsk}}}
\expandafter\newcommand\csname LN@forcedPart\endcsname{\href{https://github.com/lsernau/grh-lean/blob/72fead102f663d5c7a407964a13dc0cf9d977207/grh-compact/src/New/Curve/LiveSum.lean\#L24-L24}{{\leanfont Data.\allowbreak{}forcedPart}}}
\expandafter\newcommand\csname LN@gain\endcsname{\href{https://github.com/lsernau/grh-lean/blob/72fead102f663d5c7a407964a13dc0cf9d977207/grh-compact/src/New/Curve/LiveSum.lean\#L28-L28}{{\leanfont Data.\allowbreak{}gain}}}
\expandafter\newcommand\csname LQ@gain\endcsname{\leanquoteline{0}{def gain := D.m * (∑ τ ∈ D.tops, D.localAsk p D.a τ + D.k * D.forcedPart p)}}
\expandafter\newcommand\csname LN@IsSmallNonsplit\endcsname{\href{https://github.com/lsernau/grh-lean/blob/72fead102f663d5c7a407964a13dc0cf9d977207/grh-compact/src/New/Curve/LiveSum.lean\#L32-L32}{{\leanfont Data.\allowbreak{}IsSmallNonsplit}}}
\expandafter\newcommand\csname LN@IsShortPrime\endcsname{\href{https://github.com/lsernau/grh-lean/blob/72fead102f663d5c7a407964a13dc0cf9d977207/grh-compact/src/New/Curve/LiveSum.lean\#L35-L35}{{\leanfont Data.\allowbreak{}IsShortPrime}}}
\expandafter\newcommand\csname LN@loss\endcsname{\href{https://github.com/lsernau/grh-lean/blob/72fead102f663d5c7a407964a13dc0cf9d977207/grh-compact/src/New/Curve/LiveSum.lean\#L39-L39}{{\leanfont Data.\allowbreak{}loss}}}
\expandafter\newcommand\csname LN@ceilLogM\endcsname{\href{https://github.com/lsernau/grh-lean/blob/72fead102f663d5c7a407964a13dc0cf9d977207/grh-compact/src/New/Curve/LiveSum.lean\#L42-L42}{{\leanfont Data.\allowbreak{}ceilLogM}}}
\expandafter\newcommand\csname LN@deltaP\endcsname{\href{https://github.com/lsernau/grh-lean/blob/72fead102f663d5c7a407964a13dc0cf9d977207/grh-compact/src/New/Curve/LiveSum.lean\#L46-L46}{{\leanfont Data.\allowbreak{}deltaP}}}
\expandafter\newcommand\csname LN@jetCeil\endcsname{\href{https://github.com/lsernau/grh-lean/blob/72fead102f663d5c7a407964a13dc0cf9d977207/grh-compact/src/New/Curve/LiveSum.lean\#L49-L49}{{\leanfont Data.\allowbreak{}jetCeil}}}
\expandafter\newcommand\csname LN@ozero\endcsname{\href{https://github.com/lsernau/grh-lean/blob/72fead102f663d5c7a407964a13dc0cf9d977207/grh-compact/src/New/Curve/LiveSum.lean\#L53-L53}{{\leanfont Data.\allowbreak{}o₀}}}
\expandafter\newcommand\csname LN@IsServed\endcsname{\href{https://github.com/lsernau/grh-lean/blob/72fead102f663d5c7a407964a13dc0cf9d977207/grh-compact/src/New/Curve/LiveSum.lean\#L57-L57}{{\leanfont Data.\allowbreak{}IsServed}}}
\expandafter\newcommand\csname LN@IsCollision\endcsname{\href{https://github.com/lsernau/grh-lean/blob/72fead102f663d5c7a407964a13dc0cf9d977207/grh-compact/src/New/Curve/LiveSum.lean\#L61-L62}{{\leanfont Data.\allowbreak{}IsCollision}}}
\expandafter\newcommand\csname LN@collisionSum\endcsname{\href{https://github.com/lsernau/grh-lean/blob/72fead102f663d5c7a407964a13dc0cf9d977207/grh-compact/src/New/Curve/LiveSum.lean\#L65-L65}{{\leanfont Data.\allowbreak{}collisionSum}}}
\expandafter\newcommand\csname LN@couplingCount\endcsname{\href{https://github.com/lsernau/grh-lean/blob/72fead102f663d5c7a407964a13dc0cf9d977207/grh-compact/src/New/Curve/LiveSum.lean\#L69-L73}{{\leanfont couplingCount}}}
\expandafter\newcommand\csname LQ@couplingCount\endcsname{\leanquoteline{0}{def \_root\_.couplingCount \{Φ\} (p : ℕ) (κ : Φ → ℕ → ℕ) n' (ev : Finset Φ) σ σ' e :=}\leanquoteline{2}{let t σ := ∑ h ∈ .range n', ite (h = 1 ∧ σ ∉ ev) (κ σ h) 0}\leanquoteline{2}{let r σ := ∑ h ∈ .range n', ite (h = 1) 0 (κ σ h)}\leanquoteline{2}{let B σ := ∑ h ∈ .range n', h * ite (h = 1) 0 (κ σ h)}\leanquoteline{2}{(t σ / p \textasciicircum{} e : ℝ) ⊓ (r σ' + t σ') + r σ / p \textasciicircum{} e + (B σ / p \textasciicircum{} e : ℝ) ⊓ B σ'}}
\expandafter\newcommand\csname LN@kap\endcsname{\href{https://github.com/lsernau/grh-lean/blob/72fead102f663d5c7a407964a13dc0cf9d977207/grh-compact/src/New/Curve/LiveSum.lean\#L77-L78}{{\leanfont Data.\allowbreak{}κ}}}
\expandafter\newcommand\csname LN@kaple\endcsname{\href{https://github.com/lsernau/grh-lean/blob/72fead102f663d5c7a407964a13dc0cf9d977207/grh-compact/src/New/Curve/LiveSum.lean\#L80-L81}{{\leanfont Data.\allowbreak{}κ\_le}}}
\expandafter\newcommand\csname LN@lastEntry\endcsname{\href{https://github.com/lsernau/grh-lean/blob/72fead102f663d5c7a407964a13dc0cf9d977207/grh-compact/src/New/Curve/LiveSum.lean\#L85-L85}{{\leanfont Data.\allowbreak{}lastEntry}}}
\expandafter\newcommand\csname LN@keptColumns\endcsname{\href{https://github.com/lsernau/grh-lean/blob/72fead102f663d5c7a407964a13dc0cf9d977207/grh-compact/src/New/Curve/LiveSum.lean\#L89-L89}{{\leanfont Data.\allowbreak{}keptColumns}}}
\expandafter\newcommand\csname LN@columnsY\endcsname{\href{https://github.com/lsernau/grh-lean/blob/72fead102f663d5c7a407964a13dc0cf9d977207/grh-compact/src/New/Curve/LiveSum.lean\#L92-L92}{{\leanfont Data.\allowbreak{}columnsY}}}
\expandafter\newcommand\csname LN@hasSumliveinvpow\endcsname{\href{https://github.com/lsernau/grh-lean/blob/72fead102f663d5c7a407964a13dc0cf9d977207/grh-compact/src/New/Curve/LiveSum.lean\#L96-L102}{{\leanfont hasSum\_live\_inv\_pow}}}
\expandafter\newcommand\csname LN@vacuumRoot\endcsname{\href{https://github.com/lsernau/grh-lean/blob/72fead102f663d5c7a407964a13dc0cf9d977207/grh-compact/src/New/Curve/PeriodDefs.lean\#L13-L13}{{\leanfont Parameters.\allowbreak{}vacuumRoot}}}
\expandafter\newcommand\csname LN@normgalvacuumRoot\endcsname{\href{https://github.com/lsernau/grh-lean/blob/72fead102f663d5c7a407964a13dc0cf9d977207/grh-compact/src/New/Curve/PeriodDefs.lean\#L19-L30}{{\leanfont Parameters.\allowbreak{}norm\_gal\_vacuumRoot}}}
\expandafter\newcommand\csname LN@cmCurve\endcsname{\href{https://github.com/lsernau/grh-lean/blob/72fead102f663d5c7a407964a13dc0cf9d977207/grh-compact/src/New/Curve/PeriodDefs.lean\#L33-L34}{{\leanfont Parameters.\allowbreak{}cmCurve}}}
\expandafter\newcommand\csname LN@Omz\endcsname{\href{https://github.com/lsernau/grh-lean/blob/72fead102f663d5c7a407964a13dc0cf9d977207/grh-compact/src/New/Curve/PeriodDefs.lean\#L37-L38}{{\leanfont Parameters.\allowbreak{}Ω₀}}}
\expandafter\newcommand\csname LN@areaScale\endcsname{\href{https://github.com/lsernau/grh-lean/blob/72fead102f663d5c7a407964a13dc0cf9d977207/grh-compact/src/New/Curve/PeriodDefs.lean\#L41-L42}{{\leanfont Parameters.\allowbreak{}areaScale}}}
\expandafter\newcommand\csname LN@periodFactor\endcsname{\href{https://github.com/lsernau/grh-lean/blob/72fead102f663d5c7a407964a13dc0cf9d977207/grh-compact/src/New/Curve/PeriodDefs.lean\#L46-L47}{{\leanfont Parameters.\allowbreak{}periodFactor}}}
\expandafter\newcommand\csname LN@jetScale\endcsname{\href{https://github.com/lsernau/grh-lean/blob/72fead102f663d5c7a407964a13dc0cf9d977207/grh-compact/src/New/Curve/PeriodDefs.lean\#L51-L51}{{\leanfont Parameters.\allowbreak{}jetScale}}}
\expandafter\newcommand\csname LN@sumpoweqzeroofnull\endcsname{\href{https://github.com/lsernau/grh-lean/blob/72fead102f663d5c7a407964a13dc0cf9d977207/grh-compact/src/New/Curve/ThetaProducts.lean\#L19-L25}{{\leanfont sum\_pow\_eq\_zero\_of\_null}}}
\expandafter\newcommand\csname LN@coeffsumfrobnezeroiff\endcsname{\href{https://github.com/lsernau/grh-lean/blob/72fead102f663d5c7a407964a13dc0cf9d977207/grh-compact/src/New/Curve/ThetaProducts.lean\#L30-L39}{{\leanfont coeff\_sum\_frob\_ne\_zero\_iff}}}
\expandafter\newcommand\csname LN@thetaProducts\endcsname{\href{https://github.com/lsernau/grh-lean/blob/72fead102f663d5c7a407964a13dc0cf9d977207/grh-compact/src/New/Curve/ThetaProducts.lean\#L48-L51}{{\leanfont Parameters.\allowbreak{}thetaProducts}}}
\expandafter\newcommand\csname LQ@thetaProducts\endcsname{\leanquoteline{0}{def thetaProducts (δ : 𝓞 P.K) r c [Fintype c] := \{Θ | ∃ J, ∃ (\_ : Fintype J) (w : J → P.E s), ∃ M n}\leanquoteline{2}{b, ∃ y : J × c → ℂ, IsTightFrame w 1 M ∧ Θ = ∏ i, (taylorVia ιₚ (P.cmCurve.twistedTranslate 1 (y}\leanquoteline{2}{i) P.cmCurve.θ)).subst (linearSeries fun (τ, k) ↦ ite (k = i.2) (ιₚ (b * conj (Φ τ (w i.1)))) 0) ∧}\leanquoteline{2}{(p.Coprime M ∧ p.Coprime n) ∧ ‖ιₚ b‖ = r ∧ ∀ i, n * δ * y i ∈ P.cmCurve.Λ\}}}
\expandafter\newcommand\csname LN@existsframerowsfrob\endcsname{\href{https://github.com/lsernau/grh-lean/blob/72fead102f663d5c7a407964a13dc0cf9d977207/grh-compact/src/New/Curve/ThetaProducts.lean\#L57-L77}{{\leanfont Parameters.\allowbreak{}exists\_frame\_rows\_frob}}}
\expandafter\newcommand\csname LN@thetaProductsnormalform\endcsname{\href{https://github.com/lsernau/grh-lean/blob/72fead102f663d5c7a407964a13dc0cf9d977207/grh-compact/src/New/Curve/ThetaProducts.lean\#L86-L143}{{\leanfont Parameters.\allowbreak{}thetaProducts\_normal\_form}}}
\expandafter\newcommand\csname LQ@thetaProductsnormalform\endcsname{\leanquoteline{0}{lemma thetaProducts\_normal\_form (hp : chiK P.d p ≠ 1) : ∃ ϖ π, ‖ϖ‖ = p \textasciicircum{} (-nonsplitWeight P.d}\leanquoteline{2}{p) ∧ ‖π.1‖ = p \textasciicircum{} ((chiK P.d p - 1) / 2 : ℝ) ∧ ∀ \{I\} [Fintype I], ∀ r ≤ 1, ∀ Θ ∈ P.thetaProducts}\leanquoteline{2}{s 1 r I, (∀ B, IsIntegralOnNullCone (fun j v ↦ P.placeOfInt s (-p \textasciicircum{} j) * v) \{j ≤ p.log B | Odd j ∨}\leanquoteline{2}{p = P.d\} ϖ π B Θ) ∧ (chiK P.d p = -1 → r = 1 → ∃ (κ : ℂ\_[p]) (Ψ V : MvPowerSeries \_}\leanquoteline{2}{(valRingCp p)), ‖κ‖ = p \textasciicircum{} (1 - 1 / (p - 1) : ℝ) ∧ constantCoeff Ψ = 0 ∧ Θ.rescale (fun \_ ↦ ϖ) =}\leanquoteline{2}{expSubst (κ • seriesToCp Ψ) * seriesToCp V ∧ ∀ x d, (pderiv x (Ψ.map (residue \_))).coeff d ≠ 0 ↔}\leanquoteline{2}{∃ j y, Odd j ∧ (P.placeOfInt s (-p \textasciicircum{} j) * y.1, y.2) = x ∧ d = .single y (p \textasciicircum{} j))}}
\expandafter\newcommand\csname LN@coetraceeqsum\endcsname{\href{https://github.com/lsernau/grh-lean/blob/72fead102f663d5c7a407964a13dc0cf9d977207/grh-compact/src/New/Curve/TightFrameExists.lean\#L12-L13}{{\leanfont Parameters.\allowbreak{}coe\_trace\_eq\_sum}}}
\expandafter\newcommand\csname LN@hermmemquadInt\endcsname{\href{https://github.com/lsernau/grh-lean/blob/72fead102f663d5c7a407964a13dc0cf9d977207/grh-compact/src/New/Curve/TightFrameExists.lean\#L15-L23}{{\leanfont Parameters.\allowbreak{}herm\_mem\_quadInt}}}
\expandafter\newcommand\csname LN@adeleExtensionalpha\endcsname{\href{https://github.com/lsernau/grh-lean/blob/72fead102f663d5c7a407964a13dc0cf9d977207/grh-compact/src/New/Curve/TightFrameExists.lean\#L29-L31}{{\leanfont Parameters.\allowbreak{}adeleExtension\_alpha}}}
\expandafter\newcommand\csname LN@normiotaperiodFactor\endcsname{\href{https://github.com/lsernau/grh-lean/blob/72fead102f663d5c7a407964a13dc0cf9d977207/grh-compact/src/New/Curve/TightFrameExists.lean\#L36-L43}{{\leanfont Parameters.\allowbreak{}norm\_iota\_periodFactor}}}
\expandafter\newcommand\csname LN@normperiodFactormulconjleone\endcsname{\href{https://github.com/lsernau/grh-lean/blob/72fead102f663d5c7a407964a13dc0cf9d977207/grh-compact/src/New/Curve/TightFrameExists.lean\#L45-L48}{{\leanfont Parameters.\allowbreak{}norm\_periodFactor\_mul\_conj\_le\_one}}}
\expandafter\newcommand\csname LN@normiotaperiodFactoreqone\endcsname{\href{https://github.com/lsernau/grh-lean/blob/72fead102f663d5c7a407964a13dc0cf9d977207/grh-compact/src/New/Curve/TightFrameExists.lean\#L52-L55}{{\leanfont Parameters.\allowbreak{}norm\_iota\_periodFactor\_eq\_one}}}
\expandafter\newcommand\csname LN@isTightFramezetapow\endcsname{\href{https://github.com/lsernau/grh-lean/blob/72fead102f663d5c7a407964a13dc0cf9d977207/grh-compact/src/New/Curve/TightFrameExists.lean\#L59-L66}{{\leanfont Parameters.\allowbreak{}isTightFrame\_zeta\_pow}}}
\expandafter\newcommand\csname LN@existscyclotomicdifferentfactor\endcsname{\href{https://github.com/lsernau/grh-lean/blob/72fead102f663d5c7a407964a13dc0cf9d977207/grh-compact/src/New/Curve/TightFrameExists.lean\#L71-L85}{{\leanfont Parameters.\allowbreak{}exists\_cyclotomic\_different\_factor}}}
\expandafter\newcommand\csname LN@existsisTightFrame\endcsname{\href{https://github.com/lsernau/grh-lean/blob/72fead102f663d5c7a407964a13dc0cf9d977207/grh-compact/src/New/Curve/TightFrameExists.lean\#L90-L148}{{\leanfont Parameters.\allowbreak{}exists\_isTightFrame}}}
\newcommand{\ConcordanceRows}{\crow{bigOmegapos}{\href{https://github.com/lsernau/grh-lean/blob/72fead102f663d5c7a407964a13dc0cf9d977207/grh-compact/src/New/Curve/ConjugateGenerator.lean\#L15-L15}{{\leanfont Parameters.\allowbreak{}bigOmega\_pos}}}{ConjugateGenerator.lean}{15--15}
\crow{jetScalenezero}{\href{https://github.com/lsernau/grh-lean/blob/72fead102f663d5c7a407964a13dc0cf9d977207/grh-compact/src/New/Curve/ConjugateGenerator.lean\#L18-L20}{{\leanfont Parameters.\allowbreak{}jetScale\_ne\_zero}}}{ConjugateGenerator.lean}{18--20}
\crow{existsisJetConjugategenerator}{\href{https://github.com/lsernau/grh-lean/blob/72fead102f663d5c7a407964a13dc0cf9d977207/grh-compact/src/New/Curve/ConjugateGenerator.lean\#L26-L85}{{\leanfont Parameters.\allowbreak{}exists\_isJetConjugate\_generator}}}{ConjugateGenerator.lean}{26--85}
\crow{ofintegral}{\href{https://github.com/lsernau/grh-lean/blob/72fead102f663d5c7a407964a13dc0cf9d977207/grh-compact/src/New/Curve/ConjugatesOfJets.lean\#L25-L26}{{\leanfont Parameters.\allowbreak{}integralAfter.\allowbreak{}of\_integral}}}{ConjugatesOfJets.lean}{25--26}
\crow{finHeisOpmemstandardData}{\href{https://github.com/lsernau/grh-lean/blob/72fead102f663d5c7a407964a13dc0cf9d977207/grh-compact/src/New/Curve/ConjugatesOfJets.lean\#L31-L36}{{\leanfont Parameters.\allowbreak{}finHeisOp\_mem\_standardData}}}{ConjugatesOfJets.lean}{31--36}
\crow{omegamemstandardData}{\href{https://github.com/lsernau/grh-lean/blob/72fead102f663d5c7a407964a13dc0cf9d977207/grh-compact/src/New/Curve/ConjugatesOfJets.lean\#L40-L50}{{\leanfont Parameters.\allowbreak{}omega\_mem\_standardData}}}{ConjugatesOfJets.lean}{40--50}
\crow{normjetSeriesle}{\href{https://github.com/lsernau/grh-lean/blob/72fead102f663d5c7a407964a13dc0cf9d977207/grh-compact/src/New/Curve/ConjugatesOfJets.lean\#L55-L84}{{\leanfont Parameters.\allowbreak{}norm\_jetSeries\_le}}}{ConjugatesOfJets.lean}{55--84}
\crow{normjetjetSeriesleone}{\href{https://github.com/lsernau/grh-lean/blob/72fead102f663d5c7a407964a13dc0cf9d977207/grh-compact/src/New/Curve/ConjugatesOfJets.lean\#L89-L96}{{\leanfont Parameters.\allowbreak{}norm\_jet\_jetSeries\_le\_one}}}{ConjugatesOfJets.lean}{89--96}
\crow{raises}{\href{https://github.com/lsernau/grh-lean/blob/72fead102f663d5c7a407964a13dc0cf9d977207/grh-compact/src/New/Curve/ConjugatesOfJets.lean\#L100-L102}{{\leanfont Kernel.\allowbreak{}raises}}}{ConjugatesOfJets.lean}{100--102}
\crow{smulmemraises}{\href{https://github.com/lsernau/grh-lean/blob/72fead102f663d5c7a407964a13dc0cf9d977207/grh-compact/src/New/Curve/ConjugatesOfJets.lean\#L106-L108}{{\leanfont Kernel.\allowbreak{}smul\_mem\_raises}}}{ConjugatesOfJets.lean}{106--108}
\crow{pairWeilmemstandardData}{\href{https://github.com/lsernau/grh-lean/blob/72fead102f663d5c7a407964a13dc0cf9d977207/grh-compact/src/New/Curve/ConjugatesOfJets.lean\#L112-L116}{{\leanfont Parameters.\allowbreak{}pairWeil\_mem\_standardData}}}{ConjugatesOfJets.lean}{112--116}
\crow{memraisesofjetbound}{\href{https://github.com/lsernau/grh-lean/blob/72fead102f663d5c7a407964a13dc0cf9d977207/grh-compact/src/New/Curve/ConjugatesOfJets.lean\#L121-L133}{{\leanfont Parameters.\allowbreak{}mem\_raises\_of\_jet\_bound}}}{ConjugatesOfJets.lean}{121--133}
\crow{memraisesofstandardData}{\href{https://github.com/lsernau/grh-lean/blob/72fead102f663d5c7a407964a13dc0cf9d977207/grh-compact/src/New/Curve/ConjugatesOfJets.lean\#L137-L143}{{\leanfont Parameters.\allowbreak{}mem\_raises\_of\_standardData}}}{ConjugatesOfJets.lean}{137--143}
\crow{memraisesofdvd}{\href{https://github.com/lsernau/grh-lean/blob/72fead102f663d5c7a407964a13dc0cf9d977207/grh-compact/src/New/Curve/ConjugatesOfJets.lean\#L148-L153}{{\leanfont Parameters.\allowbreak{}mem\_raises\_of\_dvd}}}{ConjugatesOfJets.lean}{148--153}
\crow{existssemilinearisJetConjugate}{\href{https://github.com/lsernau/grh-lean/blob/72fead102f663d5c7a407964a13dc0cf9d977207/grh-compact/src/New/Curve/ConjugatesOfJets.lean\#L164-L225}{{\leanfont Parameters.\allowbreak{}exists\_semilinear\_isJetConjugate}}}{ConjugatesOfJets.lean}{164--225}
\crow{IsFrobPacking}{\href{https://github.com/lsernau/grh-lean/blob/72fead102f663d5c7a407964a13dc0cf9d977207/grh-compact/src/New/Curve/DepthPairs.lean\#L14-L16}{{\leanfont IsFrobPacking}}}{DepthPairs.lean}{14--16}
\crow{nedofplaceClasseq}{\href{https://github.com/lsernau/grh-lean/blob/72fead102f663d5c7a407964a13dc0cf9d977207/grh-compact/src/New/Curve/DepthPairs.lean\#L18-L21}{{\leanfont Data.\allowbreak{}ne\_d\_of\_placeClass\_eq}}}{DepthPairs.lean}{18--21}
\crow{summablelocalAsk}{\href{https://github.com/lsernau/grh-lean/blob/72fead102f663d5c7a407964a13dc0cf9d977207/grh-compact/src/New/Curve/DepthPairs.lean\#L23-L26}{{\leanfont Data.\allowbreak{}summable\_localAsk}}}{DepthPairs.lean}{23--26}
\crow{depthLossInt}{\href{https://github.com/lsernau/grh-lean/blob/72fead102f663d5c7a407964a13dc0cf9d977207/grh-compact/src/New/Curve/DepthPairs.lean\#L29-L29}{{\leanfont Data.\allowbreak{}depthLossInt}}}{DepthPairs.lean}{29--29}
\crow{existsisFrobPacking}{\href{https://github.com/lsernau/grh-lean/blob/72fead102f663d5c7a407964a13dc0cf9d977207/grh-compact/src/New/Curve/DepthPairs.lean\#L34-L92}{{\leanfont Data.\allowbreak{}exists\_isFrobPacking}}}{DepthPairs.lean}{34--92}
\crow{existshasLargeSplitofpairs}{\href{https://github.com/lsernau/grh-lean/blob/72fead102f663d5c7a407964a13dc0cf9d977207/grh-compact/src/New/Curve/DepthPairs.lean\#L101-L146}{{\leanfont exists\_hasLargeSplit\_of\_pairs}}}{DepthPairs.lean}{101--146}
\crow{jetthetaprodintegral}{\href{https://github.com/lsernau/grh-lean/blob/72fead102f663d5c7a407964a13dc0cf9d977207/grh-compact/src/New/Curve/FormalParameter.lean\#L27-L55}{{\leanfont CMLattice.\allowbreak{}jet\_theta\_prod\_integral}}}{FormalParameter.lean}{27--55}
\crow{existsnormeqinvpowbounds}{\href{https://github.com/lsernau/grh-lean/blob/72fead102f663d5c7a407964a13dc0cf9d977207/grh-compact/src/New/Curve/FormalParameter.lean\#L59-L66}{{\leanfont exists\_norm\_eq\_inv\_pow\_bounds}}}{FormalParameter.lean}{59--66}
\crow{normcoeffleofderiv}{\href{https://github.com/lsernau/grh-lean/blob/72fead102f663d5c7a407964a13dc0cf9d977207/grh-compact/src/New/Curve/FormalParameter.lean\#L72-L95}{{\leanfont norm\_coeff\_le\_of\_deriv}}}{FormalParameter.lean}{72--95}
\crow{IsIntegralOnNullCone}{\href{https://github.com/lsernau/grh-lean/blob/72fead102f663d5c7a407964a13dc0cf9d977207/grh-compact/src/New/Curve/FormalParameter.lean\#L101-L105}{{\leanfont IsIntegralOnNullCone}}}{FormalParameter.lean}{101--105}
\crow{coeffsubsttangent}{\href{https://github.com/lsernau/grh-lean/blob/72fead102f663d5c7a407964a13dc0cf9d977207/grh-compact/src/New/Curve/FormalParameter.lean\#L119-L123}{{\leanfont coeff\_subst\_tangent}}}{FormalParameter.lean}{119--123}
\crow{existsformalparameter}{\href{https://github.com/lsernau/grh-lean/blob/72fead102f663d5c7a407964a13dc0cf9d977207/grh-compact/src/New/Curve/FormalParameter.lean\#L132-L248}{{\leanfont CMLattice.\allowbreak{}exists\_formal\_parameter}}}{FormalParameter.lean}{132--248}
\crow{taylorlineeqsubstjetSeries}{\href{https://github.com/lsernau/grh-lean/blob/72fead102f663d5c7a407964a13dc0cf9d977207/grh-compact/src/New/Curve/FrameSections.lean\#L22-L38}{{\leanfont Space.\allowbreak{}taylor\_line\_eq\_subst\_jetSeries}}}{FrameSections.lean}{22--38}
\crow{eqzeroofjetSerieseqzero}{\href{https://github.com/lsernau/grh-lean/blob/72fead102f663d5c7a407964a13dc0cf9d977207/grh-compact/src/New/Curve/FrameSections.lean\#L42-L46}{{\leanfont Space.\allowbreak{}eq\_zero\_of\_jetSeries\_eq\_zero}}}{FrameSections.lean}{42--46}
\crow{jetSerieseqprod}{\href{https://github.com/lsernau/grh-lean/blob/72fead102f663d5c7a407964a13dc0cf9d977207/grh-compact/src/New/Curve/FrameSections.lean\#L51-L67}{{\leanfont Space.\allowbreak{}jetSeries\_eq\_prod}}}{FrameSections.lean}{51--67}
\crow{latticeModulecmCurve}{\href{https://github.com/lsernau/grh-lean/blob/72fead102f663d5c7a407964a13dc0cf9d977207/grh-compact/src/New/Curve/FrameSections.lean\#L72-L73}{{\leanfont Parameters.\allowbreak{}latticeModule\_cmCurve}}}{FrameSections.lean}{72--73}
\crow{omegamemcmCurvelattice}{\href{https://github.com/lsernau/grh-lean/blob/72fead102f663d5c7a407964a13dc0cf9d977207/grh-compact/src/New/Curve/FrameSections.lean\#L76-L77}{{\leanfont Parameters.\allowbreak{}omega₀\_mem\_cmCurve\_lattice}}}{FrameSections.lean}{76--77}
\crow{normmappsi}{\href{https://github.com/lsernau/grh-lean/blob/72fead102f663d5c7a407964a13dc0cf9d977207/grh-compact/src/New/Curve/FrameSections.lean\#L79-L81}{{\leanfont Parameters.\allowbreak{}norm\_map\_psi₀}}}{FrameSections.lean}{79--81}
\crow{existsfinitePartofisTightFrame}{\href{https://github.com/lsernau/grh-lean/blob/72fead102f663d5c7a407964a13dc0cf9d977207/grh-compact/src/New/Curve/FrameSections.lean\#L89-L126}{{\leanfont Parameters.\allowbreak{}exists\_finitePart\_of\_isTightFrame}}}{FrameSections.lean}{89--126}
\crow{isAdmissibleCombname}{\href{https://github.com/lsernau/grh-lean/blob/72fead102f663d5c7a407964a13dc0cf9d977207/grh-compact/src/New/Curve/IntegralVacuum.lean\#L19-L39}{{\leanfont isAdmissibleComb\_name}}}{IntegralVacuum.lean}{19--39}
\crow{vacuumScale}{\href{https://github.com/lsernau/grh-lean/blob/72fead102f663d5c7a407964a13dc0cf9d977207/grh-compact/src/New/Curve/IntegralVacuum.lean\#L43-L44}{{\leanfont Parameters.\allowbreak{}vacuumScale}}}{IntegralVacuum.lean}{43--44}
\crow{intVacuum}{\href{https://github.com/lsernau/grh-lean/blob/72fead102f663d5c7a407964a13dc0cf9d977207/grh-compact/src/New/Curve/IntegralVacuum.lean\#L49-L49}{{\leanfont Parameters.\allowbreak{}intVacuum}}}{IntegralVacuum.lean}{49--49}
\crow{integralAfter}{\href{https://github.com/lsernau/grh-lean/blob/72fead102f663d5c7a407964a13dc0cf9d977207/grh-compact/src/New/Curve/IntegralVacuum.lean\#L53-L55}{{\leanfont Parameters.\allowbreak{}integralAfter}}}{IntegralVacuum.lean}{53--55}
\crow{standardData}{\href{https://github.com/lsernau/grh-lean/blob/72fead102f663d5c7a407964a13dc0cf9d977207/grh-compact/src/New/Curve/IntegralVacuum.lean\#L59-L60}{{\leanfont Parameters.\allowbreak{}standardData}}}{IntegralVacuum.lean}{59--60}
\crow{kappanezero}{\href{https://github.com/lsernau/grh-lean/blob/72fead102f663d5c7a407964a13dc0cf9d977207/grh-compact/src/New/Curve/IntegralVacuum.lean\#L63-L63}{{\leanfont Parameters.\allowbreak{}kappa₀\_ne\_zero}}}{IntegralVacuum.lean}{63--63}
\crow{intVacuumfinVeceqprod}{\href{https://github.com/lsernau/grh-lean/blob/72fead102f663d5c7a407964a13dc0cf9d977207/grh-compact/src/New/Curve/IntegralVacuum.lean\#L69-L72}{{\leanfont Parameters.\allowbreak{}intVacuum\_finVec\_eq\_prod}}}{IntegralVacuum.lean}{69--72}
\crow{normintVacuumle}{\href{https://github.com/lsernau/grh-lean/blob/72fead102f663d5c7a407964a13dc0cf9d977207/grh-compact/src/New/Curve/IntegralVacuum.lean\#L77-L80}{{\leanfont Parameters.\allowbreak{}norm\_intVacuum\_le}}}{IntegralVacuum.lean}{77--80}
\crow{normintVacuumleone}{\href{https://github.com/lsernau/grh-lean/blob/72fead102f663d5c7a407964a13dc0cf9d977207/grh-compact/src/New/Curve/IntegralVacuum.lean\#L84-L86}{{\leanfont Parameters.\allowbreak{}norm\_intVacuum\_le\_one}}}{IntegralVacuum.lean}{84--86}
\crow{normdataInnerle}{\href{https://github.com/lsernau/grh-lean/blob/72fead102f663d5c7a407964a13dc0cf9d977207/grh-compact/src/New/Curve/IntegralVacuum.lean\#L90-L95}{{\leanfont Parameters.\allowbreak{}norm\_dataInner\_le}}}{IntegralVacuum.lean}{90--95}
\crow{existsintVacuumnezero}{\href{https://github.com/lsernau/grh-lean/blob/72fead102f663d5c7a407964a13dc0cf9d977207/grh-compact/src/New/Curve/IntegralVacuum.lean\#L97-L98}{{\leanfont Parameters.\allowbreak{}exists\_intVacuum\_ne\_zero}}}{IntegralVacuum.lean}{97--98}
\crow{existsnullcouplingframes}{\href{https://github.com/lsernau/grh-lean/blob/72fead102f663d5c7a407964a13dc0cf9d977207/grh-compact/src/New/Curve/JetDivisibility.lean\#L24-L53}{{\leanfont exists\_null\_coupling\_frames}}}{JetDivisibility.lean}{24--53}
\crow{normschurJetle}{\href{https://github.com/lsernau/grh-lean/blob/72fead102f663d5c7a407964a13dc0cf9d977207/grh-compact/src/New/Curve/JetDivisibility.lean\#L65-L111}{{\leanfont norm\_schurJet\_le}}}{JetDivisibility.lean}{65--111}
\crow{katzdivisibilitylimit}{\href{https://github.com/lsernau/grh-lean/blob/72fead102f663d5c7a407964a13dc0cf9d977207/grh-compact/src/New/Curve/KatzDivisibility.lean\#L15-L42}{{\leanfont CMLattice.\allowbreak{}katz\_divisibility\_limit}}}{KatzDivisibility.lean}{15--42}
\crow{localAsk}{\href{https://github.com/lsernau/grh-lean/blob/72fead102f663d5c7a407964a13dc0cf9d977207/grh-compact/src/New/Curve/LiveSum.lean\#L19-L20}{{\leanfont Data.\allowbreak{}localAsk}}}{LiveSum.lean}{19--20}
\crow{forcedPart}{\href{https://github.com/lsernau/grh-lean/blob/72fead102f663d5c7a407964a13dc0cf9d977207/grh-compact/src/New/Curve/LiveSum.lean\#L24-L24}{{\leanfont Data.\allowbreak{}forcedPart}}}{LiveSum.lean}{24--24}
\crow{gain}{\href{https://github.com/lsernau/grh-lean/blob/72fead102f663d5c7a407964a13dc0cf9d977207/grh-compact/src/New/Curve/LiveSum.lean\#L28-L28}{{\leanfont Data.\allowbreak{}gain}}}{LiveSum.lean}{28--28}
\crow{IsSmallNonsplit}{\href{https://github.com/lsernau/grh-lean/blob/72fead102f663d5c7a407964a13dc0cf9d977207/grh-compact/src/New/Curve/LiveSum.lean\#L32-L32}{{\leanfont Data.\allowbreak{}IsSmallNonsplit}}}{LiveSum.lean}{32--32}
\crow{IsShortPrime}{\href{https://github.com/lsernau/grh-lean/blob/72fead102f663d5c7a407964a13dc0cf9d977207/grh-compact/src/New/Curve/LiveSum.lean\#L35-L35}{{\leanfont Data.\allowbreak{}IsShortPrime}}}{LiveSum.lean}{35--35}
\crow{loss}{\href{https://github.com/lsernau/grh-lean/blob/72fead102f663d5c7a407964a13dc0cf9d977207/grh-compact/src/New/Curve/LiveSum.lean\#L39-L39}{{\leanfont Data.\allowbreak{}loss}}}{LiveSum.lean}{39--39}
\crow{ceilLogM}{\href{https://github.com/lsernau/grh-lean/blob/72fead102f663d5c7a407964a13dc0cf9d977207/grh-compact/src/New/Curve/LiveSum.lean\#L42-L42}{{\leanfont Data.\allowbreak{}ceilLogM}}}{LiveSum.lean}{42--42}
\crow{deltaP}{\href{https://github.com/lsernau/grh-lean/blob/72fead102f663d5c7a407964a13dc0cf9d977207/grh-compact/src/New/Curve/LiveSum.lean\#L46-L46}{{\leanfont Data.\allowbreak{}deltaP}}}{LiveSum.lean}{46--46}
\crow{jetCeil}{\href{https://github.com/lsernau/grh-lean/blob/72fead102f663d5c7a407964a13dc0cf9d977207/grh-compact/src/New/Curve/LiveSum.lean\#L49-L49}{{\leanfont Data.\allowbreak{}jetCeil}}}{LiveSum.lean}{49--49}
\crow{ozero}{\href{https://github.com/lsernau/grh-lean/blob/72fead102f663d5c7a407964a13dc0cf9d977207/grh-compact/src/New/Curve/LiveSum.lean\#L53-L53}{{\leanfont Data.\allowbreak{}o₀}}}{LiveSum.lean}{53--53}
\crow{IsServed}{\href{https://github.com/lsernau/grh-lean/blob/72fead102f663d5c7a407964a13dc0cf9d977207/grh-compact/src/New/Curve/LiveSum.lean\#L57-L57}{{\leanfont Data.\allowbreak{}IsServed}}}{LiveSum.lean}{57--57}
\crow{IsCollision}{\href{https://github.com/lsernau/grh-lean/blob/72fead102f663d5c7a407964a13dc0cf9d977207/grh-compact/src/New/Curve/LiveSum.lean\#L61-L62}{{\leanfont Data.\allowbreak{}IsCollision}}}{LiveSum.lean}{61--62}
\crow{collisionSum}{\href{https://github.com/lsernau/grh-lean/blob/72fead102f663d5c7a407964a13dc0cf9d977207/grh-compact/src/New/Curve/LiveSum.lean\#L65-L65}{{\leanfont Data.\allowbreak{}collisionSum}}}{LiveSum.lean}{65--65}
\crow{couplingCount}{\href{https://github.com/lsernau/grh-lean/blob/72fead102f663d5c7a407964a13dc0cf9d977207/grh-compact/src/New/Curve/LiveSum.lean\#L69-L73}{{\leanfont couplingCount}}}{LiveSum.lean}{69--73}
\crow{kap}{\href{https://github.com/lsernau/grh-lean/blob/72fead102f663d5c7a407964a13dc0cf9d977207/grh-compact/src/New/Curve/LiveSum.lean\#L77-L78}{{\leanfont Data.\allowbreak{}κ}}}{LiveSum.lean}{77--78}
\crow{kaple}{\href{https://github.com/lsernau/grh-lean/blob/72fead102f663d5c7a407964a13dc0cf9d977207/grh-compact/src/New/Curve/LiveSum.lean\#L80-L81}{{\leanfont Data.\allowbreak{}κ\_le}}}{LiveSum.lean}{80--81}
\crow{lastEntry}{\href{https://github.com/lsernau/grh-lean/blob/72fead102f663d5c7a407964a13dc0cf9d977207/grh-compact/src/New/Curve/LiveSum.lean\#L85-L85}{{\leanfont Data.\allowbreak{}lastEntry}}}{LiveSum.lean}{85--85}
\crow{keptColumns}{\href{https://github.com/lsernau/grh-lean/blob/72fead102f663d5c7a407964a13dc0cf9d977207/grh-compact/src/New/Curve/LiveSum.lean\#L89-L89}{{\leanfont Data.\allowbreak{}keptColumns}}}{LiveSum.lean}{89--89}
\crow{columnsY}{\href{https://github.com/lsernau/grh-lean/blob/72fead102f663d5c7a407964a13dc0cf9d977207/grh-compact/src/New/Curve/LiveSum.lean\#L92-L92}{{\leanfont Data.\allowbreak{}columnsY}}}{LiveSum.lean}{92--92}
\crow{hasSumliveinvpow}{\href{https://github.com/lsernau/grh-lean/blob/72fead102f663d5c7a407964a13dc0cf9d977207/grh-compact/src/New/Curve/LiveSum.lean\#L96-L102}{{\leanfont hasSum\_live\_inv\_pow}}}{LiveSum.lean}{96--102}
\crow{vacuumRoot}{\href{https://github.com/lsernau/grh-lean/blob/72fead102f663d5c7a407964a13dc0cf9d977207/grh-compact/src/New/Curve/PeriodDefs.lean\#L13-L13}{{\leanfont Parameters.\allowbreak{}vacuumRoot}}}{PeriodDefs.lean}{13--13}
\crow{normgalvacuumRoot}{\href{https://github.com/lsernau/grh-lean/blob/72fead102f663d5c7a407964a13dc0cf9d977207/grh-compact/src/New/Curve/PeriodDefs.lean\#L19-L30}{{\leanfont Parameters.\allowbreak{}norm\_gal\_vacuumRoot}}}{PeriodDefs.lean}{19--30}
\crow{cmCurve}{\href{https://github.com/lsernau/grh-lean/blob/72fead102f663d5c7a407964a13dc0cf9d977207/grh-compact/src/New/Curve/PeriodDefs.lean\#L33-L34}{{\leanfont Parameters.\allowbreak{}cmCurve}}}{PeriodDefs.lean}{33--34}
\crow{Omz}{\href{https://github.com/lsernau/grh-lean/blob/72fead102f663d5c7a407964a13dc0cf9d977207/grh-compact/src/New/Curve/PeriodDefs.lean\#L37-L38}{{\leanfont Parameters.\allowbreak{}Ω₀}}}{PeriodDefs.lean}{37--38}
\crow{areaScale}{\href{https://github.com/lsernau/grh-lean/blob/72fead102f663d5c7a407964a13dc0cf9d977207/grh-compact/src/New/Curve/PeriodDefs.lean\#L41-L42}{{\leanfont Parameters.\allowbreak{}areaScale}}}{PeriodDefs.lean}{41--42}
\crow{periodFactor}{\href{https://github.com/lsernau/grh-lean/blob/72fead102f663d5c7a407964a13dc0cf9d977207/grh-compact/src/New/Curve/PeriodDefs.lean\#L46-L47}{{\leanfont Parameters.\allowbreak{}periodFactor}}}{PeriodDefs.lean}{46--47}
\crow{jetScale}{\href{https://github.com/lsernau/grh-lean/blob/72fead102f663d5c7a407964a13dc0cf9d977207/grh-compact/src/New/Curve/PeriodDefs.lean\#L51-L51}{{\leanfont Parameters.\allowbreak{}jetScale}}}{PeriodDefs.lean}{51--51}
\crow{sumpoweqzeroofnull}{\href{https://github.com/lsernau/grh-lean/blob/72fead102f663d5c7a407964a13dc0cf9d977207/grh-compact/src/New/Curve/ThetaProducts.lean\#L19-L25}{{\leanfont sum\_pow\_eq\_zero\_of\_null}}}{ThetaProducts.lean}{19--25}
\crow{coeffsumfrobnezeroiff}{\href{https://github.com/lsernau/grh-lean/blob/72fead102f663d5c7a407964a13dc0cf9d977207/grh-compact/src/New/Curve/ThetaProducts.lean\#L30-L39}{{\leanfont coeff\_sum\_frob\_ne\_zero\_iff}}}{ThetaProducts.lean}{30--39}
\crow{thetaProducts}{\href{https://github.com/lsernau/grh-lean/blob/72fead102f663d5c7a407964a13dc0cf9d977207/grh-compact/src/New/Curve/ThetaProducts.lean\#L48-L51}{{\leanfont Parameters.\allowbreak{}thetaProducts}}}{ThetaProducts.lean}{48--51}
\crow{existsframerowsfrob}{\href{https://github.com/lsernau/grh-lean/blob/72fead102f663d5c7a407964a13dc0cf9d977207/grh-compact/src/New/Curve/ThetaProducts.lean\#L57-L77}{{\leanfont Parameters.\allowbreak{}exists\_frame\_rows\_frob}}}{ThetaProducts.lean}{57--77}
\crow{thetaProductsnormalform}{\href{https://github.com/lsernau/grh-lean/blob/72fead102f663d5c7a407964a13dc0cf9d977207/grh-compact/src/New/Curve/ThetaProducts.lean\#L86-L143}{{\leanfont Parameters.\allowbreak{}thetaProducts\_normal\_form}}}{ThetaProducts.lean}{86--143}
\crow{coetraceeqsum}{\href{https://github.com/lsernau/grh-lean/blob/72fead102f663d5c7a407964a13dc0cf9d977207/grh-compact/src/New/Curve/TightFrameExists.lean\#L12-L13}{{\leanfont Parameters.\allowbreak{}coe\_trace\_eq\_sum}}}{TightFrameExists.lean}{12--13}
\crow{hermmemquadInt}{\href{https://github.com/lsernau/grh-lean/blob/72fead102f663d5c7a407964a13dc0cf9d977207/grh-compact/src/New/Curve/TightFrameExists.lean\#L15-L23}{{\leanfont Parameters.\allowbreak{}herm\_mem\_quadInt}}}{TightFrameExists.lean}{15--23}
\crow{adeleExtensionalpha}{\href{https://github.com/lsernau/grh-lean/blob/72fead102f663d5c7a407964a13dc0cf9d977207/grh-compact/src/New/Curve/TightFrameExists.lean\#L29-L31}{{\leanfont Parameters.\allowbreak{}adeleExtension\_alpha}}}{TightFrameExists.lean}{29--31}
\crow{normiotaperiodFactor}{\href{https://github.com/lsernau/grh-lean/blob/72fead102f663d5c7a407964a13dc0cf9d977207/grh-compact/src/New/Curve/TightFrameExists.lean\#L36-L43}{{\leanfont Parameters.\allowbreak{}norm\_iota\_periodFactor}}}{TightFrameExists.lean}{36--43}
\crow{normperiodFactormulconjleone}{\href{https://github.com/lsernau/grh-lean/blob/72fead102f663d5c7a407964a13dc0cf9d977207/grh-compact/src/New/Curve/TightFrameExists.lean\#L45-L48}{{\leanfont Parameters.\allowbreak{}norm\_periodFactor\_mul\_conj\_le\_one}}}{TightFrameExists.lean}{45--48}
\crow{normiotaperiodFactoreqone}{\href{https://github.com/lsernau/grh-lean/blob/72fead102f663d5c7a407964a13dc0cf9d977207/grh-compact/src/New/Curve/TightFrameExists.lean\#L52-L55}{{\leanfont Parameters.\allowbreak{}norm\_iota\_periodFactor\_eq\_one}}}{TightFrameExists.lean}{52--55}
\crow{isTightFramezetapow}{\href{https://github.com/lsernau/grh-lean/blob/72fead102f663d5c7a407964a13dc0cf9d977207/grh-compact/src/New/Curve/TightFrameExists.lean\#L59-L66}{{\leanfont Parameters.\allowbreak{}isTightFrame\_zeta\_pow}}}{TightFrameExists.lean}{59--66}
\crow{existscyclotomicdifferentfactor}{\href{https://github.com/lsernau/grh-lean/blob/72fead102f663d5c7a407964a13dc0cf9d977207/grh-compact/src/New/Curve/TightFrameExists.lean\#L71-L85}{{\leanfont Parameters.\allowbreak{}exists\_cyclotomic\_different\_factor}}}{TightFrameExists.lean}{71--85}
\crow{existsisTightFrame}{\href{https://github.com/lsernau/grh-lean/blob/72fead102f663d5c7a407964a13dc0cf9d977207/grh-compact/src/New/Curve/TightFrameExists.lean\#L90-L148}{{\leanfont Parameters.\allowbreak{}exists\_isTightFrame}}}{TightFrameExists.lean}{90--148}}
% the notations of the defined functions: one slot for each explicit argument of the Lean declaration
\newcommand{\Rais}[4]{\mathrm{Raise}_{#1}(#2,#3;#4)}
\newcommand{\packing}[4]{(#3,#4)\text{ is a Frobenius packing of }#2\text{ at }#1}
\newcommand{\dl}[2]{\mathrm{dl}_{#1}(#2)}
\newcommand{\nullint}[6]{#6\text{ is integral on the null cone of }(#1,#2;#3,#4)\text{ up to degree }#5}
\newcommand{\vsc}[2]{\mathrm{vsc}(#1,#2)}
\newcommand{\intvac}[1]{\theta^{\circ}_{#1}}
\newcommand{\Lat}[2]{\mathrm{Lat}(#1;#2)}
\newcommand{\Std}[2]{\mathrm{Std}_{#2}(#1)}
\newcommand{\ask}[4]{\mathrm{ask}_{#1,#2}(#3,#4)}
\newcommand{\forced}[2]{\mathrm{forced}(#1,#2)}
\newcommand{\smallns}[2]{#2\text{ is small and not split for }#1}
\newcommand{\shortp}[2]{#2\text{ is short for }#1}
\newcommand{\loss}[2]{\mathrm{loss}(#1,#2)}
\newcommand{\clm}[2]{\mathrm{lg}(#1,#2)}
\newcommand{\deltaP}[2]{\Delta(#1,#2)}
\newcommand{\jc}[2]{\mathrm{jc}(#1,#2)}
\newcommand{\ozero}[2]{o_0(#1,#2)}
\newcommand{\served}[2]{#2\text{ is served for }#1}
\newcommand{\collision}[5]{(#4,#5)\text{ is a collision of }#2\text{ outside }#3\text{ for }#1}
\newcommand{\kap}[3]{\varkappa_{#1}(#2,#3)}
\newcommand{\last}[3]{\mathrm{last}_{#1}(#2,#3)}
\newcommand{\kept}[2]{\mathrm{kept}_{#1}(#2)}
\newcommand{\kapY}[3]{\varkappa^{Y}_{#1}(#2,#3)}
\newcommand{\szero}[2]{s_0(#1,#2)}
\newcommand{\cmE}[1]{\mathcal{C}(#1)}
\newcommand{\Omz}[1]{\Omega_0(#1)}
\newcommand{\asc}[1]{\mathrm{asc}(#1)}
\newcommand{\sone}[2]{s_1(#1,#2)}
\newcommand{\per}[2]{\mathrm{per}(#1,#2)}
\newcommand{\ThProd}[5]{\mathrm{ThProd}_{#1,#2}(#3,#4;#5)}
\newcommand{\UnitMembers}{89}
\newcommand{\UnitInNote}{89}
\newcommand{\UnitProofs}{36}
\newcommand{\UnitSteps}{407}
\newcommand{\UnitAllSteps}{435}


\linespread{1.08}                       % 11 pt on 14.7 pt
\setlength{\parindent}{0pt}
\setlength{\parskip}{5.5pt plus 1.5pt minus 0.5pt}
\setlength{\abovedisplayskip}{7.5pt plus 2pt minus 3pt}
\setlength{\belowdisplayskip}{7.5pt plus 2pt minus 3pt}
\setlength{\abovedisplayshortskip}{3pt plus 2pt}
\setlength{\belowdisplayshortskip}{6pt plus 2pt minus 2pt}
\emergencystretch=1.5em
\clubpenalty=10000 \widowpenalty=10000 \displaywidowpenalty=10000
\relpenalty=7000 \binoppenalty=9000
\raggedbottom
\pagestyle{fancy}\fancyhf{}\renewcommand{\headrulewidth}{0pt}
\fancyfoot[C]{\sffamily\fontsize{7.6}{9}\selectfont\color{inkfaint}\thepage\,/\,\pageref*{LastPage}}
\setlist{topsep=2pt,partopsep=0pt,parsep=0pt,itemsep=4pt,leftmargin=2.1em,labelsep=0.55em}

\newcommand{\Q}{\mathbb{Q}}
\newcommand{\Z}{\mathbb{Z}}
\newcommand{\F}{\mathbb{F}}
\newcommand{\len}{\operatorname{len}}
\newcommand{\cok}{\operatorname{cok}}
\newcommand{\End}{\operatorname{End}}
\newcommand{\Hom}{\operatorname{Hom}}
\newcommand{\Gal}{\operatorname{Gal}}
\newcommand{\diag}{\operatorname{diag}}
\newcommand{\kw}[1]{\emph{#1}}
\newcommand{\m}[1]{\mbox{$#1$}}        % a short formula that must not be broken across lines
\newcommand{\head}[1]{\par\addvspace{17pt plus 4pt minus 2pt}\Needspace{6\baselineskip}%
  \noindent{\fontsize{13.2}{16}\selectfont\bfseries\color{ink}#1}\par\nopagebreak\vspace{-1pt}}
\newcommand{\sect}[2]{\refstepcounter{section}\label{#2}\par\addvspace{19pt plus 4pt minus 2pt}\Needspace{7\baselineskip}%
  \noindent\llap{{\fontsize{13.2}{16}\selectfont\bfseries\color{inkfaint}\thesection}\hspace{9pt}}%
  {\fontsize{13.2}{16}\selectfont\bfseries\color{ink}#1}\par\nopagebreak\vspace{-1pt}}
\newcommand{\subhead}[1]{\par\addvspace{9pt plus 2pt}\Needspace{4\baselineskip}\noindent\textbf{#1}\enspace\ignorespaces}
% where a statement lives in the Lean archive
\newcommand{\leantag}{{\sffamily\bfseries\fontsize{6.3}{8}\selectfont\color{leancol}\addfontfeatures{LetterSpace=7}LEAN}}
\newcommand{\sep}{\unskip\nobreak\hspace{5pt}{\color{inkfaint}\textperiodcentered}\hspace{5pt plus 2pt}\ignorespaces}
\newcommand{\leanline}[1]{\par\nopagebreak\vspace{-2.5pt}{\fontsize{8.9}{12.4}\selectfont\color{inksoft}\raggedright
  \noindent\llap{\leantag\hspace{9pt}}#1\par}\vspace{1pt}}
% a displayed statement in a tinted box: #1 colour, #2 text
\newcommand{\statement}[2]{\par\addvspace{7pt}\begingroup\setlength{\fboxsep}{8.5pt}\setlength{\fboxrule}{0.6pt}%
  \noindent\fcolorbox{#1!30!white}{#1!3!white}{\parbox{\dimexpr\linewidth-2\fboxsep-2\fboxrule\relax}{#2}}%
  \endgroup\par\addvspace{7pt}}
\newcommand{\nt}[2]{\par\noindent\hangindent=1.1em\hangafter=1 #1\enspace{\color{inksoft}#2}\par\vspace{1pt}}
\newcommand{\legdot}[1]{\tikz[baseline=-0.62ex]\fill[#1] (0,0) circle (2.5pt);\hspace{3pt}}
\newcommand{\crow}[4]{\csname CS@#1\endcsname & #2 & \textcolor{inksoft}{#3} & \textcolor{inksoft}{#4}\\}
\newcommand{\leanquoteline}[2]{\par\noindent\hangindent=\dimexpr#1\fontdimen2\font+2.2em\relax\hangafter=1
  \hspace*{#1\fontdimen2\font}#2}
\newcommand{\quoted}[1]{\par\vspace{3pt}{\ttfamily\fontsize{7.5}{10.4}\selectfont\color{leancol}\raggedright
  \setlength{\parskip}{0pt}\leanquote{#1}\par}\vspace{4pt}}
\newcommand{\bx}[3]{{\sffamily\bfseries\fontsize{7.3}{9.2}\selectfont\color{#1!82!black}#2}\\[1.5pt]#3}
\newcommand{\cs}[2]{\expandafter\newcommand\csname CS@#1\endcsname{#2}}

% theorems
\newtheoremstyle{plainx}{7pt plus 2pt}{7pt plus 2pt}{\itshape}{}{\bfseries}{.}{0.5em}{}
\newtheoremstyle{uprightx}{7pt plus 2pt}{7pt plus 2pt}{}{}{\bfseries}{.}{0.5em}{}
\theoremstyle{plainx}
\newtheorem{theorem}{Theorem}[section]
\newtheorem{lemma}[theorem]{Lemma}
\newtheorem{proposition}[theorem]{Proposition}
\newtheorem{corollary}[theorem]{Corollary}
\theoremstyle{uprightx}
\newtheorem{definition}[theorem]{Definition}
\newtheorem{setting}[theorem]{Setting}          % hypotheses shared by several statements; each of them cites it by number
\newtheorem{remark}[theorem]{Remark}
\newtheorem{example}[theorem]{Example}
\renewcommand{\qedsymbol}{{\setlength{\fboxsep}{0pt}\setlength{\fboxrule}{0.5pt}\fbox{\rule{0pt}{0.46em}\hspace{0.46em}}}}
% what holds throughout the note: the ambient objects only (implicit and instance arguments), never a hypothesis
\newenvironment{conventions}{\par\addvspace{5pt}\noindent\textbf{Conventions.}\enspace\ignorespaces}{\par\addvspace{3pt}}
\makeatletter
\renewenvironment{proof}[1][Proof]{\par\pushQED{\qed}\normalfont\topsep0pt\partopsep0pt
  \trivlist\item[\hskip\labelsep\itshape #1.]\ignorespaces}{\popQED\endtrivlist\@endpefalse}
\makeatother

% ------------------------------------------------------------------------------- macros of this note
% Notations of the definitions of the unit come from data.tex (made from spec.json). The macros below
% write objects of other units and of the background library, each with all its arguments.
\newcommand{\CC}{\mathbb{C}}
\newcommand{\RR}{\mathbb{R}}
\newcommand{\NN}{\mathbb{N}}
\newcommand{\Cp}{\CC_p}
\newcommand{\op}{\mathfrak{o}_p}
\newcommand{\mxp}{\mathfrak{m}_p}
\newcommand{\Oj}{O_j}
\newcommand{\np}[1]{|#1|_p}
\newcommand{\cj}[1]{\overline{#1}}
\newcommand{\lnm}[1]{{\leanfont #1}}
\newcommand{\Aut}{\operatorname{Aut}}
\newcommand{\Tr}{\operatorname{Tr}}
\newcommand{\Nm}{\operatorname{N}}
\newcommand{\Af}{\mathbb{A}_{\mathrm{f}}}
\newcommand{\Ad}{\mathbb{A}}
\newcommand{\Ohat}{\widehat{\mathcal{O}}}
\newcommand{\OO}{\mathcal{O}}
% the field of the parameters (unit 152; the forms of the note on the CM field)
\newcommand{\KP}[1]{K(#1)}
\newcommand{\EP}[2]{E(#1,#2)}
\newcommand{\Nn}[2]{N_1(#1,#2)}
\newcommand{\PlP}[2]{\operatorname{Pl}(#1,#2)}
\newcommand{\emb}[1]{\iota_{#1}}
\newcommand{\pin}[3]{[#3]_{#1,#2}}
\newcommand{\chiK}[1]{\chi_{\KP{#1}}}
\newcommand{\chid}[1]{\chi_{#1}}
\newcommand{\rdE}[2]{\operatorname{rd}_{\EP{#1}{#2}}}
\newcommand{\OmP}[2]{\Omega(#1,#2)}
\newcommand{\alP}[2]{\alpha(#1,#2)}
\newcommand{\dgen}[2]{\partial(#1,#2)}
\newcommand{\rtu}[2]{\varsigma(#1,#2)}
\newcommand{\czE}[2]{c_0(#1,#2)}
\newcommand{\frl}[2]{\mathfrak{a}(#1,#2)}
\newcommand{\prm}[2]{\operatorname{perm}_{#1}(#2)}
\newcommand{\nsw}[2]{\varrho_{#1}(#2)}
\newcommand{\cP}[1]{\symit{C}(#1)}
\newcommand{\Cgen}[1]{C_{\mathrm{gen}}(#1)}
% the data (unit 152)
\newcommand{\ED}[1]{E(#1)}
\newcommand{\NoneD}[1]{N_1(#1)}
\newcommand{\PlD}[1]{\operatorname{Pl}(#1)}
\newcommand{\fD}[1]{\symit{f}(#1)}
\newcommand{\kD}[1]{\symit{k}(#1)}
\newcommand{\tD}[1]{\symit{t}(#1)}
\newcommand{\aD}[1]{\symit{a}(#1)}
\newcommand{\muD}[1]{\mu(#1)}
\newcommand{\tops}[1]{\mathrm{tops}(#1)}
\newcommand{\thick}[1]{\operatorname{thick}(#1)}
\newcommand{\nomW}[2]{w_{#1}(#2)}
\newcommand{\pinD}[2]{[#2]_{#1}}
\newcommand{\nD}[2]{n_{#2}(#1)}
\newcommand{\wt}[3]{\mathrm{wt}_{#1}(#2,#3)}
\newcommand{\gap}[1]{\mathrm{gap}(#1)}
\newcommand{\btwo}[1]{b_2(#1)}
% three definitions of the unit whose list of arguments in the census has an entry that Lean does not take
% (see the notes on the formal text): written here with exactly the explicit arguments of Lean
\newcommand{\gain}[2]{\mathrm{gain}(#1,#2)}
\newcommand{\coll}[3]{\mathrm{coll}(#1,#2,#3)}
\newcommand{\cc}[7]{\mathrm{cc}_{#1}(#2,#3;#4;#5,#6;#7)}
% the Weil representation (background library)
\newcommand{\Sw}[1]{\mathcal{S}(#1)}
\newcommand{\Wf}[1]{\mathcal{W}_{\mathrm{f}}(#1)}
\newcommand{\vac}[1]{\mathrm{vac}_{#1}}
\newcommand{\heis}[2]{\mathfrak{h}_{#1}(#2)}
\newcommand{\fheis}[2]{\mathfrak{h}^{\mathrm{f}}_{#1}(#2)}
\newcommand{\fv}[1]{#1^{\flat}}
\newcommand{\Th}[1]{\operatorname{Th}_{#1}}
\newcommand{\weil}[1]{\omega_{#1}}
\newcommand{\jet}[3]{T^{#1}_{#2}#3}
\newcommand{\Thf}[1]{\Theta_{#1}}
\newcommand{\inn}[3]{\langle #2,#3\rangle_{#1}}
\newcommand{\name}[3]{\operatorname{name}_{#1,#2}(#3)}
\newcommand{\fmult}[1]{\Delta_{\mathrm{f}}(#1)}
\newcommand{\Uf}[1]{\mathrm{U}_{#1}(\Af)}
\newcommand{\stdsp}[1]{I_{#1}}
\newcommand{\homsp}[2]{\operatorname{Hom}(\stdsp{#1},\stdsp{#2})}
% a CM datum and its functions (background library)
\newcommand{\LamC}[1]{\Lambda_{#1}}
\newcommand{\thC}[1]{\theta_{#1}}
\newcommand{\xiC}[1]{\xi_{#1}}
\newcommand{\phC}[1]{\varphi_{#1}}
\newcommand{\UC}[3]{U^{#1,#2}_{#3}}
\newcommand{\tay}[1]{\operatorname{Tay}(#1)}
\newcommand{\tayj}[1]{\operatorname{Tay}_j(#1)}
\newcommand{\phiring}[1]{\op[[\phC{#1}]]}
\newcommand{\area}[1]{\operatorname{area}(#1)}
\newcommand{\mon}[2]{\operatorname{mon}_{#1}(#2)}
\newcommand{\cpart}[3]{[#3]_{#2\bmod #1}}
% forms, gauges, kernels (background library)
\newcommand{\gauge}[4]{\|#4\|_{#1,#2;#3}}
\newcommand{\boxes}[2]{\operatorname{box}(#1,#2)}
\newcommand{\schurjet}[2]{\operatorname{SJ}_{#1}(#2)}
\newcommand{\jetalong}[2]{\operatorname{jet}_{#1}(#2)}
\newcommand{\thker}[5]{\mathrm{k}_{#1}(#2;#3,#4;#5)}
\newcommand{\pairw}[2]{\omega_{#1,#2}}

% where each declaration of the concordance is treated: \cs{key}{place}
\cs{vacuumRoot}{Definition~\ref{d:period}}
\cs{cmCurve}{Definition~\ref{d:period}}
\cs{Omz}{Definition~\ref{d:period}}
\cs{areaScale}{Definition~\ref{d:period}}
\cs{periodFactor}{Definition~\ref{d:period}}
\cs{jetScale}{Definition~\ref{d:period}}
\cs{vacuumScale}{Definition~\ref{d:vac}}
\cs{intVacuum}{Definition~\ref{d:vac}}
\cs{integralAfter}{Definition~\ref{d:vac}}
\cs{standardData}{Definition~\ref{d:vac}}
\cs{raises}{Definition~\ref{d:raises}}
\cs{thetaProducts}{Definition~\ref{d:thprod}}
\cs{IsIntegralOnNullCone}{Definition~\ref{d:null}}
\cs{localAsk}{Definition~\ref{d:ask}}
\cs{forcedPart}{Definition~\ref{d:ask}}
\cs{gain}{Definition~\ref{d:ask}}
\cs{loss}{Definition~\ref{d:ask}}
\cs{IsSmallNonsplit}{Definition~\ref{d:ask}}
\cs{IsShortPrime}{Definition~\ref{d:ask}}
\cs{ceilLogM}{Definition~\ref{d:served}}
\cs{deltaP}{Definition~\ref{d:served}}
\cs{jetCeil}{Definition~\ref{d:served}}
\cs{ozero}{Definition~\ref{d:served}}
\cs{IsServed}{Definition~\ref{d:served}}
\cs{IsCollision}{Definition~\ref{d:served}}
\cs{collisionSum}{Definition~\ref{d:served}}
\cs{couplingCount}{Definition~\ref{d:cc}}
\cs{kap}{Definition~\ref{d:kappa}}
\cs{lastEntry}{Definition~\ref{d:kappa}}
\cs{keptColumns}{Definition~\ref{d:kappa}}
\cs{columnsY}{Definition~\ref{d:kappa}}
\cs{IsFrobPacking}{Definition~\ref{d:packing}}
\cs{depthLossInt}{Definition~\ref{d:packing}}
\cs{normgalvacuumRoot}{Lemma~\ref{l:period}(a)}
\cs{bigOmegapos}{Lemma~\ref{l:period}(b)}
\cs{jetScalenezero}{Lemma~\ref{l:period}(c)}
\cs{latticeModulecmCurve}{Lemma~\ref{l:period}(d)}
\cs{omegamemcmCurvelattice}{Lemma~\ref{l:period}(e)}
\cs{adeleExtensionalpha}{Lemma~\ref{l:padic}(a)}
\cs{normiotaperiodFactor}{Lemma~\ref{l:padic}(b)}
\cs{normperiodFactormulconjleone}{Lemma~\ref{l:padic}(c)}
\cs{normiotaperiodFactoreqone}{Lemma~\ref{l:padic}(d)}
\cs{coetraceeqsum}{Lemma~\ref{l:trace}(a)}
\cs{hermmemquadInt}{Lemma~\ref{l:trace}(b)}
\cs{isTightFramezetapow}{Lemma~\ref{l:zeta}}
\cs{existscyclotomicdifferentfactor}{Lemma~\ref{l:cyclo}}
\cs{existsisTightFrame}{Theorem A}
\cs{taylorlineeqsubstjetSeries}{Lemma~\ref{l:line}(a)}
\cs{eqzeroofjetSerieseqzero}{Lemma~\ref{l:line}(b)}
\cs{jetSerieseqprod}{Lemma~\ref{l:line}(c)}
\cs{existsfinitePartofisTightFrame}{Proposition~\ref{p:F}}
\cs{kappanezero}{Lemma~\ref{l:vac}(a)}
\cs{intVacuumfinVeceqprod}{Lemma~\ref{l:vac}(b)}
\cs{normintVacuumle}{Lemma~\ref{l:vac}(c)}
\cs{normintVacuumleone}{Lemma~\ref{l:vac}(d)}
\cs{normdataInnerle}{Lemma~\ref{l:vac}(e)}
\cs{existsintVacuumnezero}{Lemma~\ref{l:vac}(f)}
\cs{isAdmissibleCombname}{Lemma~\ref{l:adm}}
\cs{existsisJetConjugategenerator}{Proposition~\ref{p:gen}}
\cs{existssemilinearisJetConjugate}{Theorem B}
\cs{ofintegral}{Lemma~\ref{l:std}(a)}
\cs{finHeisOpmemstandardData}{Lemma~\ref{l:std}(b)}
\cs{omegamemstandardData}{Lemma~\ref{l:std}(c)}
\cs{normmappsi}{Lemma~\ref{l:std}(d)}
\cs{pairWeilmemstandardData}{Lemma~\ref{l:std}(e)}
\cs{normjetSeriesle}{Proposition~\ref{p:transfer}}
\cs{jetthetaprodintegral}{Lemma~\ref{l:jetint}}
\cs{normjetjetSeriesleone}{Theorem C}
\cs{smulmemraises}{Proposition~\ref{p:raises}(a)}
\cs{memraisesofjetbound}{Proposition~\ref{p:raises}(b)}
\cs{memraisesofstandardData}{Proposition~\ref{p:raises}(c)}
\cs{memraisesofdvd}{Proposition~\ref{p:raises}(d)}
\cs{existsnormeqinvpowbounds}{Lemma~\ref{l:val}}
\cs{coeffsubsttangent}{Lemma~\ref{l:tangent}}
\cs{normcoeffleofderiv}{Lemma~\ref{l:deriv}}
\cs{existsformalparameter}{Proposition~\ref{p:formal}}
\cs{sumpoweqzeroofnull}{Lemma~\ref{l:charp}(a)}
\cs{coeffsumfrobnezeroiff}{Lemma~\ref{l:charp}(b)}
\cs{existsframerowsfrob}{Lemma~\ref{l:frob}}
\cs{thetaProductsnormalform}{Theorem D}
\cs{existsnullcouplingframes}{Lemma~\ref{l:frames}}
\cs{normschurJetle}{Theorem E}
\cs{katzdivisibilitylimit}{Proposition~\ref{p:katz}}
\cs{hasSumliveinvpow}{Lemma~\ref{l:live}(a)}
\cs{summablelocalAsk}{Lemma~\ref{l:live}(b)}
\cs{nedofplaceClasseq}{Lemma~\ref{l:live}(c)}
\cs{kaple}{Lemma~\ref{l:live}(d)}
\cs{existsisFrobPacking}{Proposition~\ref{p:packing}}
\cs{existshasLargeSplitofpairs}{Theorem F}

\begin{document}
\color{ink}

{\raggedright
{\sffamily\bfseries\fontsize{7.2}{9}\selectfont\color{leancol}\addfontfeatures{LetterSpace=8}READBACK NOTE 20 \textperiodcentered{} CHAPTER V, THE CURVE\par}
\vspace{5pt}
{\fontsize{21.5}{25.5}\selectfont\bfseries Theta functions of a CM elliptic curve: periods, tight frames, conjugates and divisibility of jets\par}
\vspace{5pt}
{\fontsize{12.4}{16}\selectfont\itshape\color{inksoft} One unit of the Lean artifact \ArtifactName, read back into
mathematics\par}}
\vspace{8pt}
{\color{rulegrey}\hrule height 0.6pt}
\vspace{7pt}
\thispagestyle{empty}

\textbf{Summary.} On the definite unitary groups of Theorem~A of the write-up, a theta lift is a finite sum of
Taylor coefficients at the origin (\emph{jets}) of theta functions. This unit shows that these theta functions are
products of translates of one function, the theta function $\thC{C}$ of an elliptic curve with complex multiplication
by $K=\Q(\sqrt{-d})$, composed with linear forms given by a \emph{tight frame} in the CM field $E$ (Theorem~A,
Proposition~\ref{p:F}). Everything else is then read off the curve: the conjugates of the jets under
$\Aut(\CC/K)$ are jets again, of bounded size (Theorem~B); suitably scaled jets are integral at every prime
(Theorem~C); at a prime $p$ that is inert or ramified in $K$ a product of theta series is, after a change of scale,
an integral series times one exponential (Theorem~D), which makes jets along products of minors divisible by an
explicit power of $p$ (Theorem~E), and this power cannot be improved by much (Theorem~F). The unit also defines the
numbers that the next unit uses to state the divisibility of the whole chain of lifts: the ask, the gain, the loss
and the column counts.

\statement{leancol}{\fontsize{9.6}{13.2}\selectfont\textbf{What this note is.} The artifact \ArtifactName{} is a
body of Lean code: the compact edition of a formal proof, attributed to L.~Alpöge with the AI system Claude, whose
final theorems state the Riemann hypothesis and its analogue for Dirichlet $L$-functions. That headline is a
\emph{claim}: the authors report that every file compiles and that Lean's kernel accepted the proofs; the record of
checks of the handout says who replayed what; nothing was compiled for this note; no mathematician has read the
argument through, and nothing here says that the Riemann hypothesis is proved. This note treats one unit, the
folder on the CM curve: all \UnitInNote{} of its \UnitMembers{} declarations. Every statement below was compared
with its Lean text, argument by argument. Proofs are given in full for what the unit itself proves; results of
other units are stated as \emph{Inputs} with the unit that owns them, and results of the background library of the
artifact (complex multiplication, the Chowla--Selberg formula, the formal group of the curve, the Weil
representation) are \emph{quoted}, with a reference where the catalogue of the handout has one. Where a formal
step is closed by automation that I could not follow without running Lean, §\ref{s:notes} says so. The Lean text
read is \PinText; Appendix~A is the concordance of \ConcDecls{} declarations.}

\sect{The statements}{s:intro}
The notation is that of §\ref{s:setting}. $P=(N_0,d)$ are parameters and $\mu$ is a natural number; $\EP{P}{\mu}$
is the CM field obtained from $\KP{P}=\Q(\sqrt{-d})$ by adjoining a root of unity of order $\Nn{P}{\mu}$,
$\PlP{P}{\mu}$ is the group of its infinite places, and $\emb{\tau}\colon\EP{P}{\mu}\to\CC$ is the chosen embedding
at the place $\tau$. A \emph{$p$-adic place} is a prime number $p$ with a ring homomorphism $j\colon\CC\to\Cp$; $\np{\,\cdot\,}$ is the
absolute value of $\Cp$ and $\op$ its ring of integers. The numbers $\sone{P}{\mu}$ and $\per{P}{\mu}$ are the
period factor and the scale of the jets (Definition~\ref{d:period}), $\jet{V}{a}{f}$ is the series of jets of a
vector $f$ of the Weil representation of a hermitian space $V$ at the scale $a$ (Background~4), and $\Std{V}{\delta}$
is the module of standard data (Definition~\ref{d:vac}).

\statement{arith}{\phantomsection\label{t:A}\textbf{Theorem A (tight frames with a unit vector).}\enspace Let $P=(N_0,d)$
be parameters, $\mu$ a natural number, $p$ a prime number and $j\colon\CC\to\Cp$ a ring homomorphism. Let $M$ and
$c_\ast$ be natural numbers with $(d\cdot\Nn{P}{\mu})^6\,(c_\ast+1)\le M$. Then there are a finite set $I$, elements
$v_i\in\EP{P}{\mu}$ and natural numbers $c_i$, for $i\in I$, such that the family $(v_i)_{i\in I}$ with the
multiplicities $(c_i)_{i\in I}$ is a tight frame of length $M$ (Background~1), $c_\ast\le c_i$ for every $i\in I$,
and for some $i\in I$
\[\np{j\bigl(\sone{P}{\mu}\cdot\cj{\emb{\tau}(v_i)}\bigr)}=1\qquad\text{for every }\tau\in\PlP{P}{\mu}.\]}
\leanline{\lean{existsisTightFrame}}

\statement{arith}{\phantomsection\label{t:B}\textbf{Theorem B (conjugates of jets).}\enspace Let $P=(N_0,d)$ be
parameters, $\mu$ a natural number, $n$ a finite set and $V$ a standard hermitian space over $\EP{P}{\mu}$ with
index set $n$. Let $R$ be a set of finite places of $\EP{P}{\mu}$ that contains the complex conjugate of none of
its members, and let $\sigma\in\Aut(\CC/\KP{P})$. Then there are $t\in\CC$, $c\in\CC$, a finite adele $z$ of
$\EP{P}{\mu}$ with $z\cj{z}=1$ and a $\sigma$-semilinear map $L\colon\Sw{V}\to\Sw{V}$ with values in $\Wf{V}$ such
that:
\begin{enumerate}[label=(\roman*),itemsep=1pt]
\item $|t|^2=\OmP{P}{\mu}^{-1}$, and $c\ne0$;
\item for every $f\in\Wf{V}$, applying $\sigma$ to the coefficients of $\jet{V}{\per{P}{\mu}}{f}$ gives the series
$\jet{V}{t}{(Lf)}$ with every variable $z_{(\tau,i)}$ replaced by $z_{(\pi^{-1}\tau,i)}$, where
$\pi=\prm{\mu}{\sigma}$ is the permutation of $\PlP{P}{\mu}$ with $\sigma\circ\emb{\tau}=\emb{\pi(\tau)}$;
\item for every $f\in\Wf{V}$ and every $g\in\Uf{n}$:
$L\bigl(\weil{V}(1,g)f\bigr)=\dfrac{\sigma(\fmult{g})}{\fmult{g}}\cdot\weil{V}(1,g)(Lf)$;
\item for every $x\in\Af^{\,n}$, every $g\in\Uf{n}$ and every adelic vector $y\in\Ad^{n}$:
\[\bigl|L\bigl(\heis{V}{\fv{x}}\,\weil{V}(1,g)\,\intvac{V}\bigr)(y)\bigr|\ \le\ \exp\bigl(\cP{P}\cdot(\Nn{P}{\mu}\,|n|+1)^2\bigr);\]
\item for every Schwartz--Bruhat function $\Psi$ above $R$:
$L\bigl(\name{V}{R}{\Psi}\vac{V}\bigr)=c\cdot\weil{V}(1,z\cdot 1_n)\bigl(\name{V}{R}{\sigma\circ\Psi}\vac{V}\bigr)$.
\end{enumerate}}
\leanline{\lean{existssemilinearisJetConjugate}}

\statement{arith}{\phantomsection\label{t:C}\textbf{Theorem C (integrality of jets).}\enspace Let $P=(N_0,d)$ be
parameters, $\mu$ a natural number, $n$ a finite set, $V$ a standard hermitian space over $\EP{P}{\mu}$ with index
set $n$, $p$ a prime number and $j\colon\CC\to\Cp$ a ring homomorphism. Let $\delta$ be a nonzero integer of
$\KP{P}$ and $f\in\Std{V}{\delta}$. Let $Q$ be a polynomial in the variables $z_{(\tau,i)}$,
$(\tau,i)\in\PlP{P}{\mu}\times n$, with coefficients in $\Oj$, and $Q'$ the same polynomial with its coefficients
read in $\CC$. Then
\[\np{j\Bigl(\jetalong{Q'}{\jet{V}{\cj{\delta}\cdot d\cdot\per{P}{\mu}}{f}}\Bigr)}\ \le\ 1 ,\]
where $\jetalong{Q'}{F}=(Q'(\partial)F)(0)$ is the jet of a power series $F$ along $Q'$.}
\leanline{\lean{normjetjetSeriesleone}}

\statement{arith}{\phantomsection\label{t:D}\textbf{Theorem D (normal form of theta products at a non-split prime).}\enspace
Let $P=(N_0,d)$ be parameters, $\mu$ a natural number, $p$ a prime number with $\chiK{P}(p)\ne1$ and
$j\colon\CC\to\Cp$ a ring homomorphism. Then there are $\varpi\in\Cp$ and $\pi\in\op$ with
\[\np{\varpi}=p^{-\nsw{d}{p}},\qquad\np{\pi}=p^{(\chiK{P}(p)-1)/2},\]
such that for every finite set $I$, every real number $r\le1$ and every $\Theta\in\ThProd{P}{\mu}{1}{r}{I}$:
\begin{enumerate}[label=(\alph*),itemsep=1pt]
\item for every natural number $B$,
\[\nullint{\mathrm{pr}}{S_B}{\varpi}{\pi}{B}{\Theta},\]
where $\mathrm{pr}_e(v)=\pin{P}{\mu}{-p^e}\cdot v$ and $S_B$ is the set of the natural numbers
$e\le\lfloor\log_pB\rfloor$ such that $e$ is odd or $p=d$;
\item if $\chiK{P}(p)=-1$ and $r=1$, there are $\kappa\in\Cp$ and power series $\Psi$, $V$ with coefficients in
$\op$, in the variables $z_x$, $x\in\PlP{P}{\mu}\times I$, such that $\np{\kappa}=p^{1-1/(p-1)}$, $\Psi(0)=0$,
\[\Theta(\varpi z)=\exp(\kappa\Psi)\cdot V ,\]
and for every variable $x$ and every exponent $\beta$: the coefficient at $z^\beta$ of the reduction of
$\partial\Psi/\partial z_x$ modulo $\mxp$ is nonzero if and only if $\beta=p^e\cdot 1_y$ for an odd $e$ and a variable
$y=(v,i)$ with $(\pin{P}{\mu}{-p^e}\cdot v,\,i)=x$; here $1_y$ is the exponent of $z_y$.
\end{enumerate}}
\leanline{\lean{thetaProductsnormalform}}

\statement{arith}{\phantomsection\label{t:E}\textbf{Theorem E (divisibility of Schur jets).}\enspace Let $p$ be a prime
number, $\Phi$, $n$, $n'$ finite sets, and $\varkappa\colon\Phi\times\NN\to\NN$ a family of column counts that fits
the ranks $|n|$, $|n'|$ (Background~5). Let $\mathrm{pr}\colon\NN\times\Phi\to\Phi$ be a map, $S\subset\NN$ and
$\Xi\subset\Phi$ finite sets, $\rho$, $\omega$ real numbers, $\varpi\in\Cp$ with $\np{\varpi}=p^{-\rho}$,
$\pi\in\op$ with $\np{\pi}=p^{\omega}$, and $\Theta$ a power series over $\Cp$ in the variables $z_{(\tau,b,c)}$,
$(\tau,b,c)\in\Phi\times n'\times n$. Let $g\colon\Phi\times\NN\times n'\to\NN$ be a function (the \emph{charges})
such that for all $e\in S$, $\tau\in\Phi$, $h\in\NN$ and $b\in n'$:
\[\tau\in\Xi\ \text{ or }\ \bigl(h=1\text{ and }\mathrm{pr}_e(\tau)\in\Xi\bigr)\quad\Longrightarrow\quad
g(\tau,1,b)\ne0\ \text{ or }\ g(\mathrm{pr}_e(\tau),h,b)\ne0 .\]
Let $i'$ and $i$ be Weyl indices for $\varkappa$ of ranks $|n'|$ and $|n|$, and put $N=|n'|$, $B=\boxes{\varkappa}{N}$
and $m=\sum_{\tau\in\Phi}\sum_{h<N}\sum_{l<\varkappa(\tau,h)}\sum_{b<h}g\bigl(\tau,h,i'_\tau(h,l,b)\bigr)$, the total charge
of the boxes of $i'$. Suppose that
\[\nullint{\mathrm{pr}}{S}{\varpi}{\pi}{B}{\Theta}.\]
Then
\[\np{\schurjet{\varkappa}{\Theta}_{i',i}}\ \le\ p^{\,m-\bigl((\frac1{p-1}-\rho)B+\omega\sum_{e\in S}\sum_{\tau\in\Phi}
\cc{p}{\varkappa}{N}{\Xi}{\tau}{\mathrm{pr}_e(\tau)}{e}-|\Phi|\,N^3(\lfloor\log_pB\rfloor+1)\bigr)} .\]}
\leanline{\lean{normschurJetle}}

\statement{arith}{\phantomsection\label{t:F}\textbf{Theorem F (a jet that is not more divisible).}\enspace Let $P=(N_0,d)$
be parameters, $\mu$ a natural number, $p$ a prime number with $\chiK{P}(p)\ne1$ and $j\colon\CC\to\Cp$ a ring
homomorphism. Let $\mathsf{N}$ be a function from the exponents in the variables $z_\tau$, $\tau\in\PlP{P}{\mu}$,
to the non-negative real numbers, $\Theta$ a power series over $\Cp$ in these variables, $n$ an exponent, and $i$, $y$
functions from $\PlP{P}{\mu}$ to $\NN$. Suppose that $\packing{p}{n}{i}{y}$; that $\Theta(0)=1$; that $\Theta$ is
carried by one copy and $\Theta^{-1}$ by five copies of $\mathsf{N}$ (Background~6); and that
$\Theta\in\ThProd{P}{\mu}{1}{1}{\{\ast\}}$. Then there are exponents $\beta\le n$ and $\gamma\le n$ with
$|\beta|=\sum_\tau y_\tau(p^{i_\tau}+1)$ and $|\gamma|=\sum_\tau y_\tau$ such that $\beta$ is a sum of six exponents
$\alpha_1,\dots,\alpha_6$ with
\[\prod_{k=1}^{6}\mathsf{N}(\alpha_k)\ \ge\ p^{\,(1-\frac1{p-1})|\gamma|+\nsw{d}{p}\,|\beta|+\sum_{\tau}v_p(\gamma_\tau!)} .\]}
\leanline{\lean{existshasLargeSplitofpairs}}

\textbf{How the pieces fit} (motivation, read from the comments of the formal text and from the places where later
units use the results; not statements of this unit). A column on the definite group is lifted by integrating it
against a theta kernel, and on a definite group this integral is a finite sum of matrices of jets
$\schurjet{\varkappa}{\jet{V}{c}{f}}$ of data $f$ (Background~5). Theorem~A and Proposition~\ref{p:F} turn the series
$\jet{V}{c}{f}$ into products of Taylor series of translates of $\thC{C}$; the scale $\per{P}{\mu}$ is chosen so
that these products have $p$-adic absolute values governed by the period factor $\sone{P}{\mu}$ and archimedean
conjugates governed by $\OmP{P}{\mu}$. Theorem~C says that one lift cannot lower $p$-adic valuations by more than
a power of $\cj{\delta}d$ (Proposition~\ref{p:raises}), which is all one needs at the primes that split in $K$. At the
other primes the formal group of the curve has height two. There the logarithm of $\thC{C}$ keeps a denominator
only in the degrees $p+1$ and $p+q$ of a good parameter (Proposition~\ref{p:formal}); a derivative of order $a$
gains about $a/(p-1)$ powers of $p$, and the exponential $\exp(\kappa\Psi)$ takes back what the couplings of a
variable at $\tau$ with the variables at the partner places $\pin{P}{\mu}{-p^e}\cdot\tau$ allow. The balance is the
number $\ask{D}{p}{y}{\tau}$ of Definition~\ref{d:ask}, summed in the gain $\gain{D}{p}$; it is the local form, at
$p$, of the weight $\lambda_{d,1}$ of the prime sum of the main inequality. Theorem~F shows that the count is
sharp, which the unit on the rows needs for a lower bound.

\sect{The setting}{s:setting}
\begin{conventions}\phantomsection\label{c:conventions}
$\NN=\{0,1,2,\dots\}$; subtraction of natural numbers is truncated at $0$ and division of natural numbers is the
floor, as in the formal text. \emph{Parameters} are a pair $P=(N_0,d)$: a positive multiple $N_0$ of $18480$ and a
prime number $d\equiv N_0/2-1\bmod N_0$; the \emph{exponent} is a natural number $\mu$. A statement about
\emph{data} concerns admissible data $D$ for parameters $P$, with components $x$, $x_1$, $b$ (the three scales),
$m$, $r$ and the exponent $\muD{D}$. Index sets such as $n$, $n'$, $I$, $J$, $\Phi$ are finite sets; a hermitian
space $V$ has an index set $n$ and signs $\varepsilon\colon n\to\{1,-1\}$. In a statement about a prime number
$p$ and a ring homomorphism $j\colon\CC\to\Cp$, the absolute value $\np{\,\cdot\,}$ of $\Cp$ is normalised by
$\np{p}=p^{-1}$, $\op=\{x\in\Cp:\np{x}\le1\}$, $\mxp$ is its maximal ideal, and $\Oj=\{z\in\CC:\np{j(z)}\le1\}$;
such a pair $(p,j)$ is called a $p$-adic place. In a statement about a field $E$ with a CM type and Weil data,
about a commutative ring of characteristic $p$, about a normed field, or about an integral domain of
characteristic $0$, these structures are part of the statement. $v_p$ is the $p$-adic valuation of a natural
number. For a multi-index (an \emph{exponent}) $\beta$, $|\beta|$ is its degree, $\beta!$ the product of the
factorials of its entries, and $\beta\le\beta'$ is meant entry by entry.
\end{conventions}

The unit rests on four bodies of earlier material. They are recalled in six numbered paragraphs, after the
definitions and the comments of the formal text; the notes of the units that own them give the details.

\subhead{Background 1: the field, its lattice, tight frames} (unit~152, the note on the CM field; the definition of
a frame is in the background library). For parameters $P=(N_0,d)$ and $\mu\in\NN$: $\KP{P}=\Q(\sqrt{-d})\subset\CC$
with ring of integers $\OO_{\KP{P}}$; $\chiK{P}$ is its quadratic character, on integers the Legendre symbol
$\chid{d}$ modulo $d$; $\Nn{P}{\mu}$ is $N_0$ times a power of $3$ and a power of $2$ fixed by $\mu$;
$\EP{P}{\mu}=\KP{P}(\rtu{P}{\mu})$ with $\rtu{P}{\mu}$ a primitive root of unity of order $\Nn{P}{\mu}$. The set
$\PlP{P}{\mu}$ of infinite places of $\EP{P}{\mu}$ is a group, identified with $\operatorname{Gal}(\EP{P}{\mu}/\KP{P})$ and with
$(\Z/\Nn{P}{\mu})^\times$; $\emb{\tau}\colon\EP{P}{\mu}\to\CC$ is the embedding at $\tau$ that is the identity on
$\KP{P}$, and $\emb{\tau}=\emb{1}\circ\operatorname{gal}(\tau)$ for the automorphism of $\tau$. For an integer $\nu$,
$\pin{P}{\mu}{\nu}$ is the place of the class of $\nu$ if $\nu$ is prime to $\Nn{P}{\mu}$, and the unit place
otherwise. Every prime dividing $\Nn{P}{\mu}$ divides $(d+1)/4$ and splits in $\KP{P}$. Further:
$\czE{P}{\mu}=(1+\sqrt{-d})/2$ as an integer of $\EP{P}{\mu}$; $\dgen{P}{\mu}$, the value at $\rtu{P}{\mu}$ of the
derivative of the cyclotomic polynomial of order $\Nn{P}{\mu}$, which generates the different of $\EP{P}{\mu}$ over
$\KP{P}$ and divides $\Nn{P}{\mu}$; the finite idele $\alP{P}{\mu}$ of $\EP{P}{\mu}$ with component
$\dgen{P}{\mu}^{-1}$ at the primes that contain $\czE{P}{\mu}$ and $1$ elsewhere; the root discriminant
$\rdE{P}{\mu}$; and $\OmP{P}{\mu}=\exp(\mathrm{cs}(d))\cdot2\pi/\rdE{P}{\mu}$, where $\mathrm{cs}(d)$ is Euler's
constant minus $\operatorname{Re}(L'/L)(1,\chid{d})$. We write $\Ad$, $\Af$, $\Ohat$ for the adeles, the finite
adeles and the integral finite adeles of the field in question, and
\[\frl{P}{\mu}=\{x\in\EP{P}{\mu}:\ \alP{P}{\mu}^{-1}x\in\Ohat\},\]
a fractional ideal that contains the integers of $\EP{P}{\mu}$. Let $I$ be a finite set, $w_i\in\EP{P}{\mu}$ and
$m_i\in\NN$ for $i\in I$, and $M\in\NN$. The family $(w_i)$ with multiplicities $(m_i)$ is a \emph{tight frame of
length $M$} if every $w_i$ lies in $\frl{P}{\mu}$ and
\[\sum_{i\in I}m_i\,\emb{\tau}(w_i)\,\cj{\emb{\tau'}(w_i)}=\begin{cases}M&\tau=\tau',\\0&\tau\ne\tau'\end{cases}
\qquad\text{for all }\tau,\tau'\in\PlP{P}{\mu}.\]
So the vectors $(\emb{\tau}(w_i))_\tau$ of $\CC^{\PlP{P}{\mu}}$, repeated $m_i$ times, resolve $M$ times the
identity: they replace an orthonormal basis, which the lattice need not have.
\leanline{in the background library: \lnm{IsTightFrame}, \lnm{finiteImage}\sep the objects of the field are in the
concordance of the note on the CM field}

\subhead{Background 2: embeddings into $\Cp$.} Let $(p,j)$ be a $p$-adic place and $\psi\colon\EP{P}{\mu}\to\Cp$ a ring
homomorphism. It induces a prime $\mathfrak{p}_\psi$ of $\EP{P}{\mu}$, the set of integers $x$ with $\np{\psi(x)}<1$,
and extends continuously to the completion at $\mathfrak{p}_\psi$; composed with the projection of $\Af$ to that
completion this is a ring homomorphism $\hat\psi\colon\Af\to\Cp$. For $\tau\in\PlP{P}{\mu}$ put
$j_\tau(x)=j(\emb{\tau}(\cj{x}))$, so that $j(\cj{\emb{\tau}(x)})=j_\tau(x)$. Two facts are quoted: $\np{\psi(x)}\le1$
for integers $x$, and $\np{\hat\psi(a)}\le1$ for $a\in\Ohat$.
\leanline{in the background library: \lnm{PadicPlace}, \lnm{placeOfEmb}, \lnm{adeleExtension}, \lnm{weightEmbedding},
\lnm{norm\_ringOfIntegers\_le}, \lnm{norm\_adeleExtension\_le\_one}}

\subhead{Background 3: a CM datum and its theta function.} Let $d$ be a prime $\equiv7\bmod8$ and $\OO_d$ the
ring of integers of $\Q(\sqrt{-d})$, embedded in $\CC$. A \emph{CM datum for $d$} is a triple $C$ of a lattice
$\LamC{C}\subset\CC$ with $\OO_d\LamC{C}\subset\LamC{C}$, a point $t$ with $2t\in\LamC{C}$, $t\notin\LamC{C}$, and a
number $a_1$, normalised by $\wp''(t)=2$ (written $g_2=12\wp(t)^2-4$) and $a_1^2=12\wp(t)$, with
$\frac{1+\sqrt{-d}}2t\notin\LamC{C}$ and $\frac{1-\sqrt{-d}}2t\notin\LamC{C}$; $\wp$ and $\sigma_C$ are the Weierstrass functions of
$\LamC{C}$ and $\area{C}$ is the area of a fundamental parallelogram. For $m\in\NN$ and $y\in\CC$ the \emph{twisted
translation} of a function $F$ on $\CC$ is
\[(\UC{C}{m}{y}F)(z)=\exp\bigl(-m\pi(z\cj{y}+|y|^2/2)/\area{C}\bigr)\,F(z+y),\]
and the \emph{theta function} of $C$ is $\thC{C}(z)=(\UC{C}{1}{t}\sigma_C)(z)/(\UC{C}{1}{t}\sigma_C)(0)$: an entire
function with $\thC{C}(0)=1$, a section of degree one for the Riemann form of $\LamC{C}$. Further
$\xiC{C}(z)=\wp(z+t)-\wp(t)$ and $\phC{C}(z)=\tfrac12(\wp'(z+t)-a_1\xiC{C}(z))$, which vanishes at $0$ with
derivative $1$; for a triple $x=(x_1,x_2,x_3)$ of natural numbers,
$\mon{C}{x}=\thC{C}^{x_1}\sigma_C^{2x_2}(\phC{C}\thC{C}^3)^{x_3}$ is the \emph{monomial section} of degree
$x_1+2x_2+3x_3$. An automorphism $\sigma$ of $\CC$ sends $C$ to a conjugate datum $\sigma C$. $\tay{F}$ is the
Taylor series at $0$ of an entire function $F$, and for a $p$-adic place $(p,j)$, $\tayj{F}\in\Cp[[X]]$ is $\tay{F}$
with $j$ applied to its coefficients. The ring
\[\phiring{C}=\{G(\tayj{\phC{C}}):\ G\in\op[[X]]\}\subset\Cp[[X]]\]
depends on $j$. For a natural number $M\ge1$ and $r\in\NN$, $\cpart{M}{r}{F}$ is the part of a series $F\in\Cp[[X]]$
in the degrees $\equiv r\bmod M$. The following are quoted.
\begin{enumerate}[label=\textup{(Q\arabic*)},leftmargin=3.1em,itemsep=2pt]
\item \emph{Normalisation.} For every nonzero fractional ideal $\mathfrak{c}$ of $\Q(\sqrt{-d})$ there are a CM
datum $C$ for $d$ and $u\in\CC^\times$ with $\LamC{C}=u\,\mathfrak{c}$; then $\area{C}=|u|^2\Nm(\mathfrak c)\sqrt d/2$.
(Elementary.)
\item \emph{Main theorem of complex multiplication, soft form.} Let $C$ be a CM datum for $d$ with
$\LamC{C}=\Omega_\ast\OO_d$ and $\sigma\in\Aut(\CC)$. There are $\Omega\in\CC^\times$ with
$\area{\sigma C}=|\Omega|^2\sqrt d/2$, an integer $e\ge1$ with $e\Omega\in\LamC{\sigma C}$, a finite adele $\beta$ of
$\Q(\sqrt{-d})$ and elements $a_S\in\Q(\sqrt{-d})$, one for each finite set $S\subset\Q(\sqrt{-d})$, with $\beta$ a
cluster point of the $a_S$ as $S$ grows, such that for every finite $S$ and every $x\in S$
\[\sigma\bigl(\tay{\UC{C}{1}{x\Omega_\ast}\thC{C}}\bigr)=\tay{\UC{\sigma C}{1}{x\Omega a_S}\thC{\sigma C}} .\]
Reference: G.~Shimura, \emph{Introduction to the Arithmetic Theory of Automorphic Functions} (1971), Thm.~5.4;
J.~Milne, \emph{Complex Multiplication} (2020), 9.17. Relation: a \emph{weaker form} (the catalogue calls it the soft
form: the adele is not identified with the Artin map; places confirmed at second hand).
\item \emph{Chowla--Selberg formula, logarithmic form.} For a CM datum $C$ for $d$ there is $u\in\CC$ with
$\np{j(u)}=1$ at every $p$-adic place $(p,j)$ and
$\area{\sigma C}\cdot|\sigma(u)|^2=\pi\exp(\mathrm{cs}(d))\cdot2\pi/\sqrt d$ for every $\sigma\in\Aut(\CC)$. The
artifact cites Selberg and Chowla, \emph{On Epstein's zeta-function}; the catalogue has no entry for this
declaration: \emph{reference to be confirmed}.
\item \emph{Integrality of $a_1$.} $\np{j(a_1)}\le1$ at every $p$-adic place. Catalogue: a \emph{weaker form} of the
good reduction everywhere of the normalised CM curve (Milne, \emph{loc.\ cit.}; place not confirmed). Also quoted:
a complex number whose images under all embeddings into an algebraic closure of $\Q_2$ are bounded is algebraic.
\item \emph{The parameter $\phC{C}$.} $\tayj{\phC{C}}$ has constant term $0$ and linear coefficient $1$; the ring
$\phiring{C}$ is stable under $d/dX$ and contains $\tayj{\xiC{C}}$. If $T\in\phiring{C}$ has $T(0)=0$, $T'(0)=1$,
then $\phiring{C}=\{G(T)\}$, and the inverse series $Z$ of $T$ satisfies: $F\in\phiring{C}$ if and only if $F(Z)$
has coefficients in $\op$. For every $k$ the largest absolute value of the $k$-th coefficient of an element of
$\phiring{C}$ is attained at a power of $\tayj{\phC{C}}$; the same holds in two variables for the ring of integral
series in $\tayj{\phC{C}}(X)$, $\tayj{\phC{C}}(Y)$, and this ring contains $\tayj{\phC{C}}(X+Y)$ (the addition law).
\item \emph{Logarithm of theta.} There are $L\in\Cp[[X]]$ and $L_2\in\Cp$ with $\tayj{\thC{C}}=\exp(-L)$,
$L(0)=L'(0)=0$, $L''=L_2+\tayj{\xiC{C}}$, $\np{j(\sqrt{-d})\,L_2}\le1$, and for every pair $(a,N)$ in the closure
$\mathcal{N}_d\subset\Cp^2$ of the pairs $(j(\delta),j(\delta\cj{\delta}))$, $\delta\in\OO_d$: the series
$\tayj{\phC{C}}(aX)$, $\exp(NL-L(aX))$ and $(NL-L(aX))'$ lie in $\phiring{C}$. Moreover, with $q$ the norm of a
prime of $\Q(\sqrt{-d})$ above $p$: $(\zeta,\zeta^{p+1})\in\mathcal{N}_d$ for every $\zeta\in\Cp$ with $\zeta^{q}=\zeta$,
$\np{\zeta}=1$. (Substantial; the artifact points to Bannai--Kobayashi and to Mazur--Tate for comparable
statements; reference to be confirmed.)
\item \emph{Torsion translates.} Let $\delta\in\OO_d$, $\delta\ne0$, and $y\in\CC$ with $\delta y\in\LamC{C}$. Then
$\tayj{\UC{C}{1}{y}\thC{C}}(j(\cj\delta)X)=\tayj{\thC{C}}^{\delta\cj\delta}\cdot u$ for some $u\in\phiring{C}$. Also:
if $x$ is a triple of degree $m$ and $My\in\LamC{C}$ for a natural number $M\ne0$, then
$\np{j\bigl((\UC{C}{m}{y}\mon{C}{x})(0)\bigr)}\le1$. (Substantial; cf.\ Bannai--Kobayashi, Duke Math.~J.~153 (2010),
2.14--2.16, not compared; Tate's telescoping for local heights.)
\item \emph{Height of the formal group.} (i) If $\np{(F')_{p-1}}<1$ for every $F\in\phiring{C}$, then
$\np{(\tayj{\xiC{C}})_{p-1}}\ge1$ and $\chid{d}(p)\ne1$. (ii) Let $Z$ be as in (Q5), $a\in\OO_d$ with $a\cj{a}=q$ odd,
and $g\in\op[[X]]$ with $g(0)=0$ and $Z(g)=j(a)Z$; then $\np{g_k}=1$ for some $k\le q$. References: Deuring's
criterion; J.~Silverman, \emph{The Arithmetic of Elliptic Curves}, IV.7 (place confirmed for (ii)); relation:
\emph{weaker forms} (one direction of the criterion; an upper bound for the height).
\end{enumerate}
\leanline{in the background library: \lnm{CMLattice}, \lnm{CMLattice.twistedTranslate}, \lnm{CMLattice.θ},
\lnm{CMLattice.phiSeriesRing}, \lnm{congruencePart}\sep (Q1) \lnm{QuadField.exists\_latticeModule\_eq\_scaledIdeal},
\lnm{QuadField.area\_eq\_of\_eq\_scaledIdeal}\sep (Q2) \lnm{CMLattice.complex\_multiplication\_theta}\sep (Q3)
\lnm{CMLattice.chowla\_selberg}\sep (Q4) \lnm{CMLattice.norm\_a₁\_le\_one}, \lnm{isAlgebraic\_of\_padic}\sep (Q5)
\lnm{CMLattice.isSeriesGenerator\_translatedWpDeriv}, \lnm{IsSeriesGenerator.exists\_subst\_inverse},
\lnm{exists\_monomial\_max\_norm\_coeff}, \lnm{CMLattice.translatedWpDeriv\_add\_mem\_multiPhiRing}\sep (Q6)
\lnm{CMLattice.exists\_log\_theta}, \lnm{mem\_normPairClosure\_of\_pow\_eq}\sep (Q7)
\lnm{CMLattice.rescale\_translate\_eq\_theta\_pow\_mul}, \lnm{CMLattice.norm\_translate\_le\_naiveHeight\_pow}\sep (Q8)
\lnm{CMLattice.one\_le\_norm\_coeff\_translatedWp\_nonsplit}, \lnm{CMLattice.exists\_unit\_coeff\_degree\_le}}

\subhead{Background 4: the Weil representation and jets} (background library). Let $E$ be a CM field with a CM
type, whose embeddings $\emb{\tau}$ are indexed by the infinite places, and with Weil data (a splitting character
and a finite idele $\alpha$; for $\EP{P}{\mu}$ these are fixed in unit~152, with $\alpha=\alP{P}{\mu}$). A
\emph{hermitian space} $V$ over $E$ with index set $n$ and signs $\varepsilon$ is a hermitian matrix over $E$ with
an adelic point that brings it to the diagonal form $\varepsilon$ at the infinite places and to the unit form at
the finite ones; it is \emph{standard} if all signs are $1$ and the point is $1$. $\stdsp{n}$ is the standard space
with index set $n$, and $\homsp{n}{n'}$, with index set $n'\times n$, is standard too. The space $\Sw{V}$ is the
Banach space of the Weil representation: integrable functions on $\Ad^n$ with a reproducing identity. It carries
the Heisenberg operators $\heis{V}{y}$, $y\in\Ad^n$; the vacuum vector $\vac{V}$; the theta functional
$\Th{V}\colon\Sw{V}\to\CC$, invariant under the rational vectors; the Weil representation $\weil{V}$ of the adelic
unitary group, whose elements we write $(g_\infty,g)$ with $g\in\Uf{n}$; and the pairing
$\inn{V}{f}{f'}=\int\cj{f}f'$. A vector $f\in\Sw{V}$ is a function on $\Ad^n$ through
$f(y)=\Th{V}(\heis{V}{y}f)$. For a finite adelic vector $x\in\Af^{\,n}$, $\fv{x}\in\Ad^n$ is the vector with real
part $0$ and finite part $\alpha x$; $\fheis{V}{x}$ is $\heis{V}{\fv{x}}$ times a root of unity, so that
$\fheis{V}{x}\fheis{V}{x'}=\psi_0(\cdot)\fheis{V}{x+x'}$ with a value of the standard additive character $\psi_0$
of $\Af$; and
\[\Wf{V}=\text{the span over }\CC\text{ of the vectors }\fheis{V}{x}\vac{V},\ x\in\Af^{\,n} ,\]
the \emph{finite part}. $\fmult{g}$ is the multiplier of $g\in\Uf{n}$ in the Weil representation. For a set $R$ of
finite places and a function $\Psi$ on $\Af^{\,n}$, $\name{V}{R}{\Psi}$ is the operator on $\Sw{V}$ named by $\Psi$
(a combination of Heisenberg operators at the places of $R$ and their conjugates); $\Psi$ is
\emph{Schwartz--Bruhat above $R$} if it depends only on the components at the places of $R$, is locally constant,
and the parts at $R$ of its support lie in a compact set. For a modulus
$\mathfrak{m}$ (a nonzero ideal divisible by no prime together with its conjugate) $P^{\mathfrak m}$ is the mean
of the operators $\fheis{V}{x}$ over the integral $x$ supported at the primes dividing $\mathfrak{m}$ and their
conjugates. For $a\in\CC$ and $f\in\Sw{V}$ the \emph{jet series} is the power series
\[\jet{V}{a}{f}=\sum_{\beta}\frac{a^{|\beta|}}{\beta!}\,\Th{V}\bigl(v_\beta\otimes f\bigr)\,z^{\beta}\ \in\ \CC[[z_{(\tau,i)}:\tau\text{ an infinite place},\ i\in n]],\]
where $v_\beta\otimes f$ is the integral of the Heisenberg translates of $f$ against the Fock vector of the monomial
$z^\beta$. For $f\in\Sw{V}$ and $y\in\CC^{\mathrm{places}\times n}$, $\Thf{f}(y)$ is the value of $f$ at the real
vector of $y$ divided by the Gaussian; Lemma~\ref{l:line} says that $\jet{V}{a}{f}$ is the Taylor series of
$z\mapsto\Thf{f}(az)$ when the signs are $1$ and $f\in\Wf{V}$. Quoted facts are cited where they are used; the
two that carry weight are:
\begin{enumerate}[label=\textup{(Q\arabic*)},leftmargin=3.1em,itemsep=2pt,start=9]
\item \emph{Theta functions of a level are Weil translates.} Let $V$ be standard over $\EP{P}{\mu}$, $\delta$ an
integer with $\delta\cj\delta=N\ne0$, and $S$ an entire function on $\CC^{\mathrm{places}\times n}$ that has the
transformation law of $\Thf{\,\cdot\,}$ under translation by the real images of the vectors
$\cj\delta\gamma$, $\gamma\in\frl{P}{\mu}^n$. Then $S=\Thf{F}$ for some $F\in\Wf{V}$, and $F$ has a \emph{vacuum
expansion for $(\delta,N)$}: a set $A$ of $N^{2|\mathrm{places}||n|}$ indices, $x_a\in\Af^{\,n}$ with $\delta x_a$
integral, and a natural number $m_0$ with
$\sum_{a\in A}m_0\,\ell(\heis{V}{\fv{x_a}}\vac{V})\cdot\heis{V}{\fv{(-x_a)}}F=|A|\,\ell(F)\cdot\vac{V}$ for every
linear form $\ell$ on $\Sw{V}$. (Substantial: Stone--von Neumann for a finite Heisenberg group; no source cited;
reference to be confirmed.)
\item \emph{Galois action on the finite Weil representation.} Let $\sigma$ be a ring endomorphism of $\CC$,
$L\colon\Sw{V}\to\Wf{V}$ a $\sigma$-semilinear map, $B$ a finite adele and $R$ a set of finite places that contains
the conjugate of none of its members. If $\sigma(\psi_0(z))=\psi_0(B\cj{B}z)$ for all $z\in\Af$ and
$L(\heis{V}{\fv{x}}f)=\heis{V}{\fv{(Bx)}}(Lf)$ for all $x$ and all $f\in\Wf{V}$, then $B$ is a unit and there are
$c\in\CC$ and $z\in\Af$, $z\cj z=1$, with $L(\vac{V})=c\,\weil{V}(1,z\cdot1_n)\vac{V}$ and with (iii) and (v) of
Theorem~B. Reference: G.~Shimura, \emph{Theta functions with complex multiplication}, Duke Math.~J.~43 (1976),
no place given; \emph{reference to be confirmed}.
\end{enumerate}
\leanline{in the background library: \lnm{Space}, \lnm{Space.IsStd}, \lnm{𝒮}, \lnm{Space.ϱ}, \lnm{finVec},
\lnm{Space.finHeisOp}, \lnm{Space.vacElem}, \lnm{Space.θ}, \lnm{Space.ω}, \lnm{Space.dataInner},
\lnm{Space.finitePart}, \lnm{finMultiplier}, \lnm{Space.name}, \lnm{Space.projOn}, \lnm{Space.jetSeries},
\lnm{Space.entireTheta}, \lnm{Space.VacuumExpansion}\sep (Q9) \lnm{Space.exists\_entireTheta\_eq}\sep (Q10)
\lnm{Space.semilinear\_intertwines\_weil}}

\subhead{Background 5: column counts, Schur jets, gauges, kernels} (background library). A family of \emph{column
counts} on a finite set $\Phi$ is a function $\varkappa\colon\Phi\times\NN\to\NN$; $\varkappa(\tau,h)$ is the number
of columns of height $h$ at $\tau$ of a family of Young diagrams, and
$\boxes{\varkappa}{N}=\sum_{\tau\in\Phi}\sum_{h<N}h\,\varkappa(\tau,h)$ counts the boxes in columns of height below
$N$. It \emph{fits the ranks} $|n|\le|n'|$ if $\varkappa(\tau,h)=0$ for $|n|<h<|n'|$. A \emph{Weyl index} of rank
$|n|$ is a family $i=(i_\tau)_\tau$ of standard fillings: $i_\tau(h,l,b)\in n$ is the entry in the box $b<h$ of the
$l$-th column of height $h<|n|$, $l<\varkappa(\tau,h)$; the Weyl indices index the coordinates of the
representation of weight $(\varkappa,e)$ of $\prod_\tau\mathrm{GL}_n$, for $e\colon\Phi\to\Z$ the exponents of the
determinants. The \emph{Schur matrix} $S_\varkappa(z)$ of matrix variables $z_\tau=(z_{(\tau,b,c)})_{b\in n',c\in n}$
has polynomial entries indexed by a Weyl index $i'$ of rank $|n'|$ and one $i$ of rank $|n|$; its row $i'$ expresses
the product of minors of $z_\tau y_\tau$ prescribed by $i'$ in the products of minors of $y_\tau$. For a power
series $F$ in the $z_{(\tau,b,c)}$, $\schurjet{\varkappa}{F}$ is the matrix of the jets of $F$ along the entries of
$S_\varkappa$. Over $E=\EP{P}{\mu}$ with $\Phi=\PlP{P}{\mu}$: an \emph{algebraic form} of weight $(\varkappa,e)$ on
$\stdsp{n}$ is a function on $\Uf{n}$ with values in $\CC^{\text{Weyl indices}}$, equivariant under the rational
points; for a $p$-adic place $(p,j)$, $g\in\Uf{n}$ and a vector $\xi$, the \emph{gauge} $\gauge{\varkappa}{e}{g}{\xi}$ is
the largest $\np{\,\cdot\,}$ of a coordinate of $g^{-1}\cdot j(\xi)$, where $g^{-1}$ acts in weight $(\varkappa,e)$
through its images under the $\hat\jmath_\tau$; an element $g\in\Uf{n}$ is \emph{integral above $p$} if its
components at the primes above $p$ are integral with integral inverse. A \emph{kernel} $k$ from weight
$(\varkappa,e)$ on $\stdsp{n}$ to weight $(\varkappa,e_a)$ on $\stdsp{n'}$ with data in a complex vector space
$\mathcal{D}$ consists of a representation $\omega_k$ of $\Uf{n}\times\Uf{n'}$ on $\mathcal{D}$ and a linear map
$\Phi\mapsto\mathrm{jm}_k(\Phi)$ from $\mathcal{D}$ to matrices (rows: Weyl indices of rank $|n'|$; columns: of rank
$|n|$), equivariant under the rational points. Let $a$, $b$ be characters of $\mathrm{U}(1)$ over $E$ with weights
$e_a$, $e_b$ and tame parts $\psi$, $\psi'$, let $K$ be a level, and suppose that $\varkappa$ is a \emph{lift
weight}: it fits the ranks, $|n|<|n'|$, and $e_\tau=e_{b,\tau}+\varkappa(\tau,|n|)$. For $c\in\CC$ the \emph{theta
kernel} $\thker{\varkappa}{c}{a}{b}{K}$ has $\mathcal{D}=\Wf{\homsp{n}{n'}}$, $\omega_k=\pairw{a}{b}$ (the Weil
representation of the pair, twisted by $a$ and $b$) and $\mathrm{jm}_k(f)=\schurjet{\varkappa}{\jet{\homsp{n}{n'}}{c}{f}}$.
Quoted: $\schurjet{\varkappa}{\jet{\homsp{n}{n'}}{c}{f}}=c^{\boxes{\varkappa}{|n'|}}\schurjet{\varkappa}{\jet{\homsp{n}{n'}}{1}{f}}$;
the entries of $S_\varkappa$ have integral coefficients; every $g\in\Uf{n}$ is a rational point times an element
integral above $p$, and the gauge is unchanged by such elements on the right.
\leanline{in the background library: \lnm{boxes}, \lnm{FitsRanks}, \lnm{WeylIndex}, \lnm{schur}, \lnm{schurJet},
\lnm{AlgForm}, \lnm{weylGauge}, \lnm{integralAbove}, \lnm{Kernel}, \lnm{IsLiftWeight}, \lnm{thetaKernel},
\lnm{pairWeil}, \lnm{schurJet\_jetSeries\_eq\_smul}, \lnm{exists\_mul\_mem\_integralAbove}, \lnm{weylGauge\_mul}}

\subhead{Background 6: carried series} (background library). Let $\mathsf{N}$ be a function from exponents to the
non-negative real numbers and $k\in\NN$. A power series $H$ over $\Cp$ is \emph{carried by $k$ copies of
$\mathsf{N}$} if every exponent $\beta$ is a sum of $k$ exponents $\alpha_1,\dots,\alpha_k$ with
$\prod_l\mathsf{N}(\alpha_l)\ge\np{H_\beta}$, $H_\beta$ being the coefficient at $z^\beta$.
\leanline{in the background library: \lnm{HasLargeSplit}, \lnm{IsCarriedBy}}

\subhead{The definitions of the unit.}
\begin{definition}[the curve, its period, the scales]\label{d:period}
Let $P=(N_0,d)$ be parameters and $\mu$ a natural number.
\begin{enumerate}[label=(\alph*),itemsep=1pt]
\item $\szero{P}{\mu}=\bigl(\chi_{\mathrm{fin}}(\alP{P}{\mu})^{-1}\bigr)^{1/f}$, the principal value of the root,
where $f=|\PlP{P}{\mu}|$ and $\chi_{\mathrm{fin}}(a)=\upsilon(a)\,|a|^{-1/2}$ is the character of the finite ideles
of $\EP{P}{\mu}$ made from the splitting character $\upsilon$ of the Weil data and the idele norm.
\item $\cmE{P}$ is a CM datum for $d$ and $\Omz{P}\in\CC^\times$ a number with $\LamC{\cmE{P}}=\Omz{P}\,\OO_{\KP{P}}$,
both chosen by (Q1) for the unit ideal.
\item $\asc{P}=\sqrt{2\pi/\sqrt d}$, a positive real number.
\item $\sone{P}{\mu}=\szero{P}{\mu}/u_P$, where $u_P$ is the number chosen by (Q3) for the datum $\cmE{P}$.
\item $\per{P}{\mu}=\sone{P}{\mu}\cdot\asc{P}/\Omz{P}$.
\end{enumerate}
\end{definition}
\leanline{\lean{vacuumRoot}\sep\lean{cmCurve}\sep\lean{Omz}\sep\lean{areaScale}\sep\lean{periodFactor}\sep\lean{jetScale}}

\begin{definition}[the scaled vacuum and the standard data]\label{d:vac}
Let $P=(N_0,d)$ be parameters and $\mu$ a natural number.
\begin{enumerate}[label=(\alph*),itemsep=1pt]
\item $\vsc{P}{\mu}$ is the nonzero complex number $c$ that Proposition~\ref{p:gen}(i) provides (by choice) for
$\mu$ and the standard space of rank one.
\item For a hermitian space $V$ over $\EP{P}{\mu}$ with index set $n$: $\intvac{V}=\vsc{P}{\mu}^{|n|}\cdot\vac{V}$,
the \emph{scaled vacuum}.
\item For an integer $\delta$ of $\KP{P}$ and a set $S$ of finite places of $\EP{P}{\mu}$ (a predicate on them),
$\Lat{\delta}{S}$ is the group of the $x\in\Af^{\,n}$ such that $\delta x_i$ is integral at every place of $S$,
for all $i\in n$.
\item For $V$ as in (b), a $p$-adic place $(p,j)$ and an integer $\delta$ of $\KP{P}$, $\Std{V}{\delta}$ is the
$\Oj$-submodule of $\Sw{V}$ spanned by the vectors $\fheis{V}{x}\intvac{V}$ with
$x\in\Lat{\delta}{\{w:w\mid p\}}$. It depends on $j$, which is always the place of the statement.
\end{enumerate}
\end{definition}
\leanline{\lean{vacuumScale}\sep\lean{intVacuum}\sep\lean{integralAfter}\sep\lean{standardData}}

\begin{definition}[data that raise gauges]\label{d:raises}
Let $(p,j)$ be a $p$-adic place, $k$ a kernel from weight $(\varkappa,e)$ on $\stdsp{n}$ to weight $(\varkappa,e_a)$
on $\stdsp{n'}$ over $\EP{P}{\mu}$ with data in $\mathcal{D}$, $d_1\in\Uf{n}$, $d_2\in\Uf{n'}$ and $c$ a real
number. Then $\Rais{k}{d_1}{d_2}{c}$ is the set of the $\Phi\in\mathcal{D}$ such that
\[\gauge{\varkappa}{e_a}{hd_2}{\mathrm{jm}_k(\omega_k(g,h)\Phi)\,\xi}\ \le\ \max(c,0)\cdot\gauge{\varkappa}{e}{gd_1}{\xi}\]
for all $g\in\Uf{n}$, $h\in\Uf{n'}$ and all vectors $\xi$. It is closed under addition and contains $0$.
\end{definition}
\leanline{\lean{raises}}

\begin{definition}[theta products]\label{d:thprod}
Let $P=(N_0,d)$ be parameters, $\mu$ a natural number, $(p,j)$ a $p$-adic place, $\delta$ an integer of $\KP{P}$,
$r$ a real number and $c$ a finite set. Then $\ThProd{P}{\mu}{\delta}{r}{c}$ is the set of the power series
$\Theta$ over $\Cp$ in the variables $z_{(\tau,k)}$, $(\tau,k)\in\PlP{P}{\mu}\times c$, for which there are a finite
set $J$, elements $w_i\in\EP{P}{\mu}$ ($i\in J$), natural numbers $M$ and $n$, a number $b\in\CC$ and numbers
$y_{(i,k)}\in\CC$ ($(i,k)\in J\times c$) such that $(w_i)_{i\in J}$ with multiplicities $1$ is a tight frame of
length $M$;
\[\Theta=\prod_{(i,k)\in J\times c}\tayj{\UC{\cmE{P}}{1}{y_{(i,k)}}\thC{\cmE{P}}}\Bigl(\sum_{\tau}j\bigl(b\,\cj{\emb{\tau}(w_i)}\bigr)\,z_{(\tau,k)}\Bigr);\]
$p\nmid M$ and $p\nmid n$; $\np{j(b)}=r$; and $n\,\delta\,y_{(i,k)}\in\LamC{\cmE{P}}$ for all $(i,k)$.
\end{definition}
\leanline{\lean{thetaProducts}}

\begin{definition}[integrality on the null cone]\label{d:null}
Let $p$ be a prime number, $\Phi$ and $c$ finite sets, $\mathrm{pr}\colon\NN\times\Phi\to\Phi$ a map,
$S\subset\NN$ a finite set, $\varpi\in\Cp$, $\pi\in\op$, $B\in\NN$, and $\Theta$ a power series over $\Cp$ in the
variables $z_x$, $x\in\Phi\times c$. We say that
\[\nullint{\mathrm{pr}}{S}{\varpi}{\pi}{B}{\Theta}\]
if the following holds for every set $T$ of variables and every family $(Z_x)_{x\in\Phi\times c}$ of power series
$Z_x\in\op[[T]]$ without constant term: if for every commutative ring $Q$ of characteristic $p$, every ring
homomorphism $\phi\colon\op[[T]]\to Q$ with $\phi(\pi)=0$ and every $e\in S$
\[\phi\Bigl(\sum_{(\tau,k)\in\Phi\times c}Z_{(\tau,k)}^{\,p^e}\cdot Z_{(\mathrm{pr}_e(\tau),k)}\Bigr)=0 ,\]
then every coefficient of degree at most $B$ of the series $\Theta\bigl((\varpi Z_x)_x\bigr)$ has absolute value at
most $1$.
\end{definition}
\leanline{\lean{IsIntegralOnNullCone}}
The sum in the display is the \emph{coupling} of exponent $e$ of the family $Z$. The condition on it says that it
lies in the ideal generated by $p$ and $\pi$ (take for $Q$ the quotient by this ideal; if the ideal is the whole
ring there is no such $\phi$): the family is \emph{null} modulo $(p,\pi)$ for the pairings of a variable, raised to
the power $p^e$, with its partner at $\mathrm{pr}_e$.

\begin{definition}[ask, forced part, gain, loss]\label{d:ask}
Let $D$ be admissible data for parameters $P=(N_0,d)$, with components $x_1$, $b$, $m$, and let $p$ be a natural
number. For a real number $y$ and $\tau\in\PlD{D}$ put
\[\ask{D}{p}{y}{\tau}=\frac{1-\chiK{P}(p)}2\sum_{\substack{e\ge1\\ e\text{ odd or }p=d}}\Bigl(\frac{y}{p^e}-\nomW{D}{\pinD{D}{-p^e}\cdot\tau}\Bigr)_{\!+},\]
where $(\,\cdot\,)_+$ is the positive part, $\nomW{D}{\,\cdot\,}$ is the nominal weight of unit~152 ($\aD{D}/x$ at
the thick places, $\aD{D}/x_1$ at the others) and the sum over $e\in\NN$ is $0$ if the family is not summable. Then
\begin{gather*}
\forced{D}{p}=\sum_{\tau\in\PlD{D}}\ask{D}{p}{\nomW{D}{\tau}}{\tau},\\
\gain{D}{p}=m\Bigl(\sum_{\tau\in\tops{D}}\ask{D}{p}{\aD{D}}{\tau}+\kD{D}\,\forced{D}{p}\Bigr),\\
\loss{D}{p}=\fD{D}\,(\log b)^2\cdot\frac{\aD{D}}{p}\cdot m^{\ast} ,
\end{gather*}
where $m^\ast=m$ if $\shortp{D}{p}$ and $m^\ast=1$ otherwise. Here we say that $\shortp{D}{p}$ if $m$ divides $p-1$
or $m$ divides $p+1$, and that $\smallns{D}{p}$ if $p<x_1$ and $\chiK{P}(p)\ne1$. The sets $\tops{D}$ (the places of class $1$ modulo $N_0$) and the numbers $\aD{D}$, $\kD{D}$,
$\fD{D}=|\PlD{D}|$ are those of unit~152.
\end{definition}
\leanline{\lean{localAsk}\sep\lean{forcedPart}\sep\lean{gain}\sep\lean{loss}\sep\lean{IsShortPrime}\sep\lean{IsSmallNonsplit}}

\begin{definition}[served primes, collisions]\label{d:served}
Let $D$ be admissible data for parameters $P=(N_0,d)$, with component $m$, and let $p$ be a natural number.
\begin{enumerate}[label=(\alph*),itemsep=1pt]
\item $\clm{D}{p}$ is the least natural number $\mathfrak r$ with $m\le p^{\mathfrak r}$ if $p>1$ (the upper
logarithm of $m$ to the base $p$, which is $0$ for $p\le1$);
$\deltaP{D}{p}=\clm{D}{p}\cdot\fD{D}\cdot\bigl(\nD{D}{m}+\clm{D}{p}\bigr)$, where $\nD{D}{l}=\kD{D}+l(\tD{D}+1)$;
and $\jc{D}{p}=\lceil m\cdot2\fD{D}\aD{D}/p\rceil$.
\item $\ozero{D}{p}$ is the order, in $(\Z/N_0)^\times$, of the class modulo $N_0$ of the place $\pinD{D}{-p}$.
\item We say that $\served{D}{p}$ if $\chiK{P}(p)=-1$, it is not the case that $\shortp{D}{p}$, $\ozero{D}{p}$ is odd,
and $\ozero{D}{p}\le\clm{D}{p}$.
\item For a finite set $\Xi\subset\PlD{D}$, a place $\tau\in\PlD{D}$ and $e\in\NN$ we say that
\[\collision{D}{p}{\Xi}{\tau}{e}\]
if $e\ge1$, $e$ is odd or $p=d$, $\tau\in\tops{D}\smallsetminus\Xi$ and
$\pinD{D}{-p^e}\cdot\tau\in\tops{D}\smallsetminus\Xi$. Then
\[\coll{D}{p}{\Xi}=\sum_{\tau\in\PlD{D}}\ \sum_{e}\frac{\aD{D}}{p^e} ,\]
the inner sum being over the $e\in\NN$ for which this holds (it is $0$ if the family is not summable).
\end{enumerate}
\end{definition}
\leanline{\lean{ceilLogM}\sep\lean{deltaP}\sep\lean{jetCeil}\sep\lean{ozero}\sep\lean{IsServed}\sep\lean{IsCollision}\sep\lean{collisionSum}}

\begin{definition}[the coupling count]\label{d:cc}
Let $\Phi$ be a set, $p$ and $n'$ natural numbers, $\varkappa\colon\Phi\times\NN\to\NN$ a function,
$\Xi\subset\Phi$ a finite set, $\tau,\tau'\in\Phi$ and $e\in\NN$. For $\sigma\in\Phi$, summing over the heights
$h<n'$, put
\[t_\sigma=\sum_{h=1,\ \sigma\notin\Xi}\varkappa(\sigma,h),\qquad r_\sigma=\sum_{h\ne1}\varkappa(\sigma,h),\qquad
B_\sigma=\sum_{h\ne1}h\,\varkappa(\sigma,h) .\]
Then
\[\cc{p}{\varkappa}{n'}{\Xi}{\tau}{\tau'}{e}=\min\Bigl(\frac{t_\tau}{p^e},\,r_{\tau'}+t_{\tau'}\Bigr)+\frac{r_\tau}{p^e}
+\min\Bigl(\frac{B_\tau}{p^e},\,B_{\tau'}\Bigr),\]
a real number.
\end{definition}
\leanline{\lean{couplingCount}}

\begin{definition}[the column counts of the data]\label{d:kappa}
Let $D$ be admissible data for parameters $P$, with component $r$; write $k=\kD{D}$, $t=\tD{D}=5r$ and
$\gap{D}=2r^8$, and let $\wt{D}{l}{\tau}$ be the weight of unit~152 at stage $l$. For $\tau\in\PlD{D}$,
$h\in\NN$ and $l\in\NN$:
\begin{gather*}
\kap{D}{\tau}{h}=\begin{cases}
\aD{D}&h=1,\ \tau\in\tops{D},\\ 0&h=1,\ \tau\notin\tops{D},\\
\wt{D}{1}{\tau}-\lfloor(k+t+1)/2\rfloor&h=k\ne1,\\
\gap{D}-(t+1)&h\ne1,\ k<h,\ (t+1)\mid h-k,\\ 0&\text{otherwise;}\end{cases}\\
\last{D}{l}{\tau}=\begin{cases}0&l=0,\\ \lfloor\nD{D}{l-1}/2\rfloor-\wt{D}{l}{\tau}\in\Z&l\ge1;\end{cases}\\
\kept{D}{h}=\begin{cases}r^8-\lfloor(t+1)/2\rfloor&k\le h,\ (t+1)\mid h-k,\\0&\text{otherwise;}\end{cases}\\
\kapY{D}{\tau}{h}=\kept{D}{h}+\bigl(\kap{D}{\tau}{h+1}-\kept{D}{h+1}\bigr).
\end{gather*}
All subtractions are truncated, except the one in $\last{D}{l}{\tau}$.
\end{definition}
\leanline{\lean{kap}\sep\lean{lastEntry}\sep\lean{keptColumns}\sep\lean{columnsY}}

\begin{definition}[Frobenius packings, the depth loss]\label{d:packing}
\textup{(a)} Let $P$ be parameters, $\mu$ a natural number, $p$ a natural number and $n$, $i$, $y$ functions from
$\PlP{P}{\mu}$ to $\NN$; write $\mathrm{pr}_e(\tau)=\pin{P}{\mu}{-p^e}\cdot\tau$. We say that $\packing{p}{n}{i}{y}$ if
\begin{enumerate}[label=(\arabic*),itemsep=1pt]
\item for every $\tau$ with $y_\tau\ne0$: $i_\tau$ is odd and $\chiK{P}(p)=-1$;
\item for all $\tau$, $\tau'$ with $y_\tau\ne0$, $y_{\tau'}\ne0$ and every odd $e<i_\tau$:
$\mathrm{pr}_e(\tau')\ne\mathrm{pr}_{i_\tau}(\tau)$;
\item for every $\tau'$:\quad $p^{i_{\tau'}}y_{\tau'}+\sum_{\tau:\ \mathrm{pr}_{i_\tau}(\tau)=\tau'}y_\tau\ \le\ n_{\tau'}$.
\end{enumerate}
\textup{(b)} Let $D$ be admissible data, with component $b$, and $n$ a natural number. With
$c=\lceil60\,b/\kD{D}\rceil$ and $l=n-c$ (truncated),
\[\dl{D}{n}=\bigl((l+1)(l+2)\cdots(l+c)\bigr)^2 .\]
\end{definition}
\leanline{\lean{IsFrobPacking}\sep\lean{depthLossInt}}

\sect{Periods and scales}{s:period}
\statement{inksoft}{\textbf{Inputs from unit 152} (the data; note on the CM field). For parameters $P=(N_0,d)$ and
$\mu\in\NN$, with $f=|\PlP{P}{\mu}|$ and $\varpi_E$ the generator of the different of $\EP{P}{\mu}$ over $\Q$ fixed
there:

\textbf{(I1)} $\alP{P}{\mu}\,\cj{\alP{P}{\mu}}\,\dgen{P}{\mu}$ has valuation $1$ at every finite place;
$\alP{P}{\mu}^{-1}\in\Ohat$; $\Nn{P}{\mu}\alP{P}{\mu}\in\Ohat$.

\textbf{(I2)} $\prod_\tau|\emb{\tau}(\varpi_E)|=\rdE{P}{\mu}^{\,f}$ and
$|\emb{\tau}(\varpi_E)|=|\emb{\tau}(\dgen{P}{\mu})|\sqrt d$; $\prod_\tau\emb{\tau}(y)=\Nm_{\EP{P}{\mu}/\KP{P}}(y)$ for
$y\in\EP{P}{\mu}$; $\emb{\tau}(x)=x$ for $x\in\KP{P}$.

\textbf{(I3)} For every injective ring endomorphism $\phi$ of the integers of $\EP{P}{\mu}$, $\phi(\dgen{P}{\mu})$
and $\dgen{P}{\mu}$ have the same valuation at every finite place; $\dgen{P}{\mu}$ divides $\Nn{P}{\mu}$ and
generates the different of $\EP{P}{\mu}$ over $\KP{P}$.

\textbf{(I4)} A prime $p$ with $\chiK{P}(p)\ne1$ is at least $13$ and does not divide $\Nn{P}{\mu}$; $64<d$; every
prime dividing $\Nn{P}{\mu}$ divides $(d+1)/4$; $\czE{P}{\mu}(1-\czE{P}{\mu})=(d+1)/4$ and
$\cj{\czE{P}{\mu}}=1-\czE{P}{\mu}$; no power of $d$ is $\equiv-1\bmod N_0$.

\textbf{(I5)} For $\sigma\in\Aut(\CC/\KP{P})$ the permutation $\pi=\prm{\mu}{\sigma}$ of $\PlP{P}{\mu}$ satisfies
$\sigma\circ\emb{\tau}=\emb{\pi(\tau)}$.

\textbf{(I6)} \emph{Frobenius.} Let $p\nmid\Nn{P}{\mu}$ be a prime, $\phi$ a ring homomorphism from the integers
of $\EP{P}{\mu}$ to a ring of characteristic $p$, $v\in\PlP{P}{\mu}$ and $e\in\NN$. Then
$\phi(\operatorname{gal}(v)x)=\phi(\cj{x})^{p^e}$ for all integers $x$ if and only if
$\phi(\operatorname{gal}(v)\rtu{P}{\mu})=\phi(\operatorname{gal}(u)\rtu{P}{\mu})$ for the place $u=\pin{P}{\mu}{-p^e}$ and
$-\chiK{P}(p)^e\phi(\sqrt{-d})=\phi(\sqrt{-d})$. (The same holds without the conjugation, with $u=\pin{P}{\mu}{p^e}$ and
without the sign.)}
\leanline{unit 152: \lnm{Parameters.v\_alpha\_mul\_relDifferentGen}, \lnm{Parameters.inv\_α\_mem},
\lnm{Parameters.N₁\_mul\_α\_mem}\sep \lnm{Parameters.prod\_varpi\_eq\_rootDiscr\_pow},
\lnm{Parameters.infinitePlace\_apply\_varpi}, \lnm{Parameters.prod\_phi\_eq\_norm}, \lnm{Parameters.phi\_algebraMap}\sep
\lnm{Parameters.valuation\_map\_relDifferentGen}, \lnm{Parameters.relDifferent\_spec}\sep
\lnm{Parameters.thirteen\_le\_of\_chiK\_ne\_one}, \lnm{Parameters.chiK\_eq\_one\_of\_dvd}, \lnm{Parameters.lt\_d},
\lnm{Parameters.dvd\_norm\_c₀}, \lnm{Parameters.c₀InE\_mul\_one\_sub}, \lnm{Parameters.star\_c₀InE},
\lnm{Parameters.pow\_d\_ne\_neg\_one}\sep \lnm{Parameters.placePerm\_spec}\sep \lnm{Parameters.frob\_ext}}

\begin{lemma}[the period]\label{l:period}
Let $P=(N_0,d)$ be parameters and $\mu$ a natural number.
\begin{enumerate}[label=(\alph*),itemsep=1pt]
\item For every $\sigma\in\Aut(\CC/\KP{P})$: $|\sigma(\szero{P}{\mu})|^2=\rdE{P}{\mu}/\sqrt d$.
\item $\OmP{P}{\mu}>0$.
\item $\per{P}{\mu}\ne0$.
\item The lattice of $\cmE{P}$, as a module over $\OO_{\KP{P}}$, is $\Omz{P}\cdot\OO_{\KP{P}}$.
\item $\Omz{P}\in\LamC{\cmE{P}}$.
\end{enumerate}
\end{lemma}
\begin{proof}
(a) Write $\alpha=\alP{P}{\mu}$, $\partial=\dgen{P}{\mu}$, $f=|\PlP{P}{\mu}|$. Quoted from the background library:
there are $k\ge1$ and $y\in\EP{P}{\mu}^\times$ with $\chi_{\mathrm{fin}}(\alpha)^k=\prod_\tau\emb{\tau}(\cj y)$. By
(I2) this number is a norm, so it lies in $\KP{P}$ and $\sigma$ fixes it; hence
$|\sigma(\chi_{\mathrm{fin}}(\alpha))|=|\chi_{\mathrm{fin}}(\alpha)|$. Since
$\szero{P}{\mu}^f=\chi_{\mathrm{fin}}(\alpha)^{-1}$ and $|\upsilon|=1$,
\[|\sigma(\szero{P}{\mu})|^{2f}=|\chi_{\mathrm{fin}}(\alpha)|^{-2}=|\alpha| ,\]
the idele norm of $\alpha$ (with components $1$ at infinity). The idele $\alpha\cj\alpha\cdot\partial$, with
$\partial$ embedded diagonally, is a unit at every finite place by (I1), so its idele norm is
$\prod_\tau|\emb{\tau}(\partial)|^2$ (a quoted fact on ideles that are units at the finite places). By the product formula
$\partial$ has idele norm $1$, and $|\cj\alpha|=|\alpha|$; so $|\alpha|^2=\prod_\tau|\emb{\tau}(\partial)|^2$. By (I2),
$\prod_\tau|\emb{\tau}(\partial)|=\rdE{P}{\mu}^f/\sqrt d^{\,f}$. Taking $f$-th roots of non-negative numbers gives (a).

(b) All factors of $\OmP{P}{\mu}$ are positive. (c) By (a) with $\sigma=1$, $\szero{P}{\mu}\ne0$; by (Q3) with
$\sigma=1$, $u_P\ne0$; $\asc{P}\ne0$ and $\Omz{P}\ne0$. (d) is the choice of $\cmE{P}$ and $\Omz{P}$, and (e)
follows since $1\in\OO_{\KP{P}}$.
\end{proof}
\leanline{\lean{normgalvacuumRoot}\sep\lean{bigOmegapos}\sep\lean{jetScalenezero}\sep\lean{latticeModulecmCurve}\sep\lean{omegamemcmCurvelattice}\sep
quoted: \lnm{WeilData.exists\_chiFin\_pow\_eq\_prod}, \lnm{ideleAbs\_eq\_prod\_of\_v\_eq\_one}}

\begin{lemma}[the period factor at a $p$-adic place]\label{l:padic}
Let $P=(N_0,d)$ be parameters, $\mu$ a natural number and $p$ a prime number.
\begin{enumerate}[label=(\alph*),itemsep=1pt]
\item For every ring homomorphism $\psi\colon\EP{P}{\mu}\to\Cp$: $\hat\psi(\alP{P}{\mu})=\psi(\dgen{P}{\mu})^{-1}$
if $\czE{P}{\mu}\in\mathfrak{p}_\psi$, and $\hat\psi(\alP{P}{\mu})=1$ otherwise.
\item For every ring homomorphism $j\colon\CC\to\Cp$ and every $\tau\in\PlP{P}{\mu}$:
$\np{j(\sone{P}{\mu})}=\np{\hat\jmath_\tau(\alP{P}{\mu})}^{-1}$.
\item For every ring homomorphism $j\colon\CC\to\Cp$, every $y\in\frl{P}{\mu}$ and every $\tau\in\PlP{P}{\mu}$:
$\np{j\bigl(\sone{P}{\mu}\cdot\cj{\emb{\tau}(y)}\bigr)}\le1$.
\item For every ring homomorphism $j\colon\CC\to\Cp$: if $p\nmid\Nn{P}{\mu}$, then $\np{j(\sone{P}{\mu})}=1$.
\end{enumerate}
\end{lemma}
\begin{proof}
(a) is the definition of $\alP{P}{\mu}$, read at the prime $\mathfrak{p}_\psi$. (b) Let $f=|\PlP{P}{\mu}|$. By (Q3),
$\np{j(u_P)}=1$, so $\np{j(\sone{P}{\mu})}^f=\np{j(\szero{P}{\mu})^f}=\np{j(\chi_{\mathrm{fin}}(\alP{P}{\mu}))}^{-1}$. A
quoted formula gives $\np{j(\chi_{\mathrm{fin}}(a))}=\prod_{v}\np{\hat\jmath_v(a)}$ for finite ideles $a$. By (a)
the factor at $v$ is $\np{j_v(\dgen{P}{\mu})}^{-1}$ or $1$ according as $\czE{P}{\mu}\in\mathfrak{p}_{j_v}$ or not.
This factor does not depend on $v$: the condition reads $\np{j(\cj{c})}<1$ for the complex number
$c=(1+\sqrt{-d})/2$, because $\emb{v}$ is the identity on $\KP{P}$ (I2); and
$j_v(\dgen{P}{\mu})=j_1(\operatorname{gal}(v)\dgen{P}{\mu})$ has the absolute value of $j_1(\dgen{P}{\mu})$ by (I3).
So the product is the $f$-th power of the factor at $\tau$.

(c) By (b), $\np{j(\sone{P}{\mu}\cj{\emb{\tau}(y)})}=\np{\hat\jmath_\tau(\alP{P}{\mu})}^{-1}\np{j_\tau(y)}=
\np{\hat\jmath_\tau(\alP{P}{\mu}^{-1}y)}\le1$, since $\alP{P}{\mu}^{-1}y\in\Ohat$ (Background~2). (d) Use (b) at
$\tau=1$. If $\czE{P}{\mu}\notin\mathfrak{p}_{j_1}$ the factor is $1$. Otherwise it is
$\np{j_1(\dgen{P}{\mu})}^{-1}$; here $\dgen{P}{\mu}$ divides $\Nn{P}{\mu}$ (I3), which is a unit at
$\mathfrak{p}_{j_1}$ because $p\nmid\Nn{P}{\mu}$, so $\np{j_1(\dgen{P}{\mu})}=1$.
\end{proof}
\leanline{\lean{adeleExtensionalpha}\sep\lean{normiotaperiodFactor}\sep\lean{normperiodFactormulconjleone}\sep\lean{normiotaperiodFactoreqone}\sep
quoted: \lnm{WeilData.norm\_map\_chiFin}, \lnm{isEquiv\_placeOfEmb}}

\begin{lemma}[the hermitian lattice]\label{l:trace}
Let $P=(N_0,d)$ be parameters and $\mu$ a natural number.
\begin{enumerate}[label=(\alph*),itemsep=1pt]
\item For every $z\in\EP{P}{\mu}$: $\Tr_{\EP{P}{\mu}/\KP{P}}(z)=\sum_{\tau\in\PlP{P}{\mu}}\emb{\tau}(z)$ in $\CC$.
\item For all $x,y\in\frl{P}{\mu}$: $\sum_{\tau\in\PlP{P}{\mu}}\emb{\tau}(x)\,\cj{\emb{\tau}(y)}\in\OO_{\KP{P}}$.
\end{enumerate}
\end{lemma}
\begin{proof}
(a) The $\emb{\tau}$ are the composites of $\emb{1}$ with the elements of the Galois group, and $\emb1$ is the
identity on $\KP{P}$. (b) The sum is $\Tr_{\EP{P}{\mu}/\KP{P}}(x\cj y)$ by (a). The finite adele of
$x\cj y\,\dgen{P}{\mu}$ is the product of $\alP{P}{\mu}\cj{\alP{P}{\mu}}\dgen{P}{\mu}$, a unit everywhere by (I1),
with the integral adeles $\alP{P}{\mu}^{-1}x$ and $\cj{\alP{P}{\mu}^{-1}y}$. So $x\cj y\,\dgen{P}{\mu}$ is an
integer, $x\cj y$ lies in the inverse different of $\EP{P}{\mu}$ over $\KP{P}$ by (I3), and its trace is integral.
\end{proof}
\leanline{\lean{coetraceeqsum}\sep\lean{hermmemquadInt}}
In other words, $\frl{P}{\mu}$ with the form $\Tr(x\cj y)$ is a unimodular hermitian lattice over
$\OO_{\KP{P}}$: $\frl{P}{\mu}\cdot\cj{\frl{P}{\mu}}$ is the inverse different.

\sect{Tight frames}{s:frames}
\begin{lemma}[the frame of the roots of unity]\label{l:zeta}
Let $P=(N_0,d)$ be parameters, $\mu$ a natural number and $k$ a natural number. The family
$(\rtu{P}{\mu}^{\,i})_{0\le i<\Nn{P}{\mu}}$, each member with multiplicity $k$, is a tight frame of length
$k\cdot\Nn{P}{\mu}$.
\end{lemma}
\begin{proof}
$\rtu{P}{\mu}^i$ is an integer, and integers lie in $\frl{P}{\mu}$ by (I1). For places $\tau,\tau'$ put
$q=\emb{\tau}(\rtu{P}{\mu})\cj{\emb{\tau'}(\rtu{P}{\mu})}$, a root of unity with $q^{\Nn{P}{\mu}}=1$. The sum of the
frame identity is $k\sum_{i<\Nn{P}{\mu}}q^i$. If $\tau=\tau'$ then $q=1$. If $q=1$ then $\emb{\tau}$ and $\emb{\tau'}$
agree on the generator $\rtu{P}{\mu}$ of $\EP{P}{\mu}$ over $\KP{P}$, so $\tau=\tau'$. Hence the sum is
$k\,\Nn{P}{\mu}$ for $\tau=\tau'$ and $0$ otherwise.
\end{proof}
\leanline{\lean{isTightFramezetapow}}

\begin{lemma}[a factor of the cyclotomic different]\label{l:cyclo}
Let $R$ be an integral domain of characteristic $0$, $N$ a natural number, $\zeta\in R$ a primitive root of unity of
order $N$, and $l$ a prime factor of $N$. Let $\Psi_N$ be the cyclotomic polynomial of order $N$. Then there are
natural numbers $r$ and $M\le N$ with $N=l\,r$, such that a prime factor $q$ of $N$ divides $M$ if and only if
$q\ne l$, and an element $V\in R$ that divides $M$ in $R$ with
\[\Psi_N'(\zeta)\,(\zeta^{r}-1)\,M=N\,V .\]
\end{lemma}
\begin{proof}
Write $N=l\,r$ and $r=l^aM$ with $l\nmid M$; then $M\le N$ and the prime factors of $N=l^{a+1}M$ other than $l$
are those of $M$. Since $l\nmid M$, $\Psi_M(X^l)=\Psi_{Ml}(X)\Psi_M(X)$; replacing $X$ by $X^{l^a}$ and using
$\Psi_{Ml}(X^{l^a})=\Psi_N(X)$ gives $\Psi_M(X^{l^{a+1}})=\Psi_N(X)\,\Psi_M(X^{l^a})$. Differentiate at $\zeta$, a
root of $\Psi_N$:
\[l^{a+1}\zeta^{l^{a+1}-1}\,\Psi_M'(\zeta^{l^{a+1}})=\Psi_N'(\zeta)\,\Psi_M(\zeta^{l^a}) .\]
As $\Psi_M(Y)$ divides $Y^M-1$, the element $\Psi_M(\zeta^{l^a})$ divides $\zeta^{r}-1$. So $l^{a+1}$ divides
$\Psi_N'(\zeta)(\zeta^r-1)$, and $N=l^{a+1}M$ divides $\Psi_N'(\zeta)(\zeta^r-1)M$; let $V$ be the quotient. Next,
$r$ is a proper divisor of $N$, so $X^N-1=(X^r-1)\Psi_N(X)Q(X)$ for a polynomial $Q$ over $R$; differentiating at
$\zeta$ and multiplying by $\zeta$ gives $N=\Psi_N'(\zeta)(\zeta^r-1)\,Q(\zeta)\zeta$. Hence
$NM=NV\cdot Q(\zeta)\zeta$, and $M=V\,Q(\zeta)\zeta$ because $N$ is not a zero divisor: $V$ divides $M$.
\end{proof}
\leanline{\lean{existscyclotomicdifferentfactor}}

\begin{proof}[Proof of Theorem A]
Write $N=\Nn{P}{\mu}$, $\varsigma=\rtu{P}{\mu}$, $c_0=\czE{P}{\mu}$, $\partial=\dgen{P}{\mu}=\Psi_N'(\varsigma)$ and
$\nu=(d+1)/4$; so $c_0(1-c_0)=\nu$ and $\cj{c_0}=1-c_0$ (I4). The frame will be the union of the frame of
Lemma~\ref{l:zeta} with a second family, built to contain a vector that is a unit against $\sone{P}{\mu}$ at the
primes above $p$.

\emph{Step 1: auxiliary numbers.} For each prime factor $l$ of $N$ let $r_l$, $M_l$, $V_l$ be as in
Lemma~\ref{l:cyclo} for $\zeta=\varsigma$ in the integers of $\EP{P}{\mu}$, and put $a_l=\varsigma^{r_l}$, a
primitive $l$-th root of unity. Every prime factor of $N$ divides $\nu$ (I4), so $N$ divides $\nu^N$. The map
$y\mapsto\nu y^2+y$ of $\Z/N$ is injective: if $\nu y^2+y=\nu y'^2+y'$ then $(y-y')(\nu(y+y')+1)=0$, and
$\nu(y+y')+1$ is a unit because $\nu(y+y')$ is nilpotent. So the map is surjective, and there are natural
numbers $z<N$ and $e$ with $\nu z^2+z+1=Ne$. Put $t=1-c_0+\nu z$. Then
\begin{equation}\label{e:t}
t+\cj t=1+2\nu z,\qquad t\,\cj t=\nu+\nu z+\nu^2z^2=\nu Ne,\qquad (c_0+\nu z)\,t=\nu Ne .
\end{equation}

\emph{Step 2: the second family.} For a prime factor $l$ of $N$ and $0\le h<r_l$, $0\le k<l$, $0\le i<l$ put
\[u_{l,h,k,i}=\varsigma^{h}\Bigl([k=0]\,a_l^{\,i}+\lambda_l\,(a_l^{\,i}-a_l^{\,k})\Bigr),\qquad\lambda_l=\frac{t\,M_l}{N} .\]
\emph{Claim: for all places $\tau,\tau'$,
$\sum_{h,k,i}\emb{\tau}(u_{l,h,k,i})\cj{\emb{\tau'}(u_{l,h,k,i})}=[\tau=\tau']\cdot c_l$ with
$c_l=N+M_l\bigl(1+2\nu(z+l\,M_l\,e)\bigr)$.} Let $A=\emb{\tau}(a_l)$, $A'=\emb{\tau'}(a_l)$ (primitive $l$-th roots of unity),
$G=\sum_{i<l}(A\cj{A'})^i$, and let $\lambda$, $\lambda'$ be the images of $\lambda_l$ under $\emb{\tau}$, $\emb{\tau'}$.
Using $\sum_{i<l}A^i=\sum_{i<l}\cj{A'}^{\,i}=0$, the sum over $k$ and $i$ of the products of the two brackets is
\[G+\cj{\lambda'}G+\lambda G+2l\,\lambda\cj{\lambda'}G=G\,(1+\lambda+\cj{\lambda'}+2l\lambda\cj{\lambda'}) :\]
the four terms come from $k=0$ with the two first summands, from $k=0$ with one first and one second summand
(twice), and from the two second summands. The sum over $h$ contributes
$\sum_{h<r_l}(\emb{\tau}(\varsigma)\cj{\emb{\tau'}(\varsigma)})^h$, and its product with $G$ is
$\sum_{h<N}(\emb{\tau}(\varsigma)\cj{\emb{\tau'}(\varsigma)})^h=N[\tau=\tau']$ by Lemma~\ref{l:zeta}. For $\tau=\tau'$,
\eqref{e:t} gives $\lambda+\cj\lambda=M_l(1+2\nu z)/N$ and $\lambda\cj\lambda=M_l^2\nu e/N$, whence the claim.
(The formal proof establishes this identity by a Gröbner basis computation.)

\emph{Step 3: lengths.} Let $c=\sum_lc_l$. It is prime to $N$: modulo a prime factor $q$ of $N$ the summands with
$l\ne q$ vanish, because $q$ divides $N$ and $M_l$, and $c_q\equiv M_q\not\equiv0$, because $q$ divides $N$ and
$\nu$ but not $M_q$. A crude bound, from $M_l,l,z\le N$, $e\le2\nu N$ and at most $N$ prime factors, is
$c\le8\nu^2N^5$; with $9\nu^2\le d^6$ and the hypothesis,
\[c_\ast(c+N)+(N-1)(c-1)\ \le\ (c_\ast+1)\,9\nu^2N^6\ \le\ (dN)^6(c_\ast+1)\ \le\ M .\]
Since $c$ and $N$ are coprime, every integer $\ge(N-1)(c-1)$ is a combination $Ac+BN$ with natural numbers $A$,
$B$. Choose them with $Ac+BN=M-c_\ast(c+N)$. Take the vectors $u_{l,h,k,i}$ with multiplicity $A+c_\ast$ and the
vectors $\varsigma^i$, $i<N$, with multiplicity $B+c_\ast$. By the claim and Lemma~\ref{l:zeta} the frame identity
holds with length $(A+c_\ast)c+(B+c_\ast)N=M$, and all multiplicities are at least $c_\ast$.

\emph{Step 4: the vectors lie in $\frl{P}{\mu}$.} For $\varsigma^i$ this is Lemma~\ref{l:zeta}. For
$u=u_{l,h,k,i}$ write $a_l^{\,i}-a_l^{\,k}=(a_l-1)g$ with an integer $g$, and let $w$ be a finite place. If
$c_0\in w$, the component of $\alP{P}{\mu}$ is $\partial^{-1}$ and we need $\partial u$ integral at $w$; indeed
$\partial\lambda_l(a_l-1)=tV_l$ by Lemma~\ref{l:cyclo}, so $\partial u$ is an integer. If $c_0\notin w$ the
component is $1$ and $c_0$ is a unit at $w$, so it suffices that $c_0^N\lambda_l$ is an integer, that is, that $N$
divides $c_0^Nt$. Now $c_0^N-(-\nu z)^N$ is divisible by $c_0+\nu z$, and $(c_0+\nu z)t=\nu Ne$ by \eqref{e:t}; so
$c_0^Nt\equiv(-z)^N\nu^Nt\equiv0$ modulo $N$, as $N$ divides $\nu^N$.

\emph{Step 5: the unit vector.} If $\np{j(\sone{P}{\mu})}=1$, the vector $\varsigma^0=1$ will do. Otherwise $p$
divides $N$ by Lemma~\ref{l:padic}(d), and we take $u=u_{p,0,1,0}=\lambda_p(1-a_p)$. Let $\tau$ be a place. By
Lemma~\ref{l:padic}(b) and (a), $c_0\in\mathfrak{p}_{j_\tau}$ (otherwise $\np{j(\sone{P}{\mu})}=1$), and
\[\np{j\bigl(\sone{P}{\mu}\,\cj{\emb{\tau}(u)}\bigr)}=\np{j_\tau(\partial u)}=\np{j_\tau(tV_p)} ,\]
since $\partial u=-tV_p$. This is at most $1$. If it were smaller, $tV_p$ would lie in $\mathfrak{p}_{j_\tau}$;
but $t\equiv1$ modulo $\mathfrak{p}_{j_\tau}$ (which contains $c_0$ and $p$, and $p$ divides $\nu$), so $V_p$ and
then $M_p$ would lie in it, and $p$ would divide $M_p$, against Lemma~\ref{l:cyclo}.
\end{proof}
\leanline{\lean{existsisTightFrame}\sep the coin lemma is Mathlib's \lnm{Nat.exists\_add\_mul\_eq\_of\_gcd\_dvd\_of\_mul\_pred\_le}}

\sect{Jets are Taylor coefficients of products of curve theta functions}{s:jets}
\begin{lemma}[jets and Taylor series]\label{l:line}
Let $E$ be a field with a CM type and Weil data, $n$ a finite set, $V$ a hermitian space over $E$ with index set
$n$ and all signs $1$, and $f\in\Wf{V}$. Write $\Phi$ for the set of infinite places of $E$.
\begin{enumerate}[label=(\alph*),itemsep=1pt]
\item For every $a\in\CC$ and $v\in\CC^{\Phi\times n}$: the Taylor series at $0$ of the function
$u\mapsto\Thf{f}(u\cdot a\cdot v)$ of one complex variable is obtained from $\jet{V}{a}{f}$ by substituting
$v_x\,u$ for the variable $z_x$, for all $x\in\Phi\times n$.
\item If $a\in\CC$, $a\ne0$ and $\jet{V}{a}{f}=0$, then $f=0$.
\item Let $I$ be a finite set, $G_i$ ($i\in I$) entire functions on $\CC$ and $\lambda_{c,i}\in\CC^{\Phi\times n}$
for $c\in\CC$, $i\in I$, such that $\Thf{f}(c\,z)=\prod_{i\in I}G_i(\lambda_{c,i}\cdot z)$ for all $c\in\CC$ and
$z\in\CC^{\Phi\times n}$, where $\lambda\cdot z=\sum_x\lambda_xz_x$. Then $\Th{V}(f)=\prod_iG_i(0)$, and for
every $c\in\CC$
\[\jet{V}{c}{f}=\prod_{i\in I}\tay{G_i}\bigl(\textstyle\sum_x\lambda_{c,i,x}z_x\bigr) .\]
\end{enumerate}
\end{lemma}
\begin{proof}
(a) Quoted from the background library: for $f\in\Wf{V}$ the function $\Thf{f}$ is an integral
$\Thf{f}(y)=\int f(w)\,\gamma(w)\exp(\langle w,y\rangle)\,dw$ over the real vectors $w$ in Fock coordinates, with a
Gaussian $\gamma$ and a pairing that is linear in $y$; and the coefficient of $z^\beta$ in $\jet{V}{a}{f}$ is the
same integral with $\exp$ replaced by the monomial of $\beta$, times $a^{|\beta|}/\beta!$. Put
$\ell(w)=a\langle w,v\rangle$. Then $\Thf{f}(uav)=\int f\gamma\,e^{u\ell}$; expanding the exponential and exchanging
sum and integral, which is allowed by dominated convergence because $|\ell(w)|$ is at most a constant times $|w|$ and
$e^{c|w|}\gamma(w)|f(w)|$ is integrable, shows that the coefficient of $u^k$ is $\int f\gamma\,\ell^k/k!$. This is
the jet of $\jet{V}{a}{f}$ along $(\sum_xv_xz_x)^k/k!$, which is the coefficient of $u^k$ after the substitution.

(b) By (a), for every $y$ the entire function $u\mapsto\Thf{f}(u\,a\,y)$ has Taylor series $0$, so it vanishes; at
$u=a^{-1}$ this gives $\Thf{f}(y)=0$. So $f$ vanishes at all real vectors, and such a vector of $\Sw{V}$ is $0$
(quoted).

(c) At $c=0$, $z=0$ the hypothesis reads $\Th{V}(f)=\Thf{f}(0)=\prod_iG_i(0)$. Fix $c$. Two power series in the
$z_x$ are equal if for every $v$ they give the same series in $u$ after the substitution $z_x\mapsto v_xu$ (compare
the homogeneous parts as polynomial functions of $v$). After this substitution the left side is the Taylor series
of $u\mapsto\Thf{f}(u\,c\,v)=\prod_iG_i(u\,\lambda_{c,i}\cdot v)$ by (a), and the right side is the product of the
Taylor series of the functions $u\mapsto G_i(u\,\lambda_{c,i}\cdot v)$; the Taylor series of a product of entire
functions is the product of their Taylor series.
\end{proof}
\leanline{\lean{taylorlineeqsubstjetSeries}\sep\lean{eqzeroofjetSerieseqzero}\sep\lean{jetSerieseqprod}\sep quoted:
\lnm{entireTheta\_eq\_integral}, \lnm{Space.theta\_integrated}, \lnm{jet\_subst\_linearSeries}, \lnm{Space.eq\_zero\_of\_realVec}}

\begin{proposition}[a vector whose theta function is a product over a tight frame]\label{p:F}
Let $P=(N_0,d)$ be parameters, $\mu$ a natural number, $n$ a finite set and $V$ a standard hermitian space over
$\EP{P}{\mu}$ with index set $n$. Let $M$ and $N$ be natural numbers with $N\ne0$, $I$ a finite set, $w_i\in\EP{P}{\mu}$
and $m_i\in\NN$ ($i\in I$) such that $(w_i)$ with multiplicities $(m_i)$ is a tight frame of length $M$, $\delta$ an
integer of $\KP{P}$ with $\delta\cj\delta=N$, and $x_{(i,k)}$, for $(i,k)\in I\times n$, triples of natural numbers
of degree $m_i$. Let $C$ be a CM datum for $d$ and $a\in\LamC{C}$ with $\area{C}=|a|^2/N\cdot M\cdot\sqrt d/2$. Then
there is $F\in\Wf{V}$ with the following properties.
\begin{enumerate}[label=(\arabic*),itemsep=1pt]
\item For every $\gamma\in\EP{P}{\mu}^n$, with
$G_{(i,k)}=\UC{C}{m_i}{y_{(i,k)}}\mon{C}{x_{(i,k)}}$ and
$y_{(i,k)}=-\dfrac{a}{\cj\delta}\sum_\tau\emb{\tau}(\gamma_k)\,\cj{\emb{\tau}(w_i)}$:
\begin{gather*}
\Th{V}\bigl(\heis{V}{\fv{(\alpha^{-1}\gamma)}}F\bigr)=\prod_{(i,k)}G_{(i,k)}(0),\\
\jet{V}{c}{\bigl(\heis{V}{\fv{(\alpha^{-1}\gamma)}}F\bigr)}=\prod_{(i,k)}\tay{G_{(i,k)}}\Bigl(\sum_\tau\frac{a}{\cj\delta\,\asc{P}}\,c\,\cj{\emb{\tau}(w_i)}\,z_{(\tau,k)}\Bigr)
\end{gather*}
for every $c\in\CC$; here $\alpha=\alP{P}{\mu}$, so that $\fv{(\alpha^{-1}\gamma)}$ has finite part $\gamma$.
\item For every $p$-adic place $(p,j)$ and every $x\in\Af^{\,n}$: $\np{j(F(\fv{x}))}\le1$.
\item $F$ has a vacuum expansion for $(\delta,N)$ in the sense of \textup{(Q9)}.
\end{enumerate}
\end{proposition}
\begin{proof}
Put $b=a/\cj\delta$ and, for $(i,k)\in I\times n$, the linear form
$A_{(i,k)}(z)=\sum_\tau(b/\asc{P})\,\cj{\emb{\tau}(w_i)}\,z_{(\tau,k)}$. The hermitian form of $C$ is
$H_C(z,z')=\pi\cj zz'/\area{C}$. By the frame identity,
\[\sum_{(i,k)}m_i\,H_C\bigl(A_{(i,k)}u,A_{(i,k)}v\bigr)=\frac{\pi\,|b|^2M}{\area{C}\,\asc{P}^2}\sum_{(\tau,k)}\cj{u_{(\tau,k)}}v_{(\tau,k)}=\sum_{(\tau,k)}\cj{u_{(\tau,k)}}v_{(\tau,k)},\]
because $|b|^2=|a|^2/N$, $\asc{P}^2=2\pi/\sqrt d$ and $\area{C}=|a|^2M\sqrt d/(2N)$: the linear forms carry the sum of
the Riemann forms of the curve to the standard form. Let $S=\prod_{(i,k)}\mon{C}{x_{(i,k)}}\circ A_{(i,k)}$, an
entire function. For $\gamma\in\frl{P}{\mu}^n$ the form $A_{(i,k)}$ takes at the real image of $\cj\delta\gamma$ the
value $a\sum_\tau\emb{\tau}(\gamma_k)\cj{\emb{\tau}(w_i)}$, which lies in $\LamC{C}$ by Lemma~\ref{l:trace}(b)
because $a\in\LamC{C}$ and $\LamC{C}$ is a module over $\OO_{\KP{P}}$. A monomial section of degree $m_i$ is
invariant under the twisted translations $\UC{C}{m_i}{\lambda}$, $\lambda\in\LamC{C}$, up to the multiplier of
$\LamC{C}$ (quoted). So $S$ has the transformation law required in (Q9), which gives $F\in\Wf{V}$ with
$\Thf{F}=S$ and (3). For (1), the theta function of $\heis{V}{\fv{(\alpha^{-1}\gamma)}}F$ is the translate of $S$ by
the real image of $\gamma$ (quoted), that is the product of the functions
$G_{(i,k)}\circ A_{(i,k)}$; Lemma~\ref{l:line}(c) with $\lambda_{c,(i,k)}=c$ times the coefficients of $A_{(i,k)}$
gives both formulas. For (2), a quoted approximation lemma gives $\gamma$ with
$\heis{V}{\fv{x}}F=\heis{V}{\fv{(\alpha^{-1}\gamma)}}F$, so $F(\fv x)=\prod G_{(i,k)}(0)$ by (1); and some positive
integer multiple of $\gamma$ lies in $\frl{P}{\mu}^n$, so that a multiple of each $y_{(i,k)}$ lies in $\LamC{C}$ as
above. Each factor then has $\np{j(\,\cdot\,)}\le1$ by (Q7).
\end{proof}
\leanline{\lean{existsfinitePartofisTightFrame}\sep quoted: \lnm{Space.exists\_entireTheta\_eq},
\lnm{entireTheta\_translate\_finiteImage}, \lnm{exists\_finiteImage\_approx}, \lnm{exists\_forall\_mul\_mem\_Ô},
\lnm{thetaShift\_of\_mem}, \lnm{CMLattice.norm\_monomialSection\_le}}

\sect{Conjugates of jets}{s:conj}
\begin{proposition}[a generating vector and its conjugates]\label{p:gen}
Let $P=(N_0,d)$ be parameters. There is a natural number $C$ such that for every natural number $\mu$ there is a
family $(B_\sigma)_{\sigma\in\Aut(\CC/\KP{P})}$ of finite adeles of $\EP{P}{\mu}$ with the following properties, for
every finite set $n$ and every standard hermitian space $V$ over $\EP{P}{\mu}$ with index set $n$.
\begin{enumerate}[label=(\roman*),itemsep=1pt]
\item There is $c\in\CC$, $c\ne0$, such that $\np{j(c\cdot\vac{V}(\fv{x}))}\le1$ for every $p$-adic place $(p,j)$ and
every $x\in\Af^{\,n}$, and $|\sigma(c\cdot\vac{V}(\fv{x}))|\le\exp(C\,\Nn{P}{\mu}^2)^{|n|}$ for every
$\sigma\in\Aut(\CC/\KP{P})$ and every $x\in\Af^{\,n}$.
\item There is $F\in\Wf{V}$ with $\Th{V}(F)=1$ such that $\Wf{V}$ is contained in the span over $\CC$ of the vectors
$\heis{V}{\fv{x}}F$, $x\in\Af^{\,n}$, and such that for every $\sigma\in\Aut(\CC/\KP{P})$ there are $t\in\CC$ and
$F'\in\Wf{V}$ with $|t|^2=\OmP{P}{\mu}^{-1}$ and, for every $x\in\Af^{\,n}$: applying $\sigma$ to the coefficients
of $\jet{V}{\per{P}{\mu}}{(\heis{V}{\fv{x}}F)}$ gives $\jet{V}{t}{(\heis{V}{\fv{(B_\sigma x)}}F')}$ with every variable
$z_{(\tau,i)}$ replaced by $z_{(\pi^{-1}\tau,i)}$, $\pi=\prm{\mu}{\sigma}$.
\end{enumerate}
The least such $C$ is not asked for; the formal text fixes one by choice, written $\Cgen{P}$ below.
\end{proposition}
\begin{proof}
\emph{The constant.} By (Q4) the invariant $a_1$ of $\cmE{P}$ is algebraic; let $g$ be its minimal polynomial. For
each of the finitely many roots of $g$ that are invariants $a_1$ of CM data for $d$, choose such a datum; the
absolute values of all twisted translates of its theta function form a compact set (quoted: sections of an
Appell--Humbert datum), so they are bounded by one number $M_0$ for all these data. Put $C=\lceil M_0\rceil+4$.
For $\sigma\in\Aut(\CC/\KP{P})$ apply (Q2) to $\cmE{P}$, whose lattice is $\Omz{P}\OO_{\KP{P}}$
(Lemma~\ref{l:period}(d)): it gives $\Omega_\sigma$, $e_\sigma$, $\beta_\sigma$ and the $a_S$. Let $B_\sigma$ be the
image of $\beta_\sigma$ in the finite adeles of $\EP{P}{\mu}$.

\emph{The vector $F$.} Let $N=\Nn{P}{\mu}$ and take the frame $(\rtu{P}{\mu}^i)_{i<N}$ with multiplicity $N$
(Lemma~\ref{l:zeta}), written with every vector repeated $N$ times and multiplicity $1$; its length is $N^2$. For
an integer $e\ge1$, Proposition~\ref{p:F} applies with $\delta=Ne$, $\delta\cj\delta=(Ne)^2$, all triples equal to
$(1,0,0)$, and any CM datum with a lattice vector of the right area. For $e=1$, the datum $\cmE{P}$ and
$a=\Omz{P}$ (the area of $\cmE{P}$ is $|\Omz{P}|^2\sqrt d/2$ by (Q1)) it gives $F$ with a vacuum expansion
$(A,(x_a),m_0)$, $|A|=N^{2f|n|}$, $f=|\PlP{P}{\mu}|$. Formula (1) at $\gamma=0$ gives $\Th{V}(F)=1$. Pairing the
vacuum expansion for $\ell=\Th{V}$ with $\vac{V}$ shows that
$\inn{V}{\vac{V}}{\heis{V}{\fv{(-x_a)}}F}\ne0$ for some $a$; the expansion for the linear form
$\ell=m_0^{-1}\inn{V}{\vac{V}}{\heis{V}{\fv{(-x_a)}}\,\cdot\,}$ then reads $c\cdot\vac{V}=\sum_{a'}\epsilon_{a'}\heis{V}{\fv{(-x_{a'})}}F$
with $c=|A|\ell(F)\ne0$ and coefficients $\epsilon_{a'}\in\{0,1\}$ (the pairings of two translates of the vacuum are
$0$ or $1$ here). So $c\cdot\vac{V}$ is a sum of at most $|A|$ translates of $F$. This gives the spanning statement of
(ii), since $\Wf{V}$ is spanned by translates of $\vac{V}$, and the $p$-adic bound of (i), by
Proposition~\ref{p:F}(2) and the ultrametric inequality.

\emph{The conjugate vector.} Fix $\sigma$. Apply Proposition~\ref{p:F} with $e=e_\sigma$ to the datum
$\sigma\cmE{P}$ and $a=e_\sigma\Omega_\sigma$, whose area condition is that of (Q2); this gives $F'$. Put
$t=\sigma(\sone{P}{\mu})\cdot\asc{P}/\Omega_\sigma$. Then $|t|^2=\OmP{P}{\mu}^{-1}$: by Lemma~\ref{l:period}(a),
(Q3) and $\area{\sigma\cmE{P}}=|\Omega_\sigma|^2\sqrt d/2$,
\[|t|^2=\frac{\rdE{P}{\mu}}{\sqrt d}\cdot\frac{\area{\sigma\cmE{P}}}{\pi\,e^{\mathrm{cs}(d)}\,2\pi/\sqrt d}\cdot\frac{2\pi/\sqrt d}{|\Omega_\sigma|^2}
=\frac{\rdE{P}{\mu}}{2\pi\,e^{\mathrm{cs}(d)}} .\]
\emph{The jet identity.} Both sides are locally constant in $x$ (quoted: $x\mapsto\heis{V}{\fv{x}}F$ is locally
constant for $F\in\Wf{V}$), and the vectors $\alpha^{-1}\gamma$, $\gamma\in\EP{P}{\mu}^n$, are dense; so let
$x=\alpha^{-1}\gamma$. Because $\beta_\sigma$ is a cluster point of the $a_S$, there is a finite $S$ containing the
numbers $x_{i,k}=-\Tr_{\EP{P}{\mu}/\KP{P}}(\gamma_k\cj{\rtu{P}{\mu}^{i}})/N$ with
$\heis{V}{\fv{(B_\sigma x)}}F'=\heis{V}{\fv{(a_Sx)}}F'$. By Proposition~\ref{p:F}(1) and Lemma~\ref{l:trace}(a) the two
jet series are products over $(i,k)$ of
\begin{gather*}
\tay{\UC{\cmE{P}}{1}{x_{i,k}\Omz{P}}\thC{\cmE{P}}}\Bigl(\sum_\tau\frac{\per{P}{\mu}\,\Omz{P}}{N\,\asc{P}}\cj{\emb{\tau}(\rtu{P}{\mu}^i)}z_{(\tau,k)}\Bigr)
\quad\text{and}\\
\tay{\UC{\sigma\cmE{P}}{1}{x_{i,k}\Omega_\sigma a_S}\thC{\sigma\cmE{P}}}\Bigl(\sum_\tau\frac{t\,\Omega_\sigma}{N\,\asc{P}}\cj{\emb{\tau}(\rtu{P}{\mu}^i)}z_{(\tau,k)}\Bigr).
\end{gather*}
The one-variable series correspond under $\sigma$ by (Q2). The linear forms correspond after the renaming of the
places: $\per{P}{\mu}\Omz{P}/\asc{P}=\sone{P}{\mu}$ and $t\Omega_\sigma/\asc{P}=\sigma(\sone{P}{\mu})$, and
$\sigma(\cj{\emb{\tau}(w)})=\cj{\emb{\pi(\tau)}(w)}$ by (I5).

\emph{The conjugate bound of \textup{(i)}.} $\sigma(c\vac{V}(\fv x))$ is a sum of at most $N^{2f|n|}$ numbers
$\sigma\Th{V}(\heis{V}{\fv{x'}}F)$. Each is the constant term of a jet series of a translate of $F$, hence by the jet
identity the constant term of one of $F'$, which by Proposition~\ref{p:F}(1) is a product of $N^2|n|$ values of
twisted translates of the theta function of $\sigma\cmE{P}$. The invariant of $\sigma\cmE{P}$ is a root of $g$, so
each value is at most $M_0\le e^{M_0}$. With $f\le N$ and $\log N\le N$,
$N^{2f|n|}e^{M_0N^2|n|}\le\exp\bigl((\lceil M_0\rceil+4)N^2\bigr)^{|n|}$.
\end{proof}
\leanline{\lean{existsisJetConjugategenerator}\sep quoted: \lnm{IsAppellHumbert.isCompact\_range\_norm\_thetaShift},
\lnm{Space.isLocallyConstant\_rho\_finVec}, \lnm{denseRange\_finiteImage}}

\begin{lemma}[the scaled vacuum]\label{l:vac}
Let $P=(N_0,d)$ be parameters and $\mu$ a natural number.
\begin{enumerate}[label=(\alph*),itemsep=1pt]
\item $\vsc{P}{\mu}\ne0$.
\end{enumerate}
Let further $n$ be a finite set and $V$ a standard hermitian space over $\EP{P}{\mu}$ with index set $n$.
\begin{enumerate}[label=(\alph*),itemsep=1pt,start=2]
\item For every $x\in\Af^{\,n}$: $\intvac{V}(\fv x)=\prod_{i\in n}\vsc{P}{\mu}\cdot\vac{V_1}(\fv{(x_i)})$, where $V_1$
is the standard space of rank one.
\item For every $x\in\Af^{\,n}$ and every $\sigma\in\Aut(\CC/\KP{P})$:
$|\sigma(\intvac{V}(\fv x))|\le\exp\bigl(\Cgen{P}\,\Nn{P}{\mu}^2\bigr)^{|n|}$.
\item For every $x\in\Af^{\,n}$ and every $p$-adic place $(p,j)$: $\np{j(\intvac{V}(\fv x))}\le1$.
\item Let $f\in\Sw{V}$ and let $M$ be a real number with $|f(\fv x)|\le M$ for every $x\in\Af^{\,n}$. Then
$|\inn{V}{f}{f}|\le M^2$.
\item There is $x\in\Af^{\,n}$ with $\intvac{V}(\fv x)\ne0$.
\end{enumerate}
\end{lemma}
\begin{proof}
(a) is part of the choice. (b) For a standard space the vacuum is the product of the rank-one vacua of the
coordinates (quoted). (c) and (d) follow from (b) and Proposition~\ref{p:gen}(i) for rank one, factor by factor.
(e) Quoted: $\inn{V}{f}{f}$ is the mean of $|f|^2$ over a compact quotient of $\Ad^n$. The set of points where
$|f|\le M$ is closed, and it contains a dense set: by the translation formula for the values of $f$ and the
density of $\EP{P}{\mu}$ in its archimedean completion, the value of $f$ at a point of a dense set is a value
$f(\fv{x})$ ($V$ is standard, so its adelic point is $1$). (f) Otherwise (e) with $M=0$ would give
$\inn{V}{\vac{V}}{\vac{V}}=0$.
\end{proof}
\leanline{\lean{kappanezero}\sep\lean{intVacuumfinVeceqprod}\sep\lean{normintVacuumle}\sep\lean{normintVacuumleone}\sep\lean{normdataInnerle}\sep\lean{existsintVacuumnezero}\sep
quoted: \lnm{Space.vacElem\_apply\_eq\_prod\_lineTheta}, \lnm{Space.dataInner\_eq\_quotMean}, \lnm{norm\_quotMean\_le},
\lnm{Space.apply\_vecAtPoint\_add}}

\begin{proof}[Proof of Theorem B]
Let $F$, $F'$, $t$, $B=B_\sigma$ be as in Proposition~\ref{p:gen}(ii) for $V$ and $\sigma$, and write
$\Theta\mapsto\Theta^{\pi}$ for the renaming $z_{(\tau,i)}\mapsto z_{(\pi^{-1}\tau,i)}$. The map
$g\mapsto(\jet{V}{t}{g})^{\pi}$ is injective on $\Wf{V}$ by Lemma~\ref{l:line}(b), as $t\ne0$. Let $L$ be a
left inverse of it, composed with the $\sigma$-semilinear map $f\mapsto\sigma(\jet{V}{\per{P}{\mu}}{f})$. By the jet
identity of Proposition~\ref{p:gen},
\begin{equation}\label{e:key}
L\bigl(\heis{V}{\fv{x}}F\bigr)=\heis{V}{\fv{(Bx)}}F'\qquad(x\in\Af^{\,n}).
\end{equation}
Since $\Wf{V}$ is spanned by the vectors $\heis{V}{\fv x}F$ and Heisenberg operators compose up to roots of unity
given by $\psi_0$, \eqref{e:key} gives (ii) for all $f\in\Wf{V}$ and the rule
$L(\heis{V}{\fv x}f)=\heis{V}{\fv{(Bx)}}(Lf)$, provided that $\sigma(\psi_0(z))=\psi_0(B\cj Bz)$ for all $z\in\Af$.
This identity is read off the rank-one case of the jet identity: choose a natural number $N'\ne0$ such that
translation by $N'x$ and by $N'Bx$ fixes the rank-one vectors $F$ and $F'$ for $x$ in a compact set (local
constancy); writing $z=N'x$ and comparing the two sides of the jet identity at $2x$ yields the same nonzero
series multiplied by $\sigma(\psi_0(z))$ and by $\psi_0(B\cj Bz)$. Now (Q10) applies: it gives $c$, $z$, (iii) and
(v), and that $B$ is a unit. Taking constant terms in (ii) for $f=\heis{V}{\fv x}\intvac{V}$ shows
$\Th{V}(L(\heis{V}{\fv{x}}\intvac{V}))=\sigma(\intvac{V}(\fv x))$, which is nonzero for some $x$ by
Lemma~\ref{l:vac}(f); as $L\intvac{V}$ is a multiple of $c\,\weil{V}(1,z1_n)\vac{V}$, this forces $c\ne0$. (i) is
Proposition~\ref{p:gen}.

\emph{The bound \textup{(iv)}.} By (iii) and the translation rule,
$L(\heis{V}{\fv x}\weil{V}(1,g)\intvac{V})$ is a number of absolute value $1$ times
$\heis{V}{\fv{(Bx)}}\weil{V}(1,g)L\intvac{V}$, and $L\intvac{V}=c'\weil{V}(1,z1_n)\vac{V}$ for a constant $c'$.
The values of $L\intvac{V}$ at the points $\fv{(Bx)}$ are the numbers $\sigma(\intvac{V}(\fv x))$, bounded by
$\exp(\Cgen{P}\Nn{P}{\mu}^2)^{|n|}$ (Lemma~\ref{l:vac}(c)); since $B$ is a unit these are all its values at finite
vectors, and Lemma~\ref{l:vac}(e) bounds $|c'|$ by the same number. It remains to bound a translated vacuum
$\weil{V}(1,g')\vac{V}$ at an arbitrary adelic point. Its value is an inner product with a translate of the vacuum,
which a quoted series expresses as a sum of Gaussian terms over the points of a lattice in
$\CC^{\PlP{P}{\mu}\times n}$ (the lattice moved by $g'$); two of its points have distance at least $1$,
because $\Nn{P}{\mu}\alP{P}{\mu}$ is integral (I1) and nonzero integral vectors have archimedean norm at least $1$. A
packing estimate for separated points (quoted) bounds the sum by $\exp$ of a quadratic polynomial in
$f|n|\le\Nn{P}{\mu}|n|$ with coefficients controlled by $\sqrt d$. The constant $\cP{P}$ of unit~152 is
$\Cgen{P}+2d+6$ plus non-negative terms, and the product of the two bounds is at most
$\exp(\cP{P}(\Nn{P}{\mu}|n|+1)^2)$. I did not replay the last estimate line by line; see §\ref{s:notes}.
\end{proof}
\leanline{\lean{existssemilinearisJetConjugate}\sep quoted: \lnm{Space.semilinear\_intertwines\_weil},
\lnm{exists\_natCast\_smul\_mem\_of\_isCompact}, \lnm{Space.hasSum\_dataInner}, \lnm{tsum\_le\_of\_separated}\sep
unit 152: \lnm{Parameters.c}}

\begin{lemma}[names are admissible combinations]\label{l:adm}
Let $E$ be a field with a CM type and Weil data, $n$ a finite set, $V$ a hermitian space over $E$ with index set
$n$, and $\mathfrak m$ a modulus of $E$ with residue ring $\OO_E/\mathfrak m$. Let $\Psi$ be a complex function on
$(\OO_E/\mathfrak m)^n$, $\mathsf{b}$ a proposition, $B$ a real number, $u,v\in\Sw{V}$, $g$ an integer of $E$
and $q$ a natural number. Suppose that $g\,q\in\mathfrak m$; that every value $\Psi(x)$ is integral away from $q$
(that is, $|j'(\Psi(x))|_\ell\le1$ for every prime $\ell\nmid q$ and every ring homomorphism
$j'\colon\CC\to\CC_\ell$); that, if $\mathsf{b}$ holds, $\sum_x|\sigma(\Psi(x))|\le B$ for every
$\sigma\in\Aut(\CC)$; and that $P^{\mathfrak m}u=v$. Let $\eta$ be a finite adele of $E$ with $\eta q\equiv0$
modulo $\mathfrak m$ (by default the image of $g$). Then the vector $\name{V}{\mathfrak m}{\tilde\Psi}\,u$, where
$\tilde\Psi$ is the function on $\Af^{\,n}$ that is $\Psi$ of the reduction at the vectors integral at the primes
of $\mathfrak m$ and $0$ elsewhere, is a finite combination $\sum_xr(x)\,\fheis{V}{x}v$ in which every $x$ has
$\cj\eta\,x_i$ integral at all places not dividing $q$, every $\Nm(g)^{|n|}r(x)$ is integral away from $q$, and, if
$\mathsf{b}$ holds, $\sum_x|\sigma(r(x))|\le B$ for every $\sigma\in\Aut(\CC)$. (In the terms of unit~141:
an admissible combination of the $\fheis{V}{x}v$ with support, denominator $\Nm(g)^{|n|}$ and size $B$ as said.)
\end{lemma}
\begin{proof}
By Fourier inversion on the finite group $(\OO_E/\mathfrak m)^n$, with a primitive additive character
$\psi_{\mathfrak m}$,
\[\Psi=\sum_a\Psi(a)\sum_b\frac{\psi_{\mathfrak m}(-a\cdot b)}{|\OO_E/\mathfrak m|^{n}}\,\psi_{\mathfrak m}(\,\cdot\,b) .\]
Quoted: the name of the character $x\mapsto\psi_{\mathfrak m}(x\cdot b)$, applied to $u$, is a root of unity times
$\fheis{V}{x_b}P^{\mathfrak m}u=\fheis{V}{x_b}v$ for a lift $x_b$ of $b$. The outer coefficients $\Psi(a)$ are
integral away from $q$ and of size $B$; the inner ones are roots of unity divided by $|\OO_E/\mathfrak m|^n$, and
$\Nm(g)^{|n|}/|\OO_E/\mathfrak m|^n$ is integral away from $q$ because $\Nm(\mathfrak m)$ divides $\Nm(gq)$. The
composition rule for admissible combinations (unit~141) gives the result; the lifts $x_b$ can be taken with
$\cj\eta x_b$ integral away from $q$ since $\eta q\in\mathfrak m\Ohat$ and $q$ is a unit there.
\end{proof}
\leanline{\lean{isAdmissibleCombname}\sep unit 141: \lnm{IsAdmissibleComb}, \lnm{IsAdmissibleComb.bind}\sep quoted:
\lnm{Modulus.name\_char\_eq\_translate\_projOn}}

\sect{Standard data, integral jets, gauges}{s:std}
\begin{lemma}[stability of the standard data]\label{l:std}
Let $P=(N_0,d)$ be parameters, $\mu$ a natural number and $n$ a finite set.
\begin{enumerate}[label=(\alph*),itemsep=1pt]
\item For every integer $\delta$ of $\KP{P}$ and every set $S$ of finite places: $\Ohat^{\,n}\subset\Lat{\delta}{S}$.
\end{enumerate}
Let further $V$ be a hermitian space over $\EP{P}{\mu}$ with index set $n$, $(p,j)$ a $p$-adic place and $\delta$
an integer of $\KP{P}$.
\begin{enumerate}[label=(\alph*),itemsep=1pt,start=2]
\item If $x\in\Lat{\delta}{\{w:w\mid p\}}$ and $f\in\Std{V}{\delta}$, then $\fheis{V}{x}f\in\Std{V}{\delta}$.
\item If $g\in\Uf{n}$ is integral above $p$ and $f\in\Std{V}{\delta}$, then $\weil{V}(1,g)f\in\Std{V}{\delta}$.
\item For a normed field $R$, a multiplicative map $\iota\colon\CC\to R$ and $z\in\Af$:
$|\iota(\psi_0(z))|=1$.
\item Let $n'$ be a finite set and $a$, $b$ characters of $\mathrm{U}(1)$ over $\EP{P}{\mu}$ that have weights
$e$, $e_a$ and tame parts $\psi$, $\psi'$. If $f\in\Std{\homsp{n}{n'}}{\delta}$, and $g\in\Uf{n}$, $h\in\Uf{n'}$ are
integral above $p$, then $\pairw{a}{b}\bigl((1,g),(1,h)\bigr)f\in\Std{\homsp{n}{n'}}{\delta}$.
\end{enumerate}
\end{lemma}
\begin{proof}
(a) $\delta$ is integral. (d) $\psi_0(z)$ is a root of unity. (b) It suffices to treat a generator
$f=\fheis{V}{y}\intvac{V}$; then $\fheis{V}{x}f$ is $\psi_0(\cdot)\fheis{V}{x+y}\intvac{V}$, the scalar lies in $\Oj$
by (d), and $x+y\in\Lat{\delta}{\{w\mid p\}}$. (c) Multiplication by $g$ preserves $\Lat{\delta}{\{w\mid p\}}$.
Quoted: $\weil{V}(1,g)\vac{V}$ is $\fmult{g}$, of absolute value $1$ at $j$, times a finite sum of vectors
$\fheis{V}{y}\vac{V}$ with $y\in\Ohat^{\,n}$; so $\weil{V}(1,g)\fheis{V}{x}\intvac{V}$, which is
$\fheis{V}{gx}\weil{V}(1,g)\intvac{V}$, lies in $\Std{V}{\delta}$ by (a) and (b). (e) The action is a value of the
characters, of absolute value $1$ at $j$ on elements integral above $p$ (quoted), times $\weil{\homsp{n}{n'}}$ of an
element integral above $p$; use (c).
\end{proof}
\leanline{\lean{ofintegral}\sep\lean{finHeisOpmemstandardData}\sep\lean{omegamemstandardData}\sep\lean{normmappsi}\sep\lean{pairWeilmemstandardData}\sep
quoted: \lnm{Space.finHeisOp\_finHeisOp}, \lnm{Space.omega\_vacElem\_sum}, \lnm{AutChar.HasWeight.norm\_padic\_det}}

\begin{proposition}[from theta products to jets of standard data]\label{p:transfer}
Let $P=(N_0,d)$ be parameters, $\mu$ a natural number, $n$ a finite set, $V$ a standard hermitian space over
$\EP{P}{\mu}$ with index set $n$, $(p,j)$ a $p$-adic place, $\delta$ an integer of $\KP{P}$ and
$f\in\Std{V}{\delta}$. Let $c\in\CC$, let $L$ be a $\CC$-linear form on the power series in the variables
$z_{(\tau,i)}$, $(\tau,i)\in\PlP{P}{\mu}\times n$, and let $r\ge0$ be a real number. Suppose that
$\np{j(L(\Theta))}\le r$ for every power series $\Theta$ over $\CC$ whose image under $j$ lies in
$\ThProd{P}{\mu}{\delta}{\np{j(c\cdot\sone{P}{\mu})}}{n}$. Then
\[\np{j\bigl(L(\jet{V}{c\cdot\per{P}{\mu}}{f})\bigr)}\le r .\]
\end{proposition}
\begin{proof}
The set of $f\in\Sw{V}$ with the asserted bound is an $\Oj$-submodule, so let $f=\fheis{V}{x}\intvac{V}$ with
$x\in\Lat{\delta}{\{w\mid p\}}$. Put $M=p\,(d\,\Nn{P}{\mu})^6+1$, prime to $p$. By Theorem~A (with $c_\ast=0$) there
is a tight frame of length $M^2$; repeat its vectors so that the multiplicities are $1$. Proposition~\ref{p:F}
with $N=M^2$, $\delta_0=M$, triples $(1,0,0)$, the datum $\cmE{P}$ and $a=\Omz{P}$ gives $F$ with $\Th{V}(F)=1$ and
a vacuum expansion $(A,(y_a),m_0)$ with $My_a$ integral. Applied to the linear form $\vsc{P}{\mu}^{|n|}|A|^{-1}\Th{V}$
it writes $\intvac{V}$ as a combination of the vectors $\heis{V}{\fv{(-y_a)}}F$ whose coefficients are values of
$\intvac{V}$, in $\Oj$ by Lemma~\ref{l:vac}(d), times $m_0/|A|$, where $|A|$ is a power of $M$ and so a unit at $j$.
Hence it suffices to bound $L(\jet{V}{c\per{P}{\mu}}{f'})$ for $f'=\heis{V}{\fv{(x-y_a)}}F$. Approximate:
$f'=\heis{V}{\fv{(\alpha^{-1}\gamma)}}F$ with $\gamma\in\EP{P}{\mu}^n$. By Proposition~\ref{p:F}(1),
$\jet{V}{c\per{P}{\mu}}{f'}$ is a product of Taylor series of translates $\UC{\cmE{P}}{1}{y}\thC{\cmE{P}}$ at the
linear forms with coefficients $b\,\cj{\emb{\tau}(w_i)}$, where
$b=\Omz{P}\,c\,\per{P}{\mu}/(M\asc{P})=c\,\sone{P}{\mu}/M$, so $\np{j(b)}=\np{j(c\sone{P}{\mu})}$. Its image under
$j$ lies in $\ThProd{P}{\mu}{\delta}{\np{j(c\sone{P}{\mu})}}{n}$: the frame has length $M^2$, prime to $p$; and for
the lattice condition choose a positive integer $m'$ with $m'\delta M\gamma\in\frl{P}{\mu}^n$, write $m'=p^km''$
with $p\nmid m''$, and take the integer $m''M^2$, prime to $p$. Indeed $\delta M\gamma$ is integral at the places
above $p$ already ($\delta x$ is, $My_a$ is, and the approximation is by an integral vector), so
$m''\delta M\gamma\in\frl{P}{\mu}^n$, and the translation vectors, $-\Omz{P}/M$ times
$\sum_\tau\emb{\tau}(\gamma_k)\cj{\emb{\tau}(w_i)}$, are carried into $\Omz{P}\OO_{\KP{P}}=\LamC{\cmE{P}}$ by
$m''M^2\delta$ (Lemma~\ref{l:trace}(b)). The hypothesis gives the bound.
\end{proof}
\leanline{\lean{normjetSeriesle}}

\begin{lemma}[jets of products of torsion translates]\label{l:jetint}
Let $d$ be a prime number $\equiv7\bmod8$, $C$ a CM datum for $d$, $(p,j)$ a $p$-adic place, $I$ and $S$ finite
sets, $c_{i,x}\in\Cp$ for $(i,x)\in I\times S$, and $\delta\in\OO_d$ with $\delta\ne0$. Suppose that
$\np{j(\cj\delta)^{-1}c_{i,x}/d}\le1$ for all $i$, $x$, and let $y_i\in\CC$ with $\delta y_i\in\LamC{C}$ for all
$i\in I$. Then for every polynomial $Q$ in the variables $z_x$, $x\in S$, with coefficients in $\Oj$, the jet along
$j(Q)$ of the series
\[\prod_{i\in I}\tayj{\UC{C}{1}{y_i}\thC{C}}\Bigl(\sum_{x\in S}c_{i,x}z_x\Bigr)\]
lies in $\op$.
\end{lemma}
\begin{proof}
Let $L$, $L_2$ be as in (Q6) and $\ell_i=j(\cj\delta)^{-1}\sum_xc_{i,x}z_x$. Let $\mathcal{R}$ be the
$\op$-subalgebra of $\Cp[[z]]$ generated by the series $G(\ell_i)$ with $i\in I$ and
$G\in\phiring{C}\cup\{L',\exp(-L)\}$. \emph{Every element of $\mathcal{R}$ has constant term in $\op$, and
$\mathcal{R}$ is stable under the $\partial/\partial z_x$.} By the Leibniz rule it suffices to check the generators. By the chain
rule $\partial_x(G(\ell_i))=(d\,G')(\ell_i)\cdot c_{i,x}/(j(\cj\delta)d)$, and the last factor lies in $\op$. If
$G\in\phiring{C}$ then $dG'\in\phiring{C}$ by (Q5). If $G=L'$ then $dL''=dL_2+d\,\tayj{\xiC{C}}\in\phiring{C}$,
because $dL_2=-j(\sqrt{-d})\cdot j(\sqrt{-d})L_2\in\op$ by (Q6). If $G=\exp(-L)$ then $G'=-L'G$. The constant terms
are integral ($L'(0)=0$, $\exp(-L)(0)=1$, (Q5)). It follows by induction on the degree of $Q$ that jets of
elements of $\mathcal{R}$ along $j(Q)$ lie in $\op$: the jet along $z_xQ$ is the jet of $\partial_xF$ along $Q$. Finally
each factor of the product lies in $\mathcal{R}$: by (Q7),
$\tayj{\UC{C}{1}{y_i}\thC{C}}(j(\cj\delta)X)=\exp(-L)^{\delta\cj\delta}u_i$ with $u_i\in\phiring{C}$, and the factor is
this series at $X=\ell_i$.
\end{proof}
\leanline{\lean{jetthetaprodintegral}}

\begin{proof}[Proof of Theorem C]
Apply Proposition~\ref{p:transfer} with $c=\cj\delta d$, $L=\jetalong{Q'}{\,\cdot\,}$ and $r=1$. Let $\Theta$ be a
series whose image under $j$ lies in $\ThProd{P}{\mu}{\delta}{\np{j(\cj\delta d\,\sone{P}{\mu})}}{n}$, with frame
$(w_i)$, numbers $n_0$, $b$ and translation vectors $y_{(i,k)}$. Then $j(\jetalong{Q'}{\Theta})$ is the jet along
$j(Q)$ of a product as in Lemma~\ref{l:jetint}, for the datum $\cmE{P}$, the integer $n_0\delta$ (so that
$n_0\delta\,y_{(i,k)}\in\LamC{\cmE{P}}$) and the coefficients $c_{(i,k),(\tau,k')}=[k=k']\,j(b\,\cj{\emb{\tau}(w_i)})$.
The hypothesis of the lemma holds: $\np{j(b)}=\np{j(\cj\delta d\sone{P}{\mu})}$, $n_0$ is a unit at $j$, and
$\np{j(\sone{P}{\mu}\cj{\emb{\tau}(w_i)})}\le1$ by Lemma~\ref{l:padic}(c), so
$\np{j(n_0\cj\delta)^{-1}j(b\,\cj{\emb{\tau}(w_i)})/d}\le1$.
\end{proof}
\leanline{\lean{normjetjetSeriesleone}}

\begin{proposition}[one theta lift and the gauges]\label{p:raises}
Let $P=(N_0,d)$ be parameters, $\mu$ a natural number, $n$, $n'$ finite sets and $(p,j)$ a $p$-adic place.
\begin{enumerate}[label=(\alph*),itemsep=2pt]
\item Let $k$ be a kernel as in Definition~\ref{d:raises} with data in $\mathcal{D}$, $\Phi\in\mathcal{D}$,
$d_1\in\Uf{n}$, $d_2\in\Uf{n'}$, $z\in\CC$ and $c,c'$ real numbers. If $\Phi\in\Rais{k}{d_1}{d_2}{c}$ and
$\np{j(z)}\cdot c\le c'$, then $z\Phi\in\Rais{k}{d_1}{d_2}{c'}$.
\end{enumerate}
Let further $\varkappa$ be a lift weight for $n$, $n'$, $e_b$, $e$; $a$, $b$ characters of $\mathrm{U}(1)$ over
$\EP{P}{\mu}$ with weights $e_a$, $e_b$ and tame parts $\psi$, $\psi'$; $K$ a level; and $c\in\CC$.
\begin{enumerate}[label=(\alph*),itemsep=2pt,start=2]
\item Let $k=\thker{\varkappa}{c}{a}{b}{K}$, $\delta$ an integer of $\KP{P}$ and $\Psi\in\Wf{\homsp{n}{n'}}$ with
$\Psi\in\Std{\homsp{n}{n'}}{\delta}$. Let $x,y\in\Uf{n'}$ be integral above $p$, $d_0\in\Uf{n'}$, and
$u_\tau\in\Cp^{\,n'}$ ($\tau\in\PlP{P}{\mu}$) such that the image of the matrix of $d_0$ under $\hat\jmath_\tau$ is
the diagonal matrix of $u_\tau$, for every $\tau$. Let $B$ be a real number such that for every
$f\in\Std{\homsp{n}{n'}}{\delta}$ and all Weyl indices $i'$, $i$
\[\prod_\tau\Bigl(\prod_{b\in n'}\np{u_\tau(b)}\Bigr)^{e_{a,\tau}}\prod_{h,l,s}\np{u_\tau(i'_\tau(h,l,s))}\ \cdot\
\np{j\bigl(\schurjet{\varkappa}{\jet{\homsp{n}{n'}}{c}{f}}_{i',i}\bigr)}\ \le\ B .\]
Then $\omega_k(1,x\,d_0\,y)\Psi\in\Rais{k}{1}{1}{B}$.
\item Let $k=\thker{\varkappa}{c\cdot\per{P}{\mu}}{a}{b}{K}$, $\delta$ an integer of $\KP{P}$ with $\delta\ne0$, and
$\Phi\in\Wf{\homsp{n}{n'}}$ with $\Phi\in\Std{\homsp{n}{n'}}{\delta}$. Then
\[\Phi\in\Rais{k}{1}{1}{\bigl(\np{j(c)}/\np{j(\cj\delta d)}\bigr)^{\boxes{\varkappa}{|n'|}}} .\]
\item Let $k=\thker{\varkappa}{c\cdot\per{P}{\mu}}{a}{b}{K}$, let $N\ne0$, $d'$ and $M$ be natural numbers and
$f\in\Wf{\homsp{n}{n'}}$. If $d'f\in\Std{\homsp{n}{n'}}{N}$, $\np{j(c)}\le1$ and
$d'\cdot(N\cdot d)^{\boxes{\varkappa}{|n'|}}$ divides $M$, then $Mf\in\Rais{k}{1}{1}{1}$.
\end{enumerate}
\end{proposition}
\begin{proof}
(a) The gauge is homogeneous: $\gauge{\varkappa}{e_a}{g}{z\xi}=\np{j(z)}\gauge{\varkappa}{e_a}{g}{\xi}$.

(b) Let $g\in\Uf{n}$, $g'\in\Uf{n'}$ and $\xi$ be given. Write $g=\gamma\,g_0$ and $g'xd_0=\varepsilon\,h_0\,d_0$ with
rational points $\gamma$, $\varepsilon$ and $g_0$, $h_0$ integral above $p$ (Background~5). By the equivariance of
the kernel under rational points and the invariance of the gauge under elements integral above $p$, it is enough
to bound the largest coordinate of $j$ applied to
$\mathrm{jm}_k(\omega_k(g_0,h_0d_0y)\Psi)\,\xi'$, with $d_0^{-1}$ acting in weight $(\varkappa,e_a)$, against the
largest coordinate of $j(\xi')$. Now $\omega_k(g_0,h_0\,\cdot\,y)$ with integral elements preserves
$\Std{\homsp{n}{n'}}{\delta}$ by Lemma~\ref{l:std}(e), the remaining $d_0$ acts on the Schur jets through its
diagonal images $u_\tau$, which multiply the row $i'$ by the product of the $u_\tau$ over the boxes of $i'_\tau$ and
by the determinants to the power $e_{a,\tau}$, and the entries are bounded by the hypothesis.

(c) Take $x=d_0=y=1$, $u_\tau=1$ in (b). For $f\in\Std{\homsp{n}{n'}}{\delta}$ the homogeneity of the Schur jets
(Background~5) gives
\[\schurjet{\varkappa}{\jet{\homsp{n}{n'}}{c\per{P}{\mu}}{f}}=\Bigl(\frac{c}{\cj\delta d}\Bigr)^{\boxes{\varkappa}{|n'|}}\schurjet{\varkappa}{\jet{\homsp{n}{n'}}{\cj\delta d\per{P}{\mu}}{f}} ,\]
and the entries of the last matrix are jets along polynomials with integral coefficients, so Theorem~C applies to
the standard space $\homsp{n}{n'}$.

(d) Write $M=d'(Nd)^{\boxes{\varkappa}{|n'|}}s$. By (c) with $\delta=N$,
\[d'f\in\Rais{k}{1}{1}{(\np{j(c)}/\np{j(Nd)})^{\boxes{\varkappa}{|n'|}}} ,\]
and by (a), multiplying by the integer $(Nd)^{\boxes{\varkappa}{|n'|}}s$ gives the bound
$\np{j(c)}^{\boxes{\varkappa}{|n'|}}\le1$.
\end{proof}
\leanline{\lean{smulmemraises}\sep\lean{memraisesofjetbound}\sep\lean{memraisesofstandardData}\sep\lean{memraisesofdvd}\sep quoted:
\lnm{weylGauge\_embFin}, \lnm{weylGauge\_mul}, \lnm{schur\_diagonal\_eq\_prod}, \lnm{map\_schur}}

\sect{The normal form at a prime that does not split}{s:normal}
In this section $q=p^\nu$ will be the norm of a prime of $\Q(\sqrt{-d})$ above $p$: $\nu=2$ if $p$ is inert and
$\nu=1$ otherwise.

\begin{lemma}[a valuation bound]\label{l:val}
Let $p$ be a prime number and $\nu$, $Q$ natural numbers with $p^\nu=Q+1$. Let $x$, $t$, $o$, $J$, $k$ be natural
numbers with $\nu\ne0$, $o\le1$ and $x+1+QJ=pt+o+Qk$. Then there is a natural number $a$ with
$\np{x+1}=p^{-a}$ such that: if $a>1$ then $p<x+1$; either $a=0$ or $k\equiv o+J\bmod p$; and $l\,(a-\nu)\le k+t$
for every natural number $l\le p$.
\end{lemma}
\begin{proof}
Let $a=v_p(x+1)$. If $a>1$ then $x+1\ge p^2>p$. If $a\ge1$, reduce the equation modulo $p$, where $Q\equiv-1$:
$-J\equiv o-k$. If $a>\nu$ then $p^a\le x+1\le pt+1+Qk\le p^\nu(k+t)+1$, so $k+t>p^{a-\nu}-1$ and therefore
$k+t\ge p^{a-\nu}\ge p\,(a-\nu)\ge l\,(a-\nu)$.
\end{proof}
\leanline{\lean{existsnormeqinvpowbounds}}

\begin{lemma}[substitution tangent to the identity]\label{l:tangent}
Let $p$ be a prime number, $Z\in\Cp[[X]]$ and $m$ a natural number such that $X^{m+1}$ divides $Z-X$. Then for
every natural number $k\le m+1$ and every $g\in\Cp[[X]]$ the coefficient of $X^k$ in $g(Z)$ is
$g_k+g_1\cdot(Z-X)_k$.
\end{lemma}
\begin{proof}
Write $Z=X+X^{m+1}W$. Modulo $X^{m+2}$ we have $Z^i\equiv X^i$ for $i\ne1$ and $Z\equiv X+W_0X^{m+1}$; so
$g(Z)\equiv g+g_1W_0X^{m+1}$, and $(Z-X)_k$ is $W_0$ for $k=m+1$ and $0$ for $k\le m$.
\end{proof}
\leanline{\lean{coeffsubsttangent}}

\begin{lemma}[only two degrees keep a denominator]\label{l:deriv}
Let $p\ge13$ be a prime number, $\nu\in\{1,2\}$, $Q$ a natural number with $p^\nu=Q+1$, and $\varpi\in\Cp$ with
$\np{\varpi^Q}^2=p^{-\nu}$ and $\varpi\ne0$. Let $m$, $n$, $Z$, $\mu$ be power series in $\Cp[[X]]$ such that $n'$
and $Z'$ have coefficients in $\op$; $n(0)=0$; $Z_{l+1}=0$ for every $l$ not divisible by $Q$; $Z_{Q+1}=-\varpi^{-Q}$;
$m_k=0$ for $k\le2$; $m'=n\,Z'$; and
\[\mu=\bigl(\varpi^{p+1}/p\bigr)^{-1}\cdot\cpart{Q}{p+1}{m}(\varpi X) .\]
Then $\np{\mu_{p+1}}\le1$, $\np{\mu_{p+p^\nu}}\le1$, and the series
$\varpi^{-Q}\bigl(\mu-\mu_{p+1}X^{p+1}-\mu_{p+p^\nu}X^{p+p^\nu}\bigr)$ has coefficients in $\op$.
\end{lemma}
\begin{proof}
Only the degrees $k=p+1+QI$, $I\ge0$, occur in $\mu$, and $\mu_k=p\,\varpi^{QI}m_k$. We show $\np{\mu_k}\le1$ for
$I\in\{0,1\}$ (the degrees $p+1$ and $p+p^\nu$) and $\np{\mu_k}\le\np{\varpi^Q}=p^{-\nu/2}$ for $I\ge2$. From $m'=nZ'$,
$n_0=0$ and the vanishing of the coefficients of $Z$,
\[k\,m_k=\sum_{i+1+QJ=p+QI}n_{i+1}\,(QJ+1)\,Z_{QJ+1},\qquad 0\le J\le I+1 .\]
Here $\np{n_{i+1}}\le p^{c}$ with $c=v_p(i+1)$, because $n'$ is integral, and $\np{(QJ+1)Z_{QJ+1}}\le1$, because
$Z'$ is integral. With $a=v_p(k)$ the ultrametric inequality gives
\[\log_p\np{\mu_k}\ \le\ a-1-\tfrac{\nu}2I+\max_{(i,J)}c .\]
Lemma~\ref{l:val}, for $x=p+QI$ ($t=o=1$, $J=0$) and for $x=i$ ($t=1$, $o=0$), says: $a=0$ unless
$I\equiv1\bmod p$; $c=0$ unless $I\equiv J\bmod p$; and $p\,(a-\nu)\le I+1$, $p\,(c-\nu)\le I+1$.
\emph{$I=0$}: $a=0$; either $J=0$, $i+1=p$, $c=1$, or ($\nu=1$) $J=1$, $c=0$: the exponent is at most $0$.
\emph{$I=1$}: $k=p+p^\nu$ and $a=1$; the terms are $J=0$ with $i+1\equiv-1\bmod p$, $c=0$; $J=1$ with $i+1=p$,
$c=1$, where moreover $\np{(Q+1)Z_{Q+1}}=p^{-\nu/2}$; and ($\nu=1$) $J=2$ with $c=0$. The exponent is at most
$\max(-\nu/2,\,1-\nu,\,-1/2)\le0$. \emph{$2\le I\le4$}: $a=0$ since $I\not\equiv1\bmod p$; and $c\ge1$ forces
$J=I$, hence $i+1=p$ and $c=1$. The exponent is at most $-\nu I/2\le-\nu/2$. \emph{$I\ge5$}: $a$ and $c$ are at
most $\nu+(I+1)/p$, so the exponent is at most $2\nu+2(I+1)/13-1-\nu I/2\le-\nu/2$; this is where $p\ge13$ enters.
\end{proof}
\leanline{\lean{normcoeffleofderiv}}

\begin{proposition}[the formal parameter and the logarithm of theta]\label{p:formal}
Let $d$ be a prime number $\equiv7\bmod8$, $C$ a CM datum for $d$ and $(p,j)$ a $p$-adic place. Then there are a
natural number $\nu$, $\varpi\in\Cp$, $\pi\in\op$ and power series $S$, $E$ in one variable with coefficients in
$\op$ such that:
\begin{enumerate}[label=(\Roman*),itemsep=2pt]
\item $\nu\ne0$; $\np{\varpi}=p^{-\nsw{d}{p}}$; $p^{-1}\le\np{\pi}<1$; $S(0)=0$; $S-S^{p^\nu}=X+\pi E$; and
$T(\varpi X)=\varpi S$ for some $T\in\phiring{C}$.
\item If $\chid{d}(p)\ne1$ and $13\le p$, there are $u_0,u_1,u_2\in\op$, a power series $R$ with coefficients in
$\op$ and $\psi\in\Cp[[X]]$ such that $\np{\pi}=p^{(\chid{d}(p)-1)/2}$; either $\nu=1$ and $p=d$, or $\nu=2$,
$u_0=0$ and $\np{u_1}=1$; $\psi=u_0X^2+u_1S^{p+1}+u_2S^{p+p^\nu}+\pi R$; $\psi(0)=0$; and for every natural number
$n$ prime to $p$ there are $N,\kappa\in\Cp$ with $\np{N}=1$ and $\np{\kappa}=p^{(1-\chid{d}(p))/2-1/(p-1)}$ such that
for every $y\in\CC$ with $ny\in\LamC{C}$ there is a power series $V$ with coefficients in $\op$ with
\[\tayj{\UC{C}{1}{y}\thC{C}}(\varpi X)=\bigl(\exp(\kappa\psi)\cdot V\bigr)(NX) .\]
\end{enumerate}
Here $\nsw{d}{p}$ is $1/(p^2-1)$ if $\chid{d}(p)=-1$, $1/(2(p-1))$ if $p=d$, and $1/(p-1)$ if $\chid{d}(p)=1$.
\end{proposition}
\begin{proof}
\emph{Step 1: the isogeny.} Let $a=\sqrt{-d}$ if $p=d$ and $a=p$ otherwise, and let $q=p^\nu$ be the norm of a
prime above $p$. Then $\np{j(a)}^2=p^{-s}$ with $s=1$ if $p=d$ and $s=2$ otherwise, $\nsw{d}{p}=s/(2(q-1))$, and
in the two non-split cases $a\cj a=q$ is odd.

\emph{Step 2: a parameter that commutes with roots of unity.} Let $\hat\varphi=\tayj{\phC{C}}$ and let
$\zeta\in\Cp$ be a primitive root of unity of order $q-1$. By (Q6), $(\zeta^i,\zeta^{i(p+1)})\in\mathcal{N}_d$, so
$\hat\varphi(\zeta^iX)\in\phiring{C}$, and $F\mapsto F(\zeta^iX)$ maps $\phiring{C}$ to itself. Hence
\[T=\cpart{q-1}{1}{\hat\varphi}=\frac1{q-1}\sum_{i<q-1}\zeta^{-i}\hat\varphi(\zeta^iX)\in\phiring{C} ,\]
$q-1$ being a unit of $\op$. As $T(0)=0$, $T'(0)=1$, (Q5) applies: let $Z$ be the inverse series of $T$. From
$T(\zeta X)=\zeta T$ we get $Z(\zeta X)=\zeta Z$, so $Z_{l+1}=0$ unless $q-1$ divides $l$; in particular $X^q$
divides $Z-X$. Also $Z'$ is integral: $T'(Z)$ is integral with constant term $1$ and $T'(Z)Z'=1$.

\emph{Step 3: the scale.} Choose $\varpi\in\Cp$ with $\varpi^{q-1}=-Z_q^{-1}$ and put $\pi=\varpi^{q-1}$. The series
$g=T(j(a)Z)$ is integral, because $T(j(a)X)\in\phiring{C}$ by (Q6), and $Z(g)=j(a)Z$. By Lemma~\ref{l:tangent},
$g_k=j(a)^kT_k$ for $k<q$, of absolute value at most $\np{j(a)}<1$, and, using $T_q=-Z_q$ (from $T(Z)=X$),
\[g_q=j(a)^qT_q+j(a)Z_q=j(a)\bigl(1-j(a)^{q-1}\bigr)Z_q,\qquad\np{g_q}=\np{j(a)}/\np{\pi} .\]
We claim $\np{g_q}=1$. It is at most $1$. In the non-split cases (Q8)(ii) gives a coefficient of absolute value
$1$ in a degree $\le q$, which must be $q$. In the split case ($q=p$, $a=p$), if $\np{g_p}<1$ then
$\np{p\,(Z-X)_p}<1$, and then every $F\in\phiring{C}$ has $\np{(F')_{p-1}}=\np{pF_p}<1$, since by
Lemma~\ref{l:tangent} $pF_p=p(F(Z))_p-F_1\,p(Z-X)_p$ with $F(Z)$ integral; by (Q8)(i) this gives $\chid{d}(p)\ne1$, a
contradiction. So $\np{\pi}=\np{j(a)}=p^{-s/2}$, which is $p^{(\chid{d}(p)-1)/2}$ in the non-split cases, and
$\np{\varpi}=p^{-s/(2(q-1))}=p^{-\nsw{d}{p}}$.

\emph{Step 4: the series $S$.} Let $W=\varpi^{-1}Z(\varpi X)$. Its coefficients vanish except in the degrees
$1+(q-1)l$; $W_1=1$, $W_q=\pi Z_q=-1$, and for $l\ge2$ the coefficient is $\pi^{l}Z_{1+(q-1)l}$, where
$\np{Z_{1+(q-1)l}}\le p^{b}$, $b=v_p(1+(q-1)l)$, because $Z'$ is integral. Lemma~\ref{l:val} (with $x=(q-1)l$, $t=0$,
$o=1$, $J=0$, $k=l$) says that $b=0$ unless $l\equiv1\bmod p$, and that $p\,(b-\nu)\le l$. In each of the three cases
$(s,\nu)=(2,2)$, $(1,1)$, $(2,1)$ of Step~1 this gives $2b\le s(l-1)$, that is $\np{\pi^{l-1}Z_{1+(q-1)l}}\le1$. So
$W=X-X^q-\pi E'$ with $E'$ integral. Let $S$ be the inverse series of $W$; it is integral with $S(0)=0$, and
$S-S^q=X+\pi E$ with $E=E'(S)$. By definition $Z(\varpi X)=\varpi W$; substituting $S$ gives $Z(\varpi S)=\varpi X$,
and applying $T$ gives $T(\varpi X)=\varpi S$. This proves (I).

\emph{Step 5: the logarithm.} Assume now $\chid{d}(p)\ne1$ and $p\ge13$. Let $L$, $L_2$ be as in (Q6) and
$L_1=L-\tfrac12L_2X^2$, so that $L_1''=\tayj{\xiC{C}}$ and $L_1=O(X^3)$. Apply Lemma~\ref{l:deriv} with $Q=q-1$
to $m=L_1(Z)$ and $n=L_1'(Z)$: $m'=nZ'$; $n'=\tayj{\xiC{C}}(Z)\,Z'$ is integral by (Q5); $n(0)=0$; $m_k=0$ for
$k\le2$ by Lemma~\ref{l:tangent}; and $Z$ has the required coefficients by Steps 2 and 3. With
$\kappa_0=\varpi^{p+1}/p$ this gives $u_1,u_2\in\op$ and an integral $R_0$ with
\[\kappa_0^{-1}\cpart{q-1}{p+1}{L_1(Z)}(\varpi X)=u_1X^{p+1}+u_2X^{p+q}+\pi R_0 .\]
Put $\psi=\kappa_0^{-1}\cpart{q-1}{p+1}{L}(\varpi X)$. Taking the part in a residue class of degrees commutes with
substituting $T$, because $T(\zeta X)=\zeta T$; and $T(\varpi X)=\varpi S$. Hence
$\psi=u_0X^2+u_1S^{p+1}+u_2S^{p+q}+\pi R$ with $R=R_0(S)$ and with $u_0=\kappa_0^{-1}\tfrac12L_2\varpi^2$ if
$2\equiv p+1\bmod q-1$, that is if $\nu=1$, and $u_0=0$ if $\nu=2$. In the case $\nu=1$, $p=d$: $u_0=\tfrac12L_2\,p/\pi$,
and $\np{u_0}=\np{L_2}p^{-1/2}\le1$ since $\np{j(\sqrt{-d})L_2}\le1$ and $\np{j(\sqrt{-d})}=p^{-1/2}$. In the case
$\nu=2$: $p+1<q$, so $u_1=p\,(L_1)_{p+1}=(\tayj{\xiC{C}})_{p-1}/(p+1)$, and $\np{u_1}\ge1$ by (Q8)(i), whose hypothesis
holds because $(Z-X)_p=0$ (argue as in Step 3). So $\np{u_1}=1$.

\emph{Step 6: torsion translates.} Let $n$ be prime to $p$, $\delta=n(q-1)$, $N=\delta^{-1}$ (a unit of $\op$) and
$\kappa=-\delta^2\kappa_0$; then $\np{\kappa}=p\,\np{\varpi}^{p+1}=p^{(1-\chid{d}(p))/2-1/(p-1)}$ in both non-split
cases. Let $ny\in\LamC{C}$. By (Q7) for the integer $\delta$,
$\tayj{\UC{C}{1}{y}\thC{C}}(\delta X)=\exp(-\delta^2L)\cdot u$ with $u\in\phiring{C}$. By (Q6) for the pairs
$(\zeta^{-i},\zeta^{-i(p+1)})$, rotated by $\zeta^i$, the series
$\exp\bigl(\zeta^{-i(p+1)}L(\zeta^iX)-L\bigr)$ lie in $\phiring{C}$; their product over $i<q-1$ is
\begin{gather*}
\Pi=\exp\bigl((q-1)(\cpart{q-1}{p+1}{L}-L)\bigr)\in\phiring{C},\quad\text{so that}\\
\exp(-\delta^2L)\,u=\exp\bigl(-\delta^2\cpart{q-1}{p+1}{L}\bigr)\cdot\Pi^{\,n^2(q-1)}u .
\end{gather*}
Let $V=(\Pi^{n^2(q-1)}u)(\varpi X)$. It is integral: $(\Pi^{n^2(q-1)}u)(Z)$ is integral by (Q5), and
$V$ is this series at $T(\varpi X)=\varpi S$. Rescaling by $\varpi$ and by $N$,
$\tayj{\UC{C}{1}{y}\thC{C}}(\varpi X)=\bigl(\exp(\kappa\psi)V\bigr)(NX)$.
\end{proof}
\leanline{\lean{existsformalparameter}}

\begin{lemma}[two computations in characteristic $p$]\label{l:charp}
Let $p$ be a prime number, $Q$ a commutative ring of characteristic $p$ and $J$ a finite set.
\begin{enumerate}[label=(\alph*),itemsep=2pt]
\item Let $B$, $\nu$ be natural numbers and $\mathsf{b}$ a proposition such that $\nu=1$ and $\mathsf{b}$ holds, or
$\nu=2$. Let $t_i,L_i\in Q$ ($i\in J$) with $t_i-t_i^{p^\nu}=L_i$ and $t_i^{p^l}=0$ for all $l\ge B$. If
$\sum_iL_i^{p^e}L_i=0$ for every $e$ such that $e$ is odd or $\mathsf{b}$ holds, then $\sum_it_i^{p+1}=0$ and
$\sum_it_i^{p+p^\nu}=0$.
\item Let $\Phi$ be a finite set, $I$ a set, $r_{v,i}\in Q$ for $(v,i)\in\Phi\times J$, and $t_i$ ($i\in J$) power
series over $Q$ in the variables $X_{(v,c)}$, $(v,c)\in\Phi\times I$, with $t_i(0)=0$. Let $x\in\Phi$ and $c\in I$
with $t_i-t_i^{p^2}=\sum_vr_{v,i}X_{(v,c)}$ for all $i$. Then for every exponent $\beta$: the coefficient at $X^\beta$
of $\sum_it_i^{\,p}r_{x,i}$ is nonzero if and only if there are $l\in\NN$ and $y\in\Phi$ with
$\sum_ir_{y,i}^{\,p^{2l+1}}r_{x,i}\ne0$ and $\beta=p^{2l+1}\cdot1_{(y,c)}$.
\end{enumerate}
\end{lemma}
\begin{proof}
(a) Telescoping, $t_i=\sum_{l<B}L_i^{p^{\nu l}}$. So $\sum_it_i^pt_i$ and $\sum_it_i^pt_i^{p^\nu}$ are sums of terms
$\sum_iL_i^{p^{a}}L_i^{p^{b}}$ with $a-b=1+\nu(l-l')$, odd if $\nu=2$. Such a term is the $p^{\min(a,b)}$-th power of
$\sum_iL_i^{p^{|a-b|}}L_i$, which vanishes. (b) With $\ell_i=\sum_vr_{v,i}X_{(v,c)}$, telescoping gives
$t_i\equiv\sum_{l<B}\ell_i^{p^{2l}}$ up to terms of order $\ge p^{2B}$, and
$t_i^p\equiv\sum_{l}\sum_vr_{v,i}^{p^{2l+1}}X_{(v,c)}^{p^{2l+1}}$. So
$\sum_it_i^pr_{x,i}=\sum_{l,v}\bigl(\sum_ir_{v,i}^{p^{2l+1}}r_{x,i}\bigr)X_{(v,c)}^{p^{2l+1}}$, and distinct pairs
$(l,v)$ give distinct monomials.
\end{proof}
\leanline{\lean{sumpoweqzeroofnull}\sep\lean{coeffsumfrobnezeroiff}\sep quoted: \lnm{eq\_sum\_pow\_of\_sub\_pow\_eq}}

\begin{lemma}[Frobenius pairings of a tight frame]\label{l:frob}
Let $P=(N_0,d)$ be parameters, $\mu$ a natural number, $p$ a prime number with $\chiK{P}(p)\ne1$ and
$j\colon\CC\to\Cp$ a ring homomorphism. Let $J$ be a finite set, $w_i\in\EP{P}{\mu}$ ($i\in J$) and $M$ a natural
number such that $(w_i)$ with multiplicities $1$ is a tight frame of length $M$, and let $\beta\in\op$. Then there
are $z_{v,i}\in\op$, for $v\in\PlP{P}{\mu}$ and $i\in J$, with $z_{v,i}=\beta\cdot j(\cj{\emb{v}(w_i)})$, such that
for every commutative ring $Q$ of characteristic $p$, every ring homomorphism $g\colon\op\to Q$ that vanishes on
the elements of absolute value at most $p^{(\chiK{P}(p)-1)/2}$, every natural number $e$ that is odd or for which
$p=d$, and all $v,v'\in\PlP{P}{\mu}$:
\[\sum_{i\in J}g(z_{v,i})^{p^e}\,g(z_{v',i})=\begin{cases}g(\beta)^{p^e}g(\beta)\cdot M&\text{if }\pin{P}{\mu}{-p^e}\cdot v=v',\\0&\text{otherwise.}\end{cases}\]
\end{lemma}
\begin{proof}
Let $N=\Nn{P}{\mu}$; it is prime to $p$ (I4), so $\beta=N\beta'$ with $\beta'\in\op$. The elements $x_i=Nw_i$ are
integers of $\EP{P}{\mu}$ (I1). Then $z_{v,i}=\beta'j_v(x_i)$ lies in $\op$. Let
$\phi_v\colon x\mapsto g(j(\emb{v}(x)))$ on the integers of $\EP{P}{\mu}$. By (I6), $\phi_v(\operatorname{gal}(u)x)=\phi_v(\cj x)^{p^e}$
for $u=\pin{P}{\mu}{-p^e}$ and all $x$: the condition on $\sqrt{-d}$ holds because $-\chiK{P}(p)^e=1$ for $e$ odd
and $\chiK{P}(p)=-1$, and because for $p=d$ the element $j(\sqrt{-d})$ has absolute value $p^{-1/2}$ and is
killed by $g$. Hence
\[\sum_ig(z_{v,i})^{p^e}g(z_{v',i})=g(\beta')^{p^e}g(\beta')\cdot g\Bigl(j\Bigl(\sum_i\emb{u\cdot v}(x_i)\,\cj{\emb{v'}(x_i)}\Bigr)\Bigr),\]
and the inner sum is $N^2M$ if $u\cdot v=v'$ and $0$ otherwise, by the frame identity.
\end{proof}
\leanline{\lean{existsframerowsfrob}}

\begin{proof}[Proof of Theorem D]
Apply Proposition~\ref{p:formal} to $C=\cmE{P}$; since $\chiK{P}(p)\ne1$ we have $p\ge13$ (I4), so (II) holds. This
gives $\nu$, $\varpi$, $\pi$, $S$, $E$ and $u_0$, $u_1$, $u_2$, $R$, $\psi$. Let $I$, $r\le1$ and
$\Theta\in\ThProd{P}{\mu}{1}{r}{I}$ be given, with frame $(w_i)_{i\in J}$ of length $M$, numbers $n$, $b$ and
translations $y_{(i,k)}$; so $p\nmid M$, $p\nmid n$, $\np{j(b)}=r$, $n\,y_{(i,k)}\in\LamC{\cmE{P}}$. Let $N$, $\kappa$
and $V_{(i,k)}$ be as in (II) for this $n$. Lemma~\ref{l:frob} with $\beta=Nj(b)$ gives $z_{v,i}$. Put
\[\ell_{(i,k)}=\sum_vz_{v,i}\,z_{(v,k)},\qquad\Psi=\sum_{(i,k)}\psi(\ell_{(i,k)}),\qquad V=\prod_{(i,k)}V_{(i,k)}(\ell_{(i,k)}) ,\]
where $z_{(v,k)}$ are the variables. Substituting in the definition of the theta product,
$\Theta(\varpi z)=\prod_{(i,k)}(\exp(\kappa\psi)V_{(i,k)})(\ell_{(i,k)})=\exp(\kappa\Psi)\,V$, and $\Psi(0)=0$.

(a) Let $B$, a family $Z$ as in Definition~\ref{d:null} with couplings in the ideal $(p,\pi)$ for $e\in S_B$, and
an exponent $\beta_0$ of degree at most $B$ be given. Let $Q$ be the quotient of $\op[[T]]$ by the ideal of the
series whose coefficients at the exponents $\le\beta_0$ have absolute value at most $\np\pi$; it has
characteristic $p$. In $Q$ let $t_{(i,k)}$ and $L_{(i,k)}$ be the classes of $S(\ell_{(i,k)}(Z))$ and of
$\ell_{(i,k)}(Z)=\sum_vz_{v,i}Z_{(v,k)}$. Then $t-t^{p^\nu}=L$, and $t^{p^l}=0$ for $p^l>B$. By Lemma~\ref{l:frob}
(for the map $\op\to Q$, which kills the elements of absolute value $\le\np\pi$),
\[\sum_{(i,k)}L_{(i,k)}^{p^e}L_{(i,k)}=\beta^{p^e+1}M\sum_{(v,k)}Z_{(v,k)}^{p^e}Z_{(\mathrm{pr}_e(v),k)}\]
for $e$ odd or $p=d$; this is $0$ in $Q$ for $e\in S_B$ by the hypothesis on $Z$, and for $e>\log_pB$ because
then $L^{p^e}=0$. Lemma~\ref{l:charp}(a), with the proposition $p=d$, gives $\sum t^{p+1}=\sum t^{p+p^\nu}=0$.
Moreover $u_0\sum L^2=0$: either $u_0=0$, or $p=d$ and this is the coupling with $e=0\in S_B$. So the class of
$\Psi(Z)$ vanishes: its coefficients at exponents $\le\beta_0$ have absolute value at most $\np\pi$. Then the
coefficient at $\beta_0$ of $\exp(\kappa\Psi(Z))=\sum_l\kappa^l\Psi(Z)^l/l!$ is integral, since
$\np{\kappa\pi}^l=p^{-l/(p-1)}\le\np{l!}$; the same holds at every exponent $\le\beta_0$, and $V(Z)$ is integral. So
the coefficient of $\Theta(\varpi Z)$ at $\beta_0$ is integral.

(b) Now $\chiK{P}(p)=-1$ and $r=1$; so $\nu=2$, $u_0=0$, $\np{u_1}=1$, $\np\kappa=p^{1-1/(p-1)}$, and $\beta$ is a unit.
Reduce modulo $\mxp$ and let $t_{(i,k)}$ be the reduction of $S(\ell_{(i,k)})$; then
$t_{(i,k)}-t_{(i,k)}^{p^2}=\sum_v\cj{z}_{v,i}z_{(v,k)}$ with the residues $\cj z_{v,i}$, and for a variable $x=(x_1,x_2)$
\[\frac{\partial\cj\Psi}{\partial z_x}=\cj u_1\sum_{i\in J}t_{(i,x_2)}^{\,p}\,\cj z_{x_1,i} ,\]
because $(p+1)t^p\partial t=t^p\partial t$, the derivative of $t^{p^2}$ is $0$, $p+p^2\equiv0$ and $\pi\equiv0$.
By Lemma~\ref{l:charp}(b) the coefficient at $z^{\beta_0}$ is nonzero exactly for $\beta_0=p^{2l+1}1_{(y,x_2)}$ with
$\sum_i\cj z_{y,i}^{\,p^{2l+1}}\cj z_{x_1,i}\ne0$, and by Lemma~\ref{l:frob}, with $\beta$ and $M$ units, this means
$\pin{P}{\mu}{-p^{2l+1}}\cdot y=x_1$.
\end{proof}
\leanline{\lean{thetaProductsnormalform}}

\sect{Divisibility of Schur jets}{s:schur}
\begin{lemma}[frames that kill the couplings]\label{l:frames}
Let $p$ be a prime number, $\Phi$ and $n$ finite sets, $\mathrm{pr}\colon\NN\times\Phi\to\Phi$ a map, $S\subset\NN$ a
finite set, $N$ a natural number, $\varkappa\colon\Phi\times\NN\to\NN$ a function, $\Xi\subset\Phi$ a finite set,
$Y_{\tau,a,k}\in\op$ for $\tau\in\Phi$, $a\in n$, $k\in\NN$, and $\pi\in\op$. Then there are square matrices
$M_{\tau,h}$ of size $h$ over $\op$, for $\tau\in\Phi$ and $h\in\NN$, such that
\begin{enumerate}[label=(\roman*),itemsep=2pt]
\item $\displaystyle\np{\pi}^{\,\sum_{e\in S}\sum_{\tau\in\Phi}\cc{p}{\varkappa}{N}{\Xi}{\tau}{\mathrm{pr}_e(\tau)}{e}}\ \le\
\prod_{\tau\in\Phi}\prod_{h<N}\np{\det M_{\tau,h}}^{\varkappa(\tau,h)}$;
\item for every commutative ring $R$ of characteristic $p$, every ring homomorphism $g\colon\op\to R$ with
$g(\pi)=0$, every $e\in S$, $\tau\in\Phi$, heights $h$, $h'$ and rows $i<h$, $i'<h'$: unless $h=1$ and
($\tau\in\Xi$, or $h'=1$ and $\mathrm{pr}_e(\tau)\in\Xi$),
\[g\Bigl(\sum_{a\in n}\Bigl(\sum_{k<h}(M_{\tau,h})_{i,k}Y_{\tau,a,k}\Bigr)^{p^e}\cdot\sum_{k<h'}(M_{\mathrm{pr}_e(\tau),h'})_{i',k}Y_{\mathrm{pr}_e(\tau),a,k}\Bigr)=0 .\]
\end{enumerate}
\end{lemma}
\begin{proof}
Fix $p^e$-th roots $\pi^{1/p^e}\in\op$. For a coupling $(e,\tau)$ with partner $\tau'=\mathrm{pr}_e(\tau)$ and a pair
of heights $(h,h')$ the pairing in (ii) can be made divisible by $\pi$ from either end: by a condition on the rows
of $M_{\tau,h}$ that costs $\pi^{1/p^e}$ (its $p^e$-th power appears), or by one on the rows of $M_{\tau',h'}$ that
costs $\pi$. Three devices are used. (1) Multiply all of $M_{\tau,h}$ by $\pi^{1/p^e}$, or all of $M_{\tau',h'}$ by
$\pi$: this kills the couplings with all rows at the other end, and costs $h/p^e$, respectively $h'$, in the
valuation of the determinant, per column. (2) If $h'=1$, ask only that $\pi^{1/p^e}$ divide the pairings of the rows of
$M_{\tau,h}Y_\tau$ with the single row of $Y_{\tau'}$: one linear condition on the rows, which by the input below
costs $1/p^e$ per column. (3) If $h=1$, ask that $\pi$ divide the pairings of the rows of $M_{\tau',h'}Y_{\tau'}$
with the $p^e$-th power of the single row of $Y_\tau$: cost $1$ per column. For $h,h'\ne1$ use (1) at the cheaper
end, which costs $\min(B_\tau/p^e,B_{\tau'})$ in all; for $h\ne1$, $h'=1$ use (2), cost $r_\tau/p^e$; for $h=1$ and
$\tau\notin\Xi$ use (1) on the $1\times1$ matrix $M_{\tau,1}$ or (3) on all $M_{\tau',h'}$, whichever is cheaper,
cost $\min(t_\tau/p^e,r_{\tau'}+t_{\tau'})$; the couplings exempted in (ii) are not treated. The total is the
coupling count. All conditions on one matrix $M_{\tau,h}$ are imposed at once by the following input, applied to
the diagonal matrix of the required roots of $\pi$: \emph{(unit~149) if $A$ is a square matrix over a valuation
ring and $\Lambda$ a submodule containing the image of $A$, then for vectors $\Phi_1,\dots,\Phi_h$ there is a square
matrix $M$ of size $h$ with $\det M$ dividing $\det A$ and all $\sum_kM_{i,k}\Phi_k$ in $\Lambda$.} In (ii) one raises
the divisibility by $\pi^{1/p^e}$ to the power $p^e$, which in characteristic $p$ is additive. The bookkeeping that
decides which device serves which coupling is a predicate of the formal text; that its costs add up to the
coupling count is checked there by a case analysis, which I read but did not replay (see §\ref{s:notes}).
\end{proof}
\leanline{\lean{existsnullcouplingframes}\sep unit 149: \lnm{IndexDvd.of\_det\_dvd}}

\begin{proof}[Proof of Theorem E]
\emph{Step 1: from Schur entries to products of minors.} For a matrix $y_\tau$ of size $n\times N$ (rows $a\in n$,
columns $k<N$) the row $i'$ of the Schur matrix satisfies
$\sum_iS_\varkappa(z)_{i',i}\,b_i(y)=b'_{i'}(zy)$, where $b_i(y)$ are products of minors of $y$ (times powers of
determinants) and
\[b'_{i'}(zy)=\prod_{\tau}\prod_{h<N}\prod_{l<\varkappa(\tau,h)}\det\bigl((z_\tau y_\tau)_{i'_\tau(h,l,s),\,k}\bigr)_{s,k<h}\]
is the product of minors prescribed by $i'$ (quoted: this is the identity that defines the Schur matrix, and its
coefficients are integers). A quoted lemma on Gauss norms gives a point $y$ with entries in $\op$
at which the largest absolute value of the jets of $\Theta$ along the coefficients is attained; so it suffices to
bound the jet of $\Theta$ along the polynomial $z\mapsto b'_{i'}(zy)$ for this $y$.

\emph{Step 2: the null family.} Let $M_{\tau,h}$ be the matrices of Lemma~\ref{l:frames} for $Y=y$ and $\pi$. Take
new variables $X_{(h,k),(\tau,b)}$, for $k<h<N$, $\tau\in\Phi$, $b\in n'$, and the linear forms
\[Z_{(\tau,b,a)}=\sum_{k<h<N}p^{\,g(\tau,h,b)}\Bigl(\sum_{k'<h}(M_{\tau,h})_{k,k'}\,y_{\tau,a,k'}\Bigr)X_{(h,k),(\tau,b)} .\]
The coupling of exponent $e\in S$ of this family is, modulo $p$, a combination of the pairings of
Lemma~\ref{l:frames}(ii) with coefficients $p^{\,g(\tau,h,b)p^e+g(\tau',h',b)}$, $\tau'=\mathrm{pr}_e(\tau)$. The pairings that are not exempted
vanish modulo $\pi$; for the exempted ones the hypothesis on the charges makes the coefficient divisible by $p$.
So the family is null, and the hypothesis of the theorem says that the coefficients of degree at most $B$ of
$\Theta(\varpi Z)$ are integral.

\emph{Step 3: the jet.} Multiplying the columns of a minor by $M_{\tau,h}$ and its row at $b$ by
$\varpi p^{g(\tau,h,b)}$ multiplies it by $\det M_{\tau,h}\prod_b\varpi p^{g(\tau,h,b)}$. So
\[\prod_{\tau,h}(\det M_{\tau,h})^{\varkappa(\tau,h)}\cdot\varpi^{B}p^{m}\cdot b'_{i'}(zy)=\mathsf{D}\bigl((\varpi Z)\bigr),\qquad
\mathsf{D}(X)=\prod_{\tau,h,l}\det\bigl(X_{(h,k),(\tau,i'_\tau(h,l,s))}\bigr)_{s,k<h},\]
where the right side is read through the transpose of the linear substitution; by the chain rule for jets
(quoted), the jet of $\Theta$ along the left side is the jet of $\Theta(\varpi Z)$ along $\mathsf{D}$. Now
$\mathsf{D}$ is a sum of monomials $\pm X^\beta$ of degree $B$, and the jet along $X^\beta$ is $\beta!$ times a
coefficient of degree $B$, which is integral. The exponent $\beta$ has at most $|\Phi|N^3$ entries, each at most
$B$, and $v_p(l!)\ge l/(p-1)-(\lfloor\log_pB\rfloor+1)$ for $l\le B$ by Legendre's formula; so
$\np{\beta!}\le p^{-B/(p-1)+|\Phi|N^3(\lfloor\log_pB\rfloor+1)}$. By Lemma~\ref{l:frames}(i), $\np\varpi=p^{-\rho}$ and
$\np\pi=p^\omega$, the factor on the left has absolute value at least
$p^{\,\omega\sum\mathrm{cc}-\rho B-m}$. Dividing gives the theorem.
\end{proof}
\leanline{\lean{normschurJetle}\sep quoted: \lnm{sum\_schur\_smul\_nablaBasis}, \lnm{exists\_val\_coeff\_le\_val\_eval},
\lnm{jet\_subst\_linearSeries}}

\sect{Sharpness}{s:sharp}
\begin{proposition}[a limit to the divisibility of Taylor coefficients]\label{p:katz}
Let $d$ be a prime number $\equiv7\bmod8$, $C$ a CM datum for $d$, $(p,j)$ a $p$-adic place and $n$ a natural
number. Then there is $F\in\phiring{C}$ with
\[\bigl(p^{\,\nsw{d}{p}-1/(p-1)}\bigr)^{n}\ \le\ \np{n!\cdot F_n} .\]
\end{proposition}
\begin{proof}
Let $c_0=p^{\nsw{d}{p}-1/(p-1)}$ and, for $k\in\NN$, let $c(k)$ be the largest value of $\np{k!F_k}/c_0^k$ for
$F\in\phiring{C}$; by (Q5) it is attained at a power of $\hat\varphi=\tayj{\phC{C}}$. \emph{$c$ is submultiplicative.}
By (Q5), $\hat\varphi(X+Y)^l$ is an integral series in $\hat\varphi(X)$ and $\hat\varphi(Y)$; its coefficient at
$X^nY^k$ is $\binom{n+k}{n}(\hat\varphi^l)_{n+k}$, and it is bounded by the largest product of an $n$-th coefficient
of a power of $\hat\varphi$ with a $k$-th one. Multiplying by $n!\,k!$ gives $c(n+k)\le c(n)c(k)$. Suppose the
proposition fails for $n$: $c(n)<1$. Then $c(k)\to0$, so there is $l_0$ with $c(k)<\np\varpi$ for $k\ge l_0$,
where $\varpi$, $\pi$, $S$, $T$, $\nu$ are those of Proposition~\ref{p:formal}(I). From $T(\varpi X)=\varpi S$,
$\np{T_k}\le c(k)c_0^k/\np{k!}$, $\np{k!}\ge p^{-k/(p-1)}$ and $\np\varpi=p^{-\nsw{d}{p}}$ we get
$\np{S_k}=\np\varpi^{k-1}\np{T_k}\le c(k)/\np\varpi<1$ for $k\ge l_0$. So the reduction $\cj S$ of $S$ modulo $\mxp$ is
a polynomial. But $\cj S-\cj S^{\,p^\nu}=X$, because $\np\pi<1$, and no polynomial satisfies this: the degree of the
left side is $p^\nu\deg\cj S\ne1$.
\end{proof}
\leanline{\lean{katzdivisibilitylimit}}

\begin{lemma}[the live sum]\label{l:live}
\begin{enumerate}[label=(\alph*),itemsep=2pt]
\item Let $d$ and $p$ be prime numbers with $\chid{d}(p)\ne1$ and $c$ a real number. Then
\[\sum_{\substack{e\ge1\\ e\text{ odd or }p=d}}\frac{c}{p^e}=\frac{2}{1-\chid{d}(p)}\cdot c\cdot\Bigl(\frac1{p-1}-\nsw{d}{p}\Bigr) ,\]
the series being convergent.
\item Let $D$ be admissible data for parameters $P=(N_0,d)$, $p$ a prime number with $\chiK{P}(p)\ne1$, $y$ a real
number and $\tau\in\PlD{D}$. Then the family of the terms
$\bigl(y/p^e-\nomW{D}{\pinD{D}{-p^e}\cdot\tau}\bigr)_+$, over the $e\ge1$ that are odd or for which $p=d$ (and $0$
for the other $e$), is summable; so the sum in $\ask{D}{p}{y}{\tau}$ is an ordinary sum.
\item Let $D$ be admissible data for parameters $P=(N_0,d)$, $p$ and $e$ natural numbers and $\tau\in\PlD{D}$. If
$\pinD{D}{-p^e}\cdot\tau$ and $\tau$ have the same class modulo $N_0$, then $p\ne d$.
\item Let $D$ be admissible data, $\tau\in\PlD{D}$, and $h$, $a$ natural numbers. Then, with $k=\kD{D}$,
\[\begin{cases}a&h=1\\ \kap{D}{\tau}{h}&h\ne1\end{cases}\ \ \le\ \ \begin{cases}a&h=1\\0&h\ne1\end{cases}
+\begin{cases}\wt{D}{1}{\tau}&h=k\\0&h\ne k\end{cases}+\gap{D} .\]
\end{enumerate}
\end{lemma}
\begin{proof}
(a) In both cases, $\chid{d}(p)=-1$ and $\chid{d}(p)=0$ (then $p=d$), the term at $e\ge0$ is
$\frac{c}{1-\chid{d}(p)}\bigl(p^{-e}-(\chid{d}(p)/p)^e\bigr)$; sum the two geometric series. (b) The terms are
bounded by $|y|/p^e$, since the nominal weight is non-negative; use (a). (c) If $p=d$, then $-d^e$ is prime to
$\NoneD{D}$, the class of $\pinD{D}{-d^e}\cdot\tau$ is $-d^e$ times that of $\tau$, and equal classes would give
$d^e\equiv-1\bmod N_0$, against (I4). (d) Unfold Definition~\ref{d:kappa}.
\end{proof}
\leanline{\lean{hasSumliveinvpow}\sep\lean{summablelocalAsk}\sep\lean{nedofplaceClasseq}\sep\lean{kaple}}

\statement{inksoft}{\textbf{Inputs from unit 152 on the sizes of the data.} For admissible data $D$ with components
$x$, $x_1$, $b$ and $c=\lceil60b/\kD{D}\rceil$:
\textbf{(S1)} $\aD{D}/x_1\le\nomW{D}{\tau}\le\aD{D}/x$ and $\nomW{D}{\tau}+x^{42}\le2b$ for all $\tau$;
$x^{42}+1\le\aD{D}/x_1$; $x^{42}\le c$; $d\le\kD{D}$ and $0<\kD{D}$.
\textbf{(S2)} For a prime $p\le\kD{D}$ and $e\le\muD{D}$: $2(\btwo{D}+\aD{D}/x_1)\le c/(p\cdot2\cdot2^e)$; and for
$e\ge\muD{D}$: $2b/p^e\le c/(p\cdot2\cdot2^e)$. Here $\btwo{D}$ is the least multiple of $\NoneD{D}$ that is at least
$\aD{D}/x_1$.
\textbf{(S3)} The thick places have class $-1$ modulo $N_0$. \emph{Three-term condition}: if $0<i<l\le324\,\muD{D}$
and $\tau$, $\tau_1$, $\tau_2$ are thick places with $\phi(\tau_1/\tau)^l=\phi(\tau_2/\tau)^i$, then
$\phi(\tau_1/\tau)=1$, where $\phi$ is the class modulo $2^{\muD{D}+1}$.
\textbf{(S4)} (background library) If $|u|>1$, $o>0$ and $u^o\equiv1\bmod2^{s+1}$ for integers $u$, $o$, $s$, then
$2^s\le(|u|+1)\,o$.}
\leanline{unit 152: \lnm{Data.nominalWeight\_mem\_Icc}, \lnm{Data.nominalWeight\_add\_le}, \lnm{Data.le\_of\_hasLogBound},
\lnm{Data.nat\_le\_of\_hasLogBound}, \lnm{Data.k\_pos}, \lnm{Data.placeClass\_of\_thick}, the field \lnm{three\_term} of
\lnm{Data}\sep library: \lnm{two\_pow\_le\_of\_pow\_eq\_one}}

\begin{proposition}[a Frobenius packing for the nominal weights]\label{p:packing}
Let $D$ be admissible data for parameters $P=(N_0,d)$, with component $x$, and let $n\colon\PlD{D}\to\NN$ satisfy
$\nomW{D}{\tau}\le n_\tau\le\nomW{D}{\tau}+x^{42}$ for every $\tau$. Let $p$ be a prime number with $\chiK{P}(p)\ne1$.
Then there are functions $i,y\colon\PlD{D}\to\NN$ such that $\packing{p}{n}{i}{y}$ and, with
\[L=\forced{D}{p}+\sum_{\tau\in\PlD{D}}v_p\bigl(\dl{D}{n_\tau}\bigr) ,\]
either $\chiK{P}(p)=-1$, $i_\tau=1$ for all $\tau$ and $\sum_\tau\bigl(v_p(n_\tau!)-y_\tau\bigr)\le L$, or
\[\sum_{\tau\in\PlD{D}}\Bigl(\bigl(\tfrac1{p-1}-\nsw{d}{p}\bigr)n_\tau+4\,\btwo{D}-y_\tau\Bigr)\ \le\ L .\]
\end{proposition}
\begin{proof}
Write $k=\kD{D}$, $c=\lceil60b/k\rceil$, $w=\nomW{D}{\,\cdot\,}$, $[u]=\pinD{D}{u}$, and $A(l)=(l-c+1)\cdots(l-c+c)$
with truncated $l-c$, so that $\dl{D}{l}=A(l)^2$. Two facts: $l!$ divides $(l-c)!\,A(l)$; and $c!$ divides $A(l)$,
a product of $c$ consecutive integers.

\emph{Case $k<p$.} Then $p\ne d$ by (S1), so $\chiK{P}(p)=-1$. Take $i_\tau=1$ and
$y_\tau=\min\bigl(n_{[-p]\tau},\lfloor(n_\tau-c)/p\rfloor\bigr)$. By (S1), $\lfloor(n_\tau-c)/p\rfloor\le2b/k\le c$.
Packing: (1) holds, (2) is empty, and for (3) the only $\tau$ with $[-p]\tau=\tau'$ has
$y_\tau\le\min(n_{\tau'},c)$, while $p\,y_{\tau'}\le n_{\tau'}-c$ (truncated); the sum is at most $n_{\tau'}$. For the bound,
Legendre's formula and the two facts give
\begin{align*}
v_p(n_\tau!)&\le v_p\bigl((n_\tau-c)!\bigr)+v_p(A(n_\tau))\\
&=\lfloor(n_\tau-c)/p\rfloor+v_p\bigl(\lfloor(n_\tau-c)/p\rfloor!\bigr)+v_p(A(n_\tau))
\ \le\ \lfloor(n_\tau-c)/p\rfloor+v_p(\dl{D}{n_\tau}) ,
\end{align*}
and $\lfloor(n_\tau-c)/p\rfloor-y_\tau\le\bigl((n_\tau-c)/p-n_{[-p]\tau}\bigr)_+\le\bigl(w(\tau)/p-w([-p]\tau)\bigr)_+$,
because $n_\tau-c\le w(\tau)$ by (S1). The last number is the term $e=1$ of $\ask{D}{p}{w(\tau)}{\tau}$ (all terms are
non-negative, Lemma~\ref{l:live}(b)). Summing over $\tau$ gives the first alternative.

\emph{Case $p\le k$.} For odd $e$, $[-p^e]=[-p]^e$. \emph{Claim: for $0<i<l\le324\muD{D}$ the places $\tau$,
$[-p]^i\tau$, $[-p]^l\tau$ are not all thick.} Otherwise (S3) gives $(-p)^i\equiv1\bmod2^{\muD{D}+1}$, and (S4) gives
$2^{\muD{D}}\le(p+1)i$, which the sizes of the data exclude for $p\le k$. Say that $e$ \emph{pairs} $\tau$ if
$e<\muD{D}$, $e$ is odd, $p\ne d$, and $\tau$ and $[-p]^e\tau$ are thick. By the claim at most one $e$ pairs $\tau$.
If $e$ pairs $\tau$ put $i_\tau=e$, $y_\tau=\lfloor n_\tau/p^e\rfloor$; otherwise $y_\tau=0$. Packing: (1) holds since
$p\ne d$ and $\chiK{P}(p)\ne1$. (2): if $[-p]^e\tau'=[-p]^{i_\tau}\tau$ with $e<i_\tau$ odd, then
$\tau'=[-p]^{i_\tau-e}\tau$, and $\tau$, $\tau'$, $[-p]^{i_\tau}\tau$ would be thick, against the claim. (3): if no
$\tau$ with $y_\tau\ne0$ has $[-p]^{i_\tau}\tau=\tau'$, the sum is $p^{i_{\tau'}}y_{\tau'}\le n_{\tau'}$. Otherwise there
is exactly one such $\tau$, and $y_{\tau'}=0$, both by the claim; and $y_\tau\le n_\tau/p\le n_{\tau'}$ since both
places are thick and $n_\tau\le\aD{D}/x+x^{42}\le2\,n_{\tau'}$ by (S1).

\emph{The bound.} Fix $\tau$. We compare, term by term in $e\ge1$ ($e$ odd or $p=d$),
\begin{equation}\label{e:term}
\frac{n_\tau}{p^e}\ \le\ \Bigl(\frac{w(\tau)}{p^e}-w([-p^e]\tau)\Bigr)_++[e=i_\tau]\,y_\tau+\frac{c}{p\cdot2\cdot2^e} .
\end{equation}
If $e\ge\muD{D}$, the left side is at most $2b/p^e$ and (S2) applies. Let $e<\muD{D}$. If $e$ pairs $\tau$, then
$n_\tau/p^e\le y_\tau+1$ and the last term is at least $1$ by (S2). If not, then $\tau$ and $[-p^e]\tau$ are not both
thick: for $p=d$ this is Lemma~\ref{l:live}(c) with (S3), and for $p\ne d$ it is the definition. If $\tau$ is not
thick, $n_\tau\le2\aD{D}/x_1$ by (S1), and (S2) applies. If $\tau$ is thick and $[-p^e]\tau$ is not, then
$w([-p^e]\tau)=\aD{D}/x_1$ and $n_\tau/p^e\le(w(\tau)/p^e-\aD{D}/x_1)_++\aD{D}/x_1+x^{42}$, and (S1), (S2) apply.
Summing \eqref{e:term} over $e$ with Lemma~\ref{l:live}(a), and multiplying by $(1-\chiK{P}(p))/2\le1$:
\[\bigl(\tfrac1{p-1}-\nsw{d}{p}\bigr)n_\tau-y_\tau\ \le\ \ask{D}{p}{w(\tau)}{\tau}+\frac cp .\]
Finally $v_p(\dl{D}{n_\tau})\ge2v_p(c!)\ge2(c/p-1)\ge c/p+4\btwo{D}$, the last step by (S2) at $e=0$. Summing over
$\tau$ gives the second alternative.
\end{proof}
\leanline{\lean{existsisFrobPacking}}

\begin{proof}[Proof of Theorem F]
We write $z_\tau$ for the variable $z_{(\tau,\ast)}$. By (I4), $p\ge13$. Let $\varpi$ be as in Theorem~D and
$F=\Theta(\varpi z)$. If $y=0$ the assertion holds with $\beta=\gamma=0$, since
$\np{(\Theta^{-1})_0}\np{\Theta_0}=1$; so let $y\ne0$. For $\tau$ with $y_\tau\ne0$ let $\tau^\ast=\pin{P}{\mu}{-p^{i_\tau}}\cdot\tau$; by the packing
$i_\tau$ is odd and $\chiK{P}(p)=-1$, so Theorem~D(b) gives $F=\exp(\kappa\Psi)V$ with $\np\kappa=p^{1-1/(p-1)}>1$.
As $\Theta(0)=1$ and $\Psi(0)=0$, $V$ is a unit of $\op[[z]]$.

\emph{Step 1: one derivative.} Let $Q$ be an integral series and $\tau$ a place with $y_\tau\ne0$, and put
$R=\partial_{\tau^\ast}\Psi\cdot Q+\kappa^{-1}(\partial_{\tau^\ast}Q+Q\,V^{-1}\partial_{\tau^\ast}V)$, an integral series.
Then $\partial_{\tau^\ast}(QF)=\kappa\,RF$, hence for every polynomial $G$:
$\np{\jetalong{G\cdot z_{\tau^\ast}}{QF}}=\np\kappa\cdot\np{\jetalong{G}{RF}}$. Modulo $\mxp$,
$\cj R=\partial_{\tau^\ast}\cj\Psi\cdot\cj Q$. Give the variable $z_\sigma$ the weight $1$ if $y_\sigma\ne0$ and
the weight $p^{\max i}$ otherwise. By Theorem~D(b) the monomials of $\partial_{\tau^\ast}\cj\Psi$ are the
$z_\sigma^{p^e}$ with $e$ odd and $\pin{P}{\mu}{-p^e}\cdot\sigma=\tau^\ast$; by condition (2) of the packing the one
of least weight is $z_\tau^{p^{i_\tau}}$, and it is the only one of that weight.

\emph{Step 2: all derivatives.} Starting from $Q=1$ and applying Step~1 $y_\tau$ times for every $\tau$, we get an
integral $Q$ whose reduction has the lowest term $u\,z^{\lambda}$, $u\ne0$,
$\lambda=\sum_\tau y_\tau p^{i_\tau}1_\tau$ (lowest terms multiply), and with
\[\np\kappa^{|\gamma|}\,\np{\jetalong{G}{QF}}=\np{\jetalong{G\cdot z^{\gamma}}{F}}\qquad\text{for all }G,\qquad
\gamma=\sum_\tau y_\tau1_{\tau^\ast} .\]
\emph{Step 3: a large coefficient.} The coefficient $Q_\lambda$ is a unit. Since $Q=F^{-1}\cdot(QF)$, there are
exponents $\gamma_1+\beta_1=\lambda$ with $1\le\np{(F^{-1})_{\gamma_1}}\np{(QF)_{\beta_1}}$. Step~2 with $G=z^{\beta_1}$
gives $\np\kappa^{|\gamma|}\np{\beta_1!}\np{(QF)_{\beta_1}}=\np{(\beta_1+\gamma)!}\np{F_{\beta_1+\gamma}}$. With
$F_\alpha=\varpi^{|\alpha|}\Theta_\alpha$, $(F^{-1})_\alpha=\varpi^{|\alpha|}(\Theta^{-1})_\alpha$,
$\np\varpi=p^{-\nsw{d}{p}}$ and $\np{\beta_1!}/\np{(\beta_1+\gamma)!}\ge\np{\gamma!}^{-1}$, this yields, for
$\beta=\gamma_1+\beta_1+\gamma$,
\[\np{(\Theta^{-1})_{\gamma_1}}\cdot\np{\Theta_{\beta_1+\gamma}}\ \ge\ p^{\,(1-\frac1{p-1})|\gamma|+\nsw{d}{p}|\beta|+\sum_\tau v_p(\gamma_\tau!)} .\]
Split $\gamma_1$ into five exponents by the hypothesis on $\Theta^{-1}$ and keep $\beta_1+\gamma$ as the sixth, by the
hypothesis on $\Theta$. Finally $|\gamma|=\sum y_\tau$, $|\beta|=|\lambda|+|\gamma|=\sum y_\tau(p^{i_\tau}+1)$, and
$\gamma\le\beta\le n$ by condition (3) of the packing: the entry of $\beta$ at $\tau'$ is
$p^{i_{\tau'}}y_{\tau'}+\sum_{\tau^\ast=\tau'}y_\tau$.
\end{proof}
\leanline{\lean{existshasLargeSplitofpairs}}

\sect{Where the results are used}{s:use}
By name, as the dossier of the unit lists the uses (a search over the sources, not a kernel export).
The unit on the columns (unit~145) uses all the definitions of §\ref{s:setting} on asks, gains, losses, served
primes, collisions, coupling counts and column counts, Theorems~C, D and E, Proposition~\ref{p:raises} and
Theorem~A, to prove that the chain of lifts is divisible by $p^{\gain{D}{p}}$ up to the loss. The unit on the rows
uses Proposition~\ref{p:katz}, Proposition~\ref{p:packing} and Theorem~F for the lower bound of a jet (the glued
form), and Theorem~B, with Lemma~\ref{l:vac} and Lemma~\ref{l:adm}, where conjugates of lifts are bounded. The unit
on the data uses Proposition~\ref{p:gen}: its constant $\cP{P}$ contains the chosen number $\Cgen{P}$.

\sect{Notes on the formal text}{s:notes}
\begin{enumerate}[label=\textup{\arabic*.},leftmargin=1.6em,itemsep=5pt]
\item \emph{Arguments that the text of the census lists and Lean does not take.} The dossier of the unit was made
without running Lean, and its lists of arguments are those of a reading of the text. For three definitions the
reading lists a section variable that only occurs as a bound index: $\tau$ in the gain and in the collision sum,
and the height $h$ in the coupling count. The source applies them to two, three and seven arguments
(\lnm{D.gain p}, \lnm{D.collisionSum p ev}, \lnm{couplingCount p κ N ev σ (pr e σ) e}), and this note writes them
so, $\gain{D}{p}$, $\coll{D}{p}{\Xi}$, $\cc{p}{\varkappa}{n'}{\Xi}{\tau}{\tau'}{e}$, through macros of its own; the
build accounts for the extra entry by the bound letter in Definition~\ref{d:ask}, \ref{d:served}, \ref{d:cc}. The
same happens for arguments of theorems and of one more definition: the index set $J$ of a theta product, bound by
``there are''; a ring $Q$ bound inside Lemma~\ref{l:frob}; the letters $c$ and $K$ in Theorem~B; a field $E$ in
Lemma~\ref{l:std}(d) and Proposition~\ref{p:F}; the prime and its place in Lemma~\ref{l:cyclo}; a second
finiteness of $n$ in the parts of Lemma~\ref{l:vac}; and the instance ``$V$ is standard'' in
Lemma~\ref{l:line}(c), whose space is only assumed to have signs $1$. All of this should be confirmed with
\lnm{\#check} once a build is at hand.
\item \emph{Arguments that Lean takes and the reading cannot see.} Four places use the notation of the formal
text for the $p$-adic place: Lemma~\ref{l:padic}(b)--(d), Theorem~A, Definition~\ref{d:raises} and
Proposition~\ref{p:raises}. This notation makes the prime $p$ and the homomorphism $j$ arguments (an implicit one
and two instance arguments). The statements above name them. Likewise the module $\Std{V}{\delta}$ depends on $j$, and so does
the set $\ThProd{P}{\mu}{\delta}{r}{c}$; the notation does not show it, because $j$ is an instance argument.
\item \emph{Three statements that the census cuts short.} The statements of Proposition~\ref{p:raises}(b), of
Proposition~\ref{p:packing} and of Theorem~E contain a \lnm{let}, which the census takes for the end of the
statement. I read these statements in the source in full; the quotations of Appendix~A avoid them. For the same
reason three intermediate facts of the dossier are cut short (in the proofs of Lemma~\ref{l:frames}, of
Lemma~\ref{l:deriv} and of Theorem~A); its claims state them in full, and I read the lines.
\item \emph{An error in the words of the dossier.} In the structured proof of Proposition~\ref{p:formal} the
dossier says that the pairs $(\zeta^i,1)$ lie in the closure $\mathcal{N}_d$. The lemma of the library that the
formal proof applies gives $(\zeta^i,\zeta^{i(p+1)})$, as in (Q6) above, and this is what makes the proof work: the
exponential factor of the normal form is the part of $\log\thC{C}$ in the degrees $\equiv p+1$ modulo $q-1$
precisely because of the exponent $p+1$. With the pairs $(\zeta^i,1)$ one would obtain the part in the degrees
$\equiv0$. The quotations of the dossier are exact; only its words are affected (steps 7 and 35.6).
\item \emph{Proofs that are shorter here, or differently arranged.} The frame identity of Step~2 of Theorem~A is
established in the formal text by a Gröbner basis computation; the computation above replaces it. The estimate at
the end of Lemma~\ref{l:deriv} is closed in the formal text by automation from the bounds of Lemma~\ref{l:val}; the
division into the cases $I\le4$ and $I\ge5$ is mine. In Proposition~\ref{p:katz} the formal text uses
$\cj S\circ(X-X^{p^\nu})=X$ where I use $\cj S-\cj S^{\,p^\nu}=X$. In Step~4 of Proposition~\ref{p:formal} the
arithmetic in the three cases is done by automation in the formal text.
\item \emph{What I could not settle without Lean.} (a) In Lemma~\ref{l:frames} the formal proof encodes the choice
of the device for each coupling in one predicate and verifies the equality of exponents by a case analysis with
automation; I read the construction and checked the three costs against the coupling count, but did not replay the
case analysis. (b) In Theorem~B(iv) the packing estimate and the final inequality between exponentials are a chain
of automated steps; the proof above gives their structure and not their constants. (c) In
Proposition~\ref{p:formal} the final identity of Step~6 is a single simplification in the formal text; the
computation above is the mathematics that it must carry out, with the pairs of note~4. (d) The inputs (S1)--(S2)
are not separate declarations: the formal proof of Proposition~\ref{p:packing} obtains each of them in place from
the lemma of unit~152 that compares logarithmic sizes. I state them as I read them at the places of use; the
second half of (S2) is my summary of a short chain there.
\item \emph{Choices.} $\cmE{P}$, $\Omz{P}$, $u_P$, $\vsc{P}{\mu}$ and $\Cgen{P}$ are obtained by choice from existence
statements, the last two from Proposition~\ref{p:gen} itself; so the definition of the scaled vacuum depends on the
proof of that proposition having been given, not on its witnesses. The file of unit~152 that defines $\cP{P}$
imports this unit for the same reason.
\item \emph{Conventions of the formal text that the statements inherit.} Sums over $e\in\NN$ in
Definitions~\ref{d:ask} and~\ref{d:served} are unconditional sums with value $0$ when the family is not summable;
Lemma~\ref{l:live}(b) shows that this does not occur in the ask for primes $p$ with $\chiK{P}(p)\ne1$, and for
$\chiK{P}(p)=1$ the factor in front is $0$. The definitions take a natural number $p$, not a prime. In
Lemma~\ref{l:val} the last hypothesis has a default proof (the formal text supplies it by arithmetic automation
when the numbers are explicit), and in Lemma~\ref{l:adm} the adele $\eta$ has a default value. In
Lemma~\ref{l:live}(d) the number $a$ is arbitrary. The hypothesis $\nu\ne0$ of Lemma~\ref{l:val} follows from
$p^\nu=Q+1$ only if $Q\ne0$.
\item \emph{The catalogue.} The relations stated for (Q2), (Q4) and (Q8) are those of the catalogue of the handout:
the main theorem of complex multiplication in a soft form, the integrality of $a_1$ as a weaker form of good
reduction everywhere, and the height of the formal group as one direction and an upper bound. For (Q3), (Q5),
(Q6), (Q7), (Q9) and (Q10) the catalogue has no confirmed place; the files of the library name Bannai--Kobayashi
(2010), Bannai--Kobayashi--Yasuda (2017), Katz (1981), Finis (2006), Mumford (1966), Conway--Sloane (1989) and
Shimura (1976) as related work, and I did not compare the statements with these sources.
\end{enumerate}

\subhead{How this note was written.} From the dossier of the unit: the texts of its \UnitMembers{} members and its
\UnitProofs{} structured proofs, which I read in full in a condensed form (every claim, without the quoted lines).
The program check of the dossier covered \UnitSteps{} leaf steps of \UnitProofs{} proofs and found every statement
at its lines. I read all twelve Lean files of the unit, statements and proofs, and compared every statement of
this note with its Lean text, argument by argument, with the list of arguments of the census beside it; for the
statements named in notes 1 to 3 the source decided. In other units and in the library I read the definitions
that the statements use (the structure of the data and its functions, tight frames, vacuum expansions, the jet
series, Schur jets, gauges, kernels and the theta kernel, the CM datum with its functions, the ring
$\phiring{C}$, carried series, admissible combinations) and the statements of the quoted results (Q2), (Q3),
(Q6), (Q7), (Q8), of the input of unit~149 and of (I6); the other inputs are stated as the dossier reads them. I
did not read proofs of the library. No Lean was run.

\head{Appendix A. Concordance with the Lean artifact}
{\fontsize{8.6}{11.6}\selectfont\renewcommand{\leanfont}{\ttfamily\fontsize{7.7}{10}\selectfont\color{leancol}}
\setlength{\LTleft}{0pt}\setlength{\LTright}{\fill}\setlength{\tabcolsep}{0pt}
\begin{longtable}{@{}p{30mm}p{59mm}p{40mm}r@{}}
\textbf{In this note} & \textbf{Lean declaration} & \textbf{File} & \textbf{Lines}\\\midrule\endhead
\ConcordanceRows
\end{longtable}}

\subhead{Four declarations as they are written in the artifact.}
\quoted{gain}
\quoted{couplingCount}
\quoted{thetaProducts}
\quoted{thetaProductsnormalform}
In these texts {\leanfont s} is the exponent $\mu$, {\leanfont Φ τ} the embedding $\emb{\tau}$, {\leanfont ιₚ} the
homomorphism $j$, {\leanfont valRingCp p} the ring $\op$, {\leanfont chiK P.d p} the value $\chiK{P}(p)$,
{\leanfont nonsplitWeight P.d p} the number $\nsw{d}{p}$, {\leanfont P.placeOfInt s (-p \textasciicircum{} j) * v} the
place $\pin{P}{\mu}{-p^e}\cdot v$, {\leanfont κ} a family of column counts $\varkappa$, {\leanfont ev} the set
$\Xi$, {\leanfont taylorVia ιₚ} the series $\tayj{\,\cdot\,}$, {\leanfont P.cmCurve.twistedTranslate 1 y} the
translation $\UC{\cmE{P}}{1}{y}$, {\leanfont expSubst} the exponential of a series without constant term and
{\leanfont residue} the reduction modulo $\mxp$. Letters of the formal text were changed where the write-up
reserves them: the exponent {\leanfont s}, the Frobenius exponents {\leanfont j} (here $e$), the degree
{\leanfont e} of the residue field (here $\nu$), the element {\leanfont θ} of the null cone (here $\pi$) and the
column counts {\leanfont κ}.

\head{Notation}
\begingroup\fontsize{8.6}{10.9}\selectfont\raggedright
\setlength{\columnsep}{9mm}\setlength{\parskip}{0pt}
\begin{multicols}{2}
\nt{$P=(N_0,d)$, $\mu$}{parameters; the exponent}
\nt{$\KP{P}$, $\OO_{\KP{P}}$, $\chiK{P}$, $\chid{d}$}{$\Q(\sqrt{-d})$, its integers, its quadratic character; the Legendre symbol modulo $d$}
\nt{$\EP{P}{\mu}$, $\Nn{P}{\mu}$, $\rtu{P}{\mu}$}{the CM field, the order of its root of unity, that root}
\nt{$\PlP{P}{\mu}$, $\emb{\tau}$, $\pin{P}{\mu}{\nu}$}{the group of infinite places; the embedding at $\tau$; the place of an integer $\nu$}
\nt{$\tau,\tau',v,v',u,\sigma$}{infinite places; $\sigma$ also an automorphism of $\CC$, $u$ also a unit or a variable}
\nt{$\czE{P}{\mu}$, $\dgen{P}{\mu}$, $\alP{P}{\mu}$, $\rdE{P}{\mu}$}{$(1+\sqrt{-d})/2$; the generator of the relative different; the idele; the root discriminant}
\nt{$\OmP{P}{\mu}$, $\mathrm{cs}(d)$}{the period $e^{\mathrm{cs}(d)}2\pi/\rdE{P}{\mu}$; the Chowla--Selberg constant}
\nt{$\frl{P}{\mu}$}{the lattice $E\cap\alpha\Ohat$}
\nt{$\Ad$, $\Af$, $\Ohat$}{adeles, finite adeles, integral finite adeles}
\nt{$\prm{\mu}{\sigma}$, $\pi$}{the permutation of places of $\sigma$; $\pi$ also an element of $\op$, and the circle constant}
\nt{$(p,j)$, $\Cp$, $\np{\,\cdot\,}$, $\op$, $\mxp$, $\Oj$}{a $p$-adic place and its objects}
\nt{$\psi$, $\mathfrak{p}_\psi$, $\hat\psi$, $j_\tau$, $\hat\jmath_\tau$}{an embedding into $\Cp$, its prime, its extension to $\Af$; $j\circ\emb{\tau}\circ$ conjugation}
\nt{$v_p$, $\lfloor\,\rfloor$, $\lceil\,\rceil$, $(\,)_+$}{$p$-adic valuation, floor, ceiling, positive part}
\nt{$C$, $\LamC{C}$, $\area{C}$, $t$, $a_1$}{a CM datum, its lattice, area, point of order two, invariant}
\nt{$\wp$, $\sigma_C$, $\thC{C}$, $\xiC{C}$, $\phC{C}$, $\mon{C}{x}$}{Weierstrass functions; theta; $\wp(z+t)-\wp(t)$; the parameter; monomial sections}
\nt{$\UC{C}{m}{y}$, $H_C$}{twisted translation; the hermitian form of $C$}
\nt{$\tay{F}$, $\tayj{F}$, $\hat\varphi$}{Taylor series; the same read through $j$; $\tayj{\phC{C}}$}
\nt{$\phiring{C}$, $\mathcal{N}_d$, $\cpart{M}{r}{F}$}{integral series in the parameter; the closure of norm pairs; part in a residue class of degrees}
\nt{$\szero{P}{\mu}$, $\sone{P}{\mu}$, $\per{P}{\mu}$, $\asc{P}$}{root of the vacuum; period factor; scale of jets; $\sqrt{2\pi/\sqrt d}$}
\nt{$\cmE{P}$, $\Omz{P}$, $u_P$}{the chosen CM datum, its period, its Chowla--Selberg number}
\nt{$\chi_{\mathrm{fin}}$, $\upsilon$, $\psi_0$}{characters of the Weil data; the additive character of $\Af$}
\nt{$V$, $n$, $\varepsilon$, $\stdsp{n}$, $\homsp{n}{n'}$}{a hermitian space, its index set, its signs; standard spaces}
\nt{$\Sw{V}$, $\Wf{V}$, $\vac{V}$, $\Th{V}$, $\weil{V}$, $\inn{V}{f}{f'}$}{the Weil space, its finite part, vacuum, theta functional, representation, pairing}
\nt{$\heis{V}{y}$, $\fheis{V}{x}$, $\fv{x}$}{Heisenberg operators; the adelic vector of a finite vector}
\nt{$\Uf{n}$, $\fmult{g}$, $1_n$}{the finite adelic unitary group, its multiplier, the unit matrix}
\nt{$\jet{V}{a}{f}$, $\Thf{f}$, $\jetalong{Q}{F}$}{jet series; theta function of $f$; jet along $Q$}
\nt{$\name{V}{R}{\Psi}$, $P^{\mathfrak m}$, $\mathfrak m$}{named operators; the mean at a modulus; a modulus}
\nt{$\vsc{P}{\mu}$, $\intvac{V}$, $\Lat{\delta}{S}$, $\Std{V}{\delta}$}{scale of the vacuum; scaled vacuum; a group of vectors; standard data}
\nt{$\Cgen{P}$, $\cP{P}$}{the chosen constant of Proposition~\ref{p:gen}; the constant of unit~152}
\nt{$\varkappa$, $\boxes{\varkappa}{N}$, $S_\varkappa$, $\schurjet{\varkappa}{F}$}{column counts, boxes, Schur matrix, Schur jets}
\nt{$i,i'$, $i_\tau(h,l,b)$}{Weyl indices and their entries; $i$ also an index or the exponents of a packing}
\nt{$\gauge{\varkappa}{e}{g}{\xi}$, $k$, $\omega_k$, $\mathrm{jm}_k$, $\mathcal D$}{gauge; a kernel, its representation, its matrices, its data}
\nt{$\thker{\varkappa}{c}{a}{b}{K}$, $\pairw{a}{b}$, $e_a$, $e_b$}{theta kernel; Weil representation of the pair; weights of the characters $a$, $b$}
\nt{$\Rais{k}{d_1}{d_2}{c}$}{data that raise gauges by at most $c$}
\nt{$\ThProd{P}{\mu}{\delta}{r}{c}$}{theta products}
\nt{$\mathrm{pr}_e$, $S$, $S_B$, $\Xi$}{the partner map; sets of Frobenius exponents; a set of exempted places}
\nt{$\varpi$, $\kappa$, $\kappa_0$, $\Psi$, $\psi$, $\nu$, $q$}{the scale at $p$; scalars of the exponential; its exponents; the degree and the norm of the prime above $p$}
\nt{$\nsw{d}{p}$}{$1/(p^2-1)$, $1/(2(p-1))$ or $1/(p-1)$}
\nt{$D$; $x,x_1,b,m,r$; $\muD{D}$}{admissible data; its components (in Definitions \ref{d:ask}--\ref{d:packing} and §\ref{s:sharp}); its exponent}
\nt{$\PlD{D}$, $\NoneD{D}$, $\fD{D}$, $\kD{D}$, $\tD{D}$, $\aD{D}$, $\pinD{D}{\nu}$}{objects of the data (unit~152)}
\nt{$\tops{D}$, $\nomW{D}{\tau}$, $\wt{D}{l}{\tau}$, $\nD{D}{l}$, $\gap{D}$, $\btwo{D}$}{tops; nominal weight; weights; ranks; gap; a rounded weight}
\nt{$\ask{D}{p}{y}{\tau}$, $\forced{D}{p}$, $\gain{D}{p}$, $\loss{D}{p}$}{ask, forced part, gain, loss}
\nt{$\clm{D}{p}$, $\deltaP{D}{p}$, $\jc{D}{p}$, $\ozero{D}{p}$, $\coll{D}{p}{\Xi}$}{the numbers of Definition~\ref{d:served}}
\nt{$\cc{p}{\varkappa}{n'}{\Xi}{\tau}{\tau'}{e}$, $t_\sigma$, $r_\sigma$, $B_\sigma$}{coupling count and its parts}
\nt{$\kap{D}{\tau}{h}$, $\last{D}{l}{\tau}$, $\kept{D}{h}$, $\kapY{D}{\tau}{h}$}{column counts of the data}
\nt{$\dl{D}{n}$, $A(l)$, $L$}{depth loss; an ascending factorial; a bound (also a linear form, a semilinear map, $\log$ of theta)}
\nt{$\mathsf{N}$, $\mathsf{b}$, $\mathsf{D}$}{a weight on exponents; a proposition; a product of determinants}
\nt{$\alpha,\beta,\gamma,\lambda$, $1_y$, $z^\beta$}{exponents (multi-indices), the exponent of one variable, monomials; $\alpha$ also the idele, $\beta,\gamma,\lambda$ also numbers where said}
\nt{$z$, $z_x$, $X$, $Y$, $T$, $Z$, $W$}{variables and families of series; $z$ also an adele or an integer, $T$, $Z$, $W$ series in Proposition~\ref{p:formal}}
\nt{$M$, $N$, $R$, $B$, $I$, $J$, $Q$, $A$, $H$}{lengths, levels, sets, rings, bounds: local to each statement}
\nt{$a,b,c,e,f,g,h,k,l,m,n,o,s,w,y$}{numbers, indices, elements and functions local to each statement}
\nt{$E$, $F$, $G$, $S$, $U$, $O$}{series, functions and sets local to each statement; $E$ also a field}
\nt{$\delta$, $\eta$, $\zeta$, $\theta$, $\rho$, $\omega$, $\phi$, $\xi$, $\ell$, $\iota$, $\varsigma$, $\epsilon$}{an integer of $\KP{P}$; an adele; a root of unity; in $\theta_C$; real exponents; a ring map; a vector; linear forms and primes; in $\iota_\tau$; the root of unity of $E$; coefficients}
\nt{$\Phi$, $\Theta$, $\Omega$, $\Lambda$, $\Delta$, $\Pi$, $\partial$}{a set of places or a datum; a series; periods; a lattice; a multiplier; a product; derivatives and the different}
\end{multicols}
\endgroup

\end{document}
